% Traduction et production de textes informatifs sur l'environnement, la santé et le bien-être — Allemand Tle A — Cours signé OnBuch+
% Programme officiel MINESEC. Compiler avec : tectonic AL15-traduction-et-production-de-textes-informatifs-sur-l-environ.tex
\newcommand{\DOCMATIERE}{Allemand}
\newcommand{\DOCNIVEAU}{Tle A}
\newcommand{\DOCMODULE}{Chapitre 3 : Environnement, santé et bien-être}
\newcommand{\DOCLECON}{AL15}
\newcommand{\DOCTITRE}{Traduction et production de textes informatifs sur l'environnement, la santé et le bien-être}
% OnBuch+ — préambule « cours signés » ENRICHI (compilé avec tectonic / XeLaTeX)
% Système visuel « Pop » d'OnBuch+ : fond crème (jamais blanc), bordures encre
% épaisses, ombres « dures » décalées, titres Archivo Black en capitales.
% Version MATHÉMATIQUES : graphes (pgfplots/tikz), boîtes pédagogiques
% (définition, propriété, méthode, attention, le savais-tu, exercice résolu,
% exercice, TP, bilan), page de garde et signature OnBuch+ ancrée en bas.
%
% Macros attendues dans main.tex AVANT \input{preamble} :
%   \DOCMATIERE  \DOCNIVEAU  \DOCMODULE  \DOCLECON (numéro)  \DOCTITRE
\documentclass[11pt]{article}

\usepackage{fontspec}
\usepackage[ngerman,french]{babel} % français = langue principale ; l'allemand via \all{..} / textebox
\usepackage[a4paper,margin=2cm,top=3.4cm,bottom=2.8cm,headheight=1.4cm,footskip=1.3cm]{geometry}
\usepackage{amsmath,amssymb,amsfonts}
\usepackage{mathtools} % \xmapsto, \coloneqq... souvent utilisés par les modèles
\usepackage{cancel} % simplifications barrées (bilans de forces)
% \abs{z} (module, valeur absolue) et \norm{u} (norme), utilisés par les modèles
\providecommand{\abs}{}\let\abs\relax\DeclarePairedDelimiter{\abs}{\lvert}{\rvert}
\providecommand{\norm}{}\let\norm\relax\DeclarePairedDelimiter{\norm}{\lVert}{\rVert}
\usepackage[table]{xcolor} % [table] : fournit aussi \rowcolors (alternance de lignes)
\usepackage{pagecolor}
\usepackage{tikz}
\usetikzlibrary{arrows.meta,positioning,calc,shapes.geometric,shapes.misc,shapes.arrows,decorations.pathmorphing,decorations.pathreplacing,decorations.markings,patterns,shadows,mindmap,trees,fit,backgrounds,matrix,intersections,angles,quotes}
\usepackage{pgfplots}
\usepackage[version=4]{mhchem} % équations nucléaires \ce{^{235}_{92}U}
\usepackage[european,siunitx]{circuitikz} % schémas électriques
\pgfplotsset{compat=1.18}
\usepgfplotslibrary{fillbetween,groupplots} % aires entre courbes (calcul intégral)
\usepackage{enumitem}
\usepackage{fancyhdr}
% [most] charge toutes les bibliothèques tcolorbox (listings, minted, raster,
% documentation...) et gonfle inutilement la mémoire TeX sur un document long
% (~300 boîtes) : on ne charge que ce qui est réellement utilisé.
\usepackage[skins,breakable]{tcolorbox}
\usepackage{graphicx}
\usepackage{colortbl}
\usepackage{booktabs}
\usepackage{tabularx}
\usepackage{diagbox} % case d'en-tête diagonale des tableaux à double entrée
% Colonnes extensibles centrée / gauche / droite (C, L, R), souvent utilisées par les modèles
\newcolumntype{C}{>{\centering\arraybackslash}X}
\newcolumntype{L}{>{\raggedright\arraybackslash}X}
\newcolumntype{R}{>{\raggedleft\arraybackslash}X}
\usepackage{array}
\usepackage{multicol}
\usepackage{multirow}
\usepackage{siunitx}
% parse-numbers=false : accepte aussi \SI{2,5 \times 10^{-3}}{...}
\sisetup{locale=FR,output-decimal-marker={,},group-minimum-digits=4,parse-numbers=false}
\DeclareSIUnit\ohm{\ensuremath{\symup{\Omega}}} % U+2126 absent de la police
\DeclareSIUnit\division{div} % réglages d'oscilloscope (ms/div, V/div)
\DeclareSIUnit\tr{tr} % tours (vitesse de rotation, ex. \SI{1500}{\tr\per\minute})
\DeclareSIUnit\molar{M} % concentration molaire (ex. \SI{3}{\molar}, \SI{1}{\micro\molar})
\DeclareSIUnit\an{an} % année (le modèle confond parfois avec \year, un compteur TeX) — ex. \SI{5}{\kilo\meter\per\an}
\DeclareSIUnit\FCFA{FCFA} % franc CFA (ex. \SI{500}{\FCFA\per\kilo\gram})
\DeclareSIUnit\degreeBrix{\textdegree Brix} % teneur en sucre d'un sirop/jus (réfractométrie)
\DeclareSIUnit\month{mois} % le modèle utilise parfois \month, un compteur TeX, comme unité de durée
\usepackage{qrcode}
% \degree : commande fréquemment utilisée par les modèles pour un angle ou une
% température en dehors de \SI{}{} (non fournie par les paquets chargés).
\providecommand{\degree}{^{\circ}}
% \qed (fin de démonstration), utilisé par les modèles sans amsthm
\providecommand{\qed}{\hfill$\square$}
% \barre : double barre des valeurs interdites dans un tableau de variations (tkz-tab)
\providecommand{\barre}{\ensuremath{\|}}
% environnement proof (amsthm non chargé) utilisé par les modèles
\ifdefined\proof\else\newenvironment{proof}[1][Démonstration]{\par\noindent\textit{#1.}\ }{\qed\par}\fi
\usepackage{lastpage}
\usepackage{hyperref}
\hypersetup{hidelinks}

% ============================================================
% Palette « Pop » OnBuch+ — mêmes hex que la classe P (plus_ui.dart)
% ============================================================
\definecolor{poporange}{HTML}{F59321}
\definecolor{popdark}{HTML}{E07A0C}
\definecolor{popcream}{HTML}{FFF6E8}
\definecolor{popink}{HTML}{1C1714}
\definecolor{poppurple}{HTML}{7A5AE0}
\definecolor{popgreen}{HTML}{2E9E6A}
\definecolor{poppink}{HTML}{E0457B}
\definecolor{popblue}{HTML}{2F7DE1}
\definecolor{popgold}{HTML}{C98A12}
\definecolor{popmuted}{HTML}{8A7F76}
\definecolor{popline}{HTML}{EADFD2}
\definecolor{popgray}{HTML}{8A7F76}
\colorlet{popgrey}{popgray}
% Alias tolérants (le modèle nomme parfois les couleurs différemment)
\colorlet{popwhite}{white}
\colorlet{popblack}{popink}
\colorlet{popred}{poppink}
\colorlet{popyellow}{popgold}
% Teintes claires (fonds de tableaux/encadrés) — le modèle nomme parfois
% les couleurs avec un suffixe L (ex. popblueL) sans les avoir définies.
\colorlet{poporangeL}{poporange!20}
\colorlet{popdarkL}{popdark!20}
\colorlet{poppurpleL}{poppurple!20}
\colorlet{popgreenL}{popgreen!20}
\colorlet{poppinkL}{poppink!20}
\colorlet{popblueL}{popblue!20}
\colorlet{popgoldL}{popgold!20}
\colorlet{popredL}{popred!20}
\colorlet{popgrayL}{popgray!20}
% Teintes claires pour fonds de tableaux / zones
\colorlet{popblueL}{popblue!10!white}
\colorlet{poporangeL}{poporange!14!white}
\colorlet{popgreenL}{popgreen!12!white}
\colorlet{poppurpleL}{poppurple!12!white}
\colorlet{poppinkL}{poppink!12!white}
\colorlet{popgoldL}{popgold!14!white}

\pagecolor{popcream}

% ============================================================
% Typographie
% ============================================================
\defaultfontfeatures{Ligatures=TeX}
\setmainfont{PlusJakartaSans-Medium.ttf}[Path=../../../fonts/,BoldFont=PlusJakartaSans-Bold.ttf]
% Maths en sans-serif (Fira Math) avec chiffres et lettres droites de Plus
% Jakarta, pour que les formules se fondent dans le texte.
\usepackage[math-style=ISO,bold-style=ISO]{unicode-math}
\setmathfont{FiraMath-Regular.otf}[Path=../../../fonts/]
\setmathfont{PlusJakartaSans-Medium.ttf}[Path=../../../fonts/,range={up/num,up/{Latin,latin},\mathup}]
\usepackage{icomma} % virgule décimale sans espace en mode math
% Caractères mathématiques Unicode absents des polices de texte (Plus Jakarta,
% Archivo Black) : rendus par leur équivalent mathématique, en texte comme en
% formule — sinon ils s'affichent en carré vide (ex. « Arithmétique dans ℤ »).
\usepackage{newunicodechar}
\newunicodechar{ℝ}{\ensuremath{\mathbb{R}}}
\newunicodechar{ℤ}{\ensuremath{\mathbb{Z}}}
\newunicodechar{ℂ}{\ensuremath{\mathbb{C}}}
\newunicodechar{ℕ}{\ensuremath{\mathbb{N}}}
\newunicodechar{ℚ}{\ensuremath{\mathbb{Q}}}
\newunicodechar{𝔻}{\ensuremath{\mathbb{D}}}
\newunicodechar{α}{\ensuremath{\alpha}}
\newunicodechar{β}{\ensuremath{\beta}}
\newunicodechar{γ}{\ensuremath{\gamma}}
\newunicodechar{δ}{\ensuremath{\delta}}
\newunicodechar{ε}{\ensuremath{\varepsilon}}
\newunicodechar{θ}{\ensuremath{\theta}}
\newunicodechar{λ}{\ensuremath{\lambda}}
\newunicodechar{μ}{\ensuremath{\mu}}
\newunicodechar{σ}{\ensuremath{\sigma}}
\newunicodechar{τ}{\ensuremath{\tau}}
\newunicodechar{φ}{\ensuremath{\varphi}}
\newunicodechar{ψ}{\ensuremath{\psi}}
\newunicodechar{ω}{\ensuremath{\omega}}
\newunicodechar{Σ}{\ensuremath{\Sigma}}
% majuscules produites par \MakeUppercase dans les titres (α-aminés -> Α)
\newunicodechar{Α}{\ensuremath{\alpha}}
\newunicodechar{Β}{\ensuremath{\beta}}
\newunicodechar{Γ}{\ensuremath{\Gamma}}
\newunicodechar{Δ}{\ensuremath{\Delta}}
\newunicodechar{Ε}{\ensuremath{\varepsilon}}
\newunicodechar{Θ}{\ensuremath{\Theta}}
\newunicodechar{Λ}{\ensuremath{\Lambda}}
\newunicodechar{Μ}{\ensuremath{\mu}}
\newunicodechar{Τ}{\ensuremath{\tau}}
\newunicodechar{Φ}{\ensuremath{\Phi}}
\newunicodechar{Ψ}{\ensuremath{\Psi}}
\newunicodechar{Ω}{\ensuremath{\Omega}}
\newunicodechar{↦}{\ensuremath{\mapsto}}
\newunicodechar{∈}{\ensuremath{\in}}
\newunicodechar{∉}{\ensuremath{\notin}}
\newunicodechar{∧}{\ensuremath{\wedge}}
\newunicodechar{⇔}{\ensuremath{\Leftrightarrow}}
\newunicodechar{⟺}{\ensuremath{\Longleftrightarrow}}
\newunicodechar{⇒}{\ensuremath{\Rightarrow}}
\newunicodechar{⟹}{\ensuremath{\Longrightarrow}}
\newunicodechar{∘}{\ensuremath{\circ}}
\newunicodechar{∪}{\ensuremath{\cup}}
\newunicodechar{∩}{\ensuremath{\cap}}
\newunicodechar{≡}{\ensuremath{\equiv}}
\newunicodechar{⊂}{\ensuremath{\subset}}
\newunicodechar{⊥}{\ensuremath{\perp}}
\newunicodechar{∃}{\ensuremath{\exists}}
\newunicodechar{∀}{\ensuremath{\forall}}
\newunicodechar{∛}{\ensuremath{\sqrt[3]{\,}}}
\newunicodechar{∜}{\ensuremath{\sqrt[4]{\,}}}
\newunicodechar{⃗}{\ensuremath{{}^{\to}}}
\newunicodechar{ⁿ}{\ifmmode^{n}\else\textsuperscript{\ensuremath{n}}\fi}
\newunicodechar{ˣ}{\ifmmode^{x}\else\textsuperscript{\ensuremath{x}}\fi}
\newunicodechar{ᵇ}{\ifmmode^{b}\else\textsuperscript{\ensuremath{b}}\fi}
\newunicodechar{ʳ}{\ifmmode^{r}\else\textsuperscript{\ensuremath{r}}\fi}
\newunicodechar{ᵘ}{\ifmmode^{u}\else\textsuperscript{\ensuremath{u}}\fi}
\newunicodechar{ʸ}{\ifmmode^{y}\else\textsuperscript{\ensuremath{y}}\fi}
\newunicodechar{ᵃ}{\ifmmode^{a}\else\textsuperscript{\ensuremath{a}}\fi}
\newunicodechar{ᵉ}{\ifmmode^{e}\else\textsuperscript{\ensuremath{e}}\fi}
\newunicodechar{ᵏ}{\ifmmode^{k}\else\textsuperscript{\ensuremath{k}}\fi}
\newunicodechar{ᵖ}{\ifmmode^{p}\else\textsuperscript{\ensuremath{p}}\fi}
\newunicodechar{ᴺ}{\ifmmode^{N}\else\textsuperscript{\ensuremath{N}}\fi}
\newunicodechar{ᶜ}{\ifmmode^{c}\else\textsuperscript{\ensuremath{c}}\fi}
\newunicodechar{ʰ}{\ifmmode^{h}\else\textsuperscript{\ensuremath{h}}\fi}
\newunicodechar{ᵠ}{\ifmmode^{q}\else\textsuperscript{\ensuremath{q}}\fi}
\newunicodechar{ₙ}{\ifmmode_{n}\else\textsubscript{\ensuremath{n}}\fi}
\newunicodechar{ₐ}{\ifmmode_{a}\else\textsubscript{\ensuremath{a}}\fi}
\newunicodechar{ᵢ}{\ifmmode_{i}\else\textsubscript{\ensuremath{i}}\fi}
\newunicodechar{ᵣ}{\ifmmode_{r}\else\textsubscript{\ensuremath{r}}\fi}
\newunicodechar{ₖ}{\ifmmode_{k}\else\textsubscript{\ensuremath{k}}\fi}
\newunicodechar{ᵦ}{\ifmmode_{\beta}\else\textsubscript{\ensuremath{\beta}}\fi}
\newunicodechar{ₓ}{\ifmmode_{x}\else\textsubscript{\ensuremath{x}}\fi}
\newfontfamily\popdisplay{ArchivoBlack-Regular.ttf}[Path=../../../fonts/]
\newfontfamily\poplogo{SpaceGrotesk-Bold.ttf}[Path=../../../fonts/]
\newfontfamily\popsemi{PlusJakartaSans-SemiBold.ttf}[Path=../../../fonts/]
\newfontfamily\popxbold{PlusJakartaSans-ExtraBold.ttf}[Path=../../../fonts/]
\newfontfamily\popxboldls{PlusJakartaSans-ExtraBold.ttf}[Path=../../../fonts/,LetterSpace=3.0]

\setlist[enumerate]{leftmargin=1.7em,itemsep=5pt,topsep=4pt}
\setlist[itemize]{leftmargin=1.7em,itemsep=4pt,topsep=4pt,label={\color{poporange}\textbullet}}
\setlength{\parindent}{0pt}
\setlength{\parskip}{5pt}
\renewcommand{\arraystretch}{1.25}

% pgfplots : style Pop par défaut
\pgfplotsset{
  every axis/.append style={axis line style={popink,line width=1pt},tick style={popink},
    grid style={popline,line width=0.5pt},label style={font=\small},tick label style={font=\footnotesize}},
  popcurve/.style={poporange,line width=1.8pt,no markers,smooth},
}
% Même style disponible dans un \draw TikZ simple (hors pgfplots)
\tikzset{popcurve/.style={poporange,line width=1.8pt}}
\tikzset{popgrid/.style={popline,very thin}}
\colorlet{popcurve}{poporange} % le nom du style est parfois utilisé comme couleur

% ============================================================
% Pill / squiggle
% ============================================================
\newcommand{\poppill}[3]{%
  \tikz[baseline=-0.55ex]\node[draw=popink,line width=1.2pt,rounded corners=2.6pt,%
    fill=#2,inner xsep=6.5pt,inner ysep=3pt]{%
    \popxboldls\fontsize{7}{7}\selectfont\color{#3}\MakeUppercase{#1}};%
}
\newcommand{\popsquiggle}[1]{%
  \tikz[baseline=-0.4ex]{
    \draw[#1,line width=2.6pt,line cap=round]
      plot [smooth] coordinates {(0,0.09) (0.34,0.01) (0.68,0.21) (1.02,0.01) (1.36,0.10)};
  }%
}
% Mot-clé surligné (marqueur orange sous le texte)
\newcommand{\cle}[1]{{\popxbold\color{popdark}#1}}

% ============================================================
% En-tête / pied
% ============================================================
\pagestyle{fancy}
\fancyhf{}
\renewcommand{\headrulewidth}{0pt}
\renewcommand{\footrulewidth}{0pt}
\lhead{\raisebox{-0.15ex}{\poplogo\fontsize{13}{13}\selectfont\color{popink}OnBuch\color{poporange}+}}
\rhead{\poppill{\DOCMATIERE\ \textbullet\ \DOCNIVEAU\ \textbullet\ Leçon \DOCLECON}{white}{popink}}
\lfoot{\popxbold\fontsize{7.6}{7.6}\selectfont\color{popmuted}\MakeUppercase{Cours signé OnBuch+}\ \textbullet\ {\popsemi\DOCTITRE}}
\rfoot{\tikz[baseline=-0.6ex]\node[rounded corners=8.5pt,fill=popink,minimum height=17pt,minimum width=17pt,inner xsep=5pt,inner sep=0]{\popxbold\fontsize{8}{8}\selectfont\color{popcream}\thepage\,/\,\pageref*{LastPage}};}

% ============================================================
% Titres
% ============================================================
\newcommand{\chaptitre}[2]{%
  \begin{tcolorbox}[enhanced,colback=poporange,colframe=popink,boxrule=2.6pt,arc=11pt,
      drop shadow=popink,left=15pt,right=15pt,top=12pt,bottom=11pt,before skip=3pt,after skip=18pt]
    \poppill{#2}{popink}{popcream}\par\vspace{9pt}
    {\raggedright\hyphenpenalty=10000\exhyphenpenalty=10000\popdisplay\fontsize{22}{24}\selectfont\color{popcream}\MakeUppercase{#1}\par}
  \end{tcolorbox}
}
% Section numérotée : pastille orange avec le numéro + titre Archivo
\newcounter{popsec}
\newcommand{\coursec}[1]{%
  \stepcounter{popsec}\setcounter{popsub}{0}%
  \par\vspace{16pt}\noindent
  \begin{minipage}{\linewidth}
  \tikz[baseline=-0.7ex]\node[draw=popink,line width=1.6pt,fill=poporange,rounded corners=4pt,minimum size=22pt,inner sep=0]{\popdisplay\fontsize{12}{12}\selectfont\color{popcream}\thepopsec};%
  \hspace{8pt}{\popdisplay\fontsize{15}{17}\selectfont\color{popink}\MakeUppercase{#1}}\par
  \vspace{4pt}\noindent{\color{poporange}\rule{36pt}{3.6pt}}
  \end{minipage}\par\nopagebreak\vspace{8pt}
}
\newcounter{popsub}
\newcommand{\courssub}[1]{%
  \stepcounter{popsub}%
  \par\vspace{9pt}\noindent{\popxbold\fontsize{11.5}{13}\selectfont\color{popink}{\color{poporange}\thepopsec.\thepopsub}\hspace{6pt}#1}\par\nopagebreak\vspace{4pt}
}

% ============================================================
% Boîtes pédagogiques — même recette : fond blanc, bordure épaisse colorée,
% ombre dure encre, badge plein. Titre optionnel [#1].
% ============================================================
\tcbset{popbase/.style={breakable,enhanced,colback=white,boxrule=2.3pt,arc=9pt,
  shadow={3pt}{-3pt}{0pt}{popink},left=14pt,right=14pt,top=11pt,bottom=11pt,before skip=11pt,after skip=15pt,
  fontupper=\popsemi\fontsize{10.6}{14.5}\selectfont\color{popink},
  fontlower=\popsemi\fontsize{10.6}{14.5}\selectfont\color{popink}}}
% Coche dessinée (le glyphe ✓ n'existe pas dans les polices du document)
\newcommand{\popcheck}[1]{\tikz[baseline=-0.5ex]\draw[#1,line width=1.6pt,line cap=round,line join=round](-0.13,0.0)--(-0.04,-0.09)--(0.13,0.1);}
\providecommand{\coche}{\popcheck{popgreen}}
\providecommand{\resultat}[1]{\par\smallskip\noindent\fcolorbox{popgreen}{popgreenL}{\parbox{\dimexpr\linewidth-2\fboxsep-2\fboxrule\relax}{\textbf{Résultat.} #1}}\par\smallskip}
\newcommand{\popboxhead}[3]{\poppill{#1}{#2}{white}\ifx\relax#3\relax\else\hspace{6pt}{\popxbold\fontsize{10.6}{12}\selectfont\color{popink}#3}\fi\par\vspace{7pt}}

\newtcolorbox{aretenir}[1][]{popbase,colframe=popgreen,before upper={\popboxhead{À retenir}{popgreen}{#1}}}
% Texte en allemand (document de lecture, dialogue, extrait) : cadre
% d'étude. Un titre optionnel (source, auteur) s'affiche dans l'en-tête.
\newtcolorbox{textebox}[1][]{popbase,colframe=popblue,before upper={\popboxhead{Text}{popblue}{#1}\begin{otherlanguage}{ngerman}},after upper={\end{otherlanguage}}}
% Marques ✓ / ✗ (forme correcte / fautive) souvent utilisées par les modèles
\providecommand{\cross}{\textcolor{poppink}{\ensuremath{\times}}}
\providecommand{\xmark}{\cross}\providecommand{\crossmark}{\cross}\providecommand{\cmark}{\ensuremath{\checkmark}}
% Ligne de réponse / justification à compléter (inventée par les modèles)
\providecommand{\justification}[1]{\par\noindent\textit{Justification~:} \dotfill\par}
% \textipa (package tipa non chargé) : la police ne couvre pas l'API -> simple passage
\providecommand{\textipa}[1]{#1}
% Soulignement (ulem non chargé) : \uline, \ul
\providecommand{\uline}[1]{\underline{#1}}\providecommand{\ul}[1]{\underline{#1}}
% \glossaire{\item ...} inventée par les modèles : liste de glossaire
\providecommand{\glossaire}[1]{\begin{itemize}#1\end{itemize}}
% \alert (beamer) inventée par les modèles : texte mis en évidence
\providecommand{\alert}[1]{\textbf{\textcolor{poppink}{#1}}}
% variantes ASCII d'ordinaux écrites en macro
\providecommand{\iere}{\textsuperscript{re}}\providecommand{\ieme}{\textsuperscript{e}}\providecommand{\ere}{\textsuperscript{re}}\providecommand{\eme}{\textsuperscript{e}}\providecommand{\er}{\textsuperscript{er}}\providecommand{\ie}{\textsuperscript{e}}
% Ordinaux français écrits en macro par les modèles : 1\ière, 3\ième
\newcommand{\ière}{\textsuperscript{re}}\newcommand{\ième}{\textsuperscript{e}}
% \exemple (inventée par les modèles) : étiquette « Exemple : »
\providecommand{\exemple}{\textbf{Exemple~:}\ }
% \square utilisable aussi hors mode math (cases à cocher)
\AtBeginDocument{\let\squareMath\square\renewcommand{\square}{\ensuremath{\squareMath}}}
% Mot ou phrase en allemand à l'intérieur d'un paragraphe français
\newcommand{\all}[1]{\ifmmode\text{\textit{#1}}\else\begingroup\selectlanguage{ngerman}\itshape #1\endgroup\fi}
\let\esp\all % alias toléré
\newtcolorbox{exemplebox}[1][]{popbase,colframe=poporange,before upper={\popboxhead{Exemple}{poporange}{#1}}}
\newtcolorbox{definition}[1][]{popbase,colframe=popblue,before upper={\popboxhead{Définition}{popblue}{#1}}}
\newtcolorbox{propriete}[1][]{popbase,colframe=poppurple,before upper={\popboxhead{Propriété}{poppurple}{#1}}}
\newtcolorbox{methode}[1][]{popbase,colframe=popgold,before upper={\popboxhead{Méthode}{popgold}{#1}}}
\newtcolorbox{attention}[1][]{popbase,colframe=poppink,before upper={\popboxhead{Attention}{poppink}{#1}}}
\newtcolorbox{experience}[1][]{popbase,colframe=popblue,colback=popblueL,before upper={\popboxhead{Expérience}{popblue}{#1}}}
% « Le savais-tu ? » : carte violette pleine (contexte, culture, Cameroun)
\newtcolorbox{savaistu}[1][]{popbase,colback=poppurple,colframe=popink,
  fontupper=\popsemi\fontsize{10.4}{14.2}\selectfont\color{white},
  before upper={\poppill{Le savais-tu\,?}{white}{poppurple}\ifx\relax#1\relax\else\hspace{6pt}{\popxbold\color{white}#1}\fi\par\vspace{7pt}}}
% Exercice résolu : énoncé en haut, \tcblower puis la solution sur fond crème
\newcounter{popexo}\newcounter{popexr}
\newtcolorbox{exoresolu}[1][]{popbase,colframe=poporange,colbacklower=popcream,
  segmentation style={popink,line width=1.2pt,dashed},
  before upper={\stepcounter{popexr}\popboxhead{Exercice résolu \thepopexr}{poporange}{#1}},
  before lower={\poppill{Solution}{popink}{popcream}\par\vspace{6pt}}}
% Exercice (non résolu en place) — la correction est dans la partie Corrigés
\newtcolorbox{exercice}[1][]{popbase,colframe=popink,
  before upper={\stepcounter{popexo}\popboxhead{Exercice \thepopexo}{popink}{#1}}}
% Corrigé
\newtcolorbox{corrige}[1][]{popbase,colframe=popgreen,colback=popgreenL,
  before upper={\popboxhead{Corrigé}{popgreen}{#1}}}
% Objectifs / prérequis / bilan
\newtcolorbox{objectifs}{popbase,colframe=popink,colback=poporangeL,
  before upper={\popboxhead{Objectifs de la leçon}{popink}{}}}
\newtcolorbox{prerequis}{popbase,colframe=popink,colback=popblueL,
  before upper={\popboxhead{Prérequis}{popblue}{}}}
\newtcolorbox{bilan}{popbase,colframe=popink,colback=white,boxrule=2.8pt,
  before upper={\popboxhead{Fiche bilan}{poporange}{L'essentiel à connaître pour l'examen}}}
% Figure encadrée avec légende
\newtcolorbox{popfigure}[1][]{enhanced,breakable=false,colback=white,colframe=popline,boxrule=1.4pt,arc=8pt,
  left=10pt,right=10pt,top=10pt,bottom=8pt,before skip=10pt,after skip=12pt,halign=center,
  lower separated=false,halign lower=center,
  fontlower=\popsemi\fontsize{9}{11.5}\selectfont\color{popmuted},
  #1}
\newcommand{\legende}[1]{\par\vspace{6pt}{\popsemi\fontsize{9}{11.5}\selectfont\color{popmuted}#1}}

% ============================================================
% Signature OnBuch+
% ============================================================
\newcommand{\poplogoinline}[1]{{\poplogo\fontsize{#1}{#1}\selectfont\color{popink}OnBuch\color{poporange}+}}
\newcommand{\signatureonbuch}{%
  \begin{tcolorbox}[enhanced,colback=popink,colframe=popink,boxrule=2.2pt,arc=10pt,
      drop shadow=poporange,left=14pt,right=14pt,top=10pt,bottom=10pt,before skip=0pt,after skip=0pt]
    \tikz[baseline=-0.7ex]\node[circle,fill=popgreen,draw=popcream,line width=1.3pt,minimum size=22pt,inner sep=0]{\popcheck{white}};%
    \hspace{9pt}\begin{minipage}[c]{0.62\linewidth}
      {\popxbold\fontsize{12}{13}\selectfont\color{popcream}Signé\ \textbullet\ {\poplogo OnBuch\color{poporange}+}}\par\vspace{2pt}
      {\raggedright\popsemi\fontsize{8}{10.5}\selectfont\color{popline}\MakeUppercase{Équipe pédagogique OnBuch+}\\\MakeUppercase{Conforme au programme officiel MINESEC}\par}
    \end{minipage}\hfill
    {\poplogo\fontsize{22}{22}\selectfont\color{popcream}OnBuch\color{poporange}+}
  \end{tcolorbox}%
}

% ============================================================
% Page de garde
% #1 titre · #2 matière · #3 niveau · #4 module · #5 n° leçon · #6 durée ·
% #7 description / objectifs (texte ou itemize)
% ============================================================
\newcommand{\pagegarde}[7]{%
  \thispagestyle{empty}
  \newgeometry{a4paper,margin=1.7cm,top=1.6cm,bottom=1.4cm}%
  \begin{tikzpicture}[remember picture,overlay]
    \fill[poporange!18] ([xshift=-3.2cm,yshift=-3.6cm]current page.north east) circle (4.6cm);
    \fill[poppurple!14] ([xshift=2.4cm,yshift=7.6cm]current page.south west) circle (3.8cm);
  \end{tikzpicture}%
  {\poplogo\fontsize{30}{30}\selectfont\color{popink}OnBuch\color{poporange}+}\hfill
  \raisebox{0.6ex}{\poppill{Cours signé \textbullet\ Premium}{popink}{popcream}}\par
  \vspace{1pt}\popsquiggle{poppurple}\par
  \vspace{8pt}
  \begin{tcolorbox}[enhanced,colback=poporange,colframe=popink,boxrule=2.8pt,arc=13pt,
      drop shadow=popink,left=18pt,right=18pt,top=12pt,bottom=13pt,after skip=10pt]
    \poppill{#2 \textbullet\ #3}{popink}{popcream}\hspace{5pt}\poppill{#4}{white}{popink}\par\vspace{8pt}
    {\popdisplay\fontsize{14}{14}\selectfont\color{popink}LEÇON #5}\par\vspace{4pt}
    {\raggedright\hyphenpenalty=10000\exhyphenpenalty=10000\popdisplay\fontsize{25}{27}\selectfont\color{popcream}\MakeUppercase{#1}\par}
    \vspace{8pt}
    {\popxbold\fontsize{9.5}{9.5}\selectfont\color{popink}\MakeUppercase{Durée conseillée : #6}}
  \end{tcolorbox}
  \begin{tcolorbox}[enhanced,colback=white,colframe=popink,boxrule=2.2pt,arc=10pt,
      drop shadow=popink,left=16pt,right=16pt,top=10pt,bottom=10pt,after skip=8pt]
    \poppill{Dans cette leçon}{poppurple}{white}\par\vspace{5pt}
    \popsemi\fontsize{9.3}{12.6}\selectfont\color{popink}#7
  \end{tcolorbox}
  \noindent\begin{minipage}[t]{0.487\linewidth}
    \begin{tcolorbox}[enhanced,colback=popgreen,colframe=popink,boxrule=2.2pt,arc=10pt,
        drop shadow=popink,left=11pt,right=11pt,top=10pt,bottom=10pt,height=3.1cm,valign=center]
      \tcbox[colback=white,colframe=popink,boxrule=1.3pt,arc=5pt,boxsep=0pt,nobeforeafter,
        left=3pt,right=3pt,top=3pt,bottom=3pt,valign=center]{\qrcode[height=1.9cm]{https://play.google.com/store/apps/details?id=com.luvvix.onbuch&hl=fr}}%
      \hspace{8pt}\begin{minipage}[c]{0.5\linewidth}
        {\popdisplay\fontsize{11}{12.5}\selectfont\color{white}\MakeUppercase{Télécharge OnBuch}}\par\vspace{4pt}
        {\raggedright\popsemi\fontsize{8.4}{11}\selectfont\color{white}Gratuit \textbullet\ Google Play\\Tuteur IA, annales, concours, résultats\par}
      \end{minipage}
    \end{tcolorbox}
  \end{minipage}\hfill
  \begin{minipage}[t]{0.487\linewidth}
    \begin{tcolorbox}[enhanced,colback=poppurple,colframe=popink,boxrule=2.2pt,arc=10pt,
        drop shadow=popink,left=11pt,right=11pt,top=10pt,bottom=10pt,height=3.1cm,valign=center]
      \tikz[baseline=-0.7ex]\node[circle,fill=white,draw=popink,line width=1.3pt,minimum size=1.6cm,inner sep=0]{\popdisplay\fontsize{24}{24}\selectfont\color{poppurple}+};%
      \hspace{8pt}\begin{minipage}[c]{0.6\linewidth}
        {\popdisplay\fontsize{11}{12.5}\selectfont\color{white}\MakeUppercase{Passe à OnBuch+}}\par\vspace{4pt}
        {\raggedright\popsemi\fontsize{8.4}{11}\selectfont\color{white}Tous les cours signés, Tuteur IA illimité, fiches PDF — onglet OnBuch+ de l'app\par}
      \end{minipage}
    \end{tcolorbox}
  \end{minipage}
  \vspace{8pt}
  \popcontenu
  \vfill
  \signatureonbuch
  \restoregeometry
  \newpage
}

% Rangée « Ce cours contient » (page de garde)
\newcommand{\poptile}[3]{% #1 couleur #2 gros texte #3 légende
  \begin{tcolorbox}[enhanced,colback=#1,colframe=popink,boxrule=2pt,arc=9pt,shadow={2.5pt}{-2.5pt}{0pt}{popink},
      left=6pt,right=6pt,top=9pt,bottom=9pt,halign=center,valign=center,height=1.75cm,width=\linewidth,nobeforeafter]
    {\popdisplay\fontsize{12.5}{13}\selectfont\color{white}\MakeUppercase{#2}}\par\vspace{3pt}
    {\centering\popsemi\fontsize{7.8}{9.5}\selectfont\color{white}#3\par}
  \end{tcolorbox}}
\newcommand{\popcontenu}{%
  \par\noindent{\popxboldls\fontsize{7.5}{7.5}\selectfont\color{popmuted}CE COURS SIGNÉ CONTIENT}\par\vspace{7pt}
  \noindent\begin{minipage}[t]{0.235\linewidth}\poptile{popblue}{Cours}{complet, illustré, pas à pas}\end{minipage}\hfill
  \begin{minipage}[t]{0.235\linewidth}\poptile{poporange}{Méthodes}{et exercices résolus}\end{minipage}\hfill
  \begin{minipage}[t]{0.235\linewidth}\poptile{poppink}{Exercices}{type examen + corrigés}\end{minipage}\hfill
  \begin{minipage}[t]{0.235\linewidth}\poptile{popgreen}{Fiche bilan}{l'essentiel à retenir}\end{minipage}\par}

% Sommaire
\newtcolorbox{sommaire}{enhanced,colback=white,colframe=popink,boxrule=2.2pt,arc=10pt,
  drop shadow=popink,left=18pt,right=18pt,top=13pt,bottom=13pt,before skip=6pt,after skip=20pt,
  before upper={\poppill{Sommaire}{poppurple}{white}\par\vspace{9pt}\popsemi\fontsize{10.6}{14.5}\selectfont\color{popink}}}

% Signature de fin de cours — poussée tout en bas de la dernière page
\newcommand{\findecours}{%
  \par\vfill\nopagebreak
  \begin{center}\popsquiggle{poporange}\end{center}\vspace{4pt}
  \signatureonbuch
}
\AtBeginDocument{\def\llbracket{\mathopen{[\mkern-3.5mu[}}\def\rrbracket{\mathclose{]\mkern-3.5mu]}}} % intervalles d'entiers (glyphe ⟦ absent de la police)
% Glyphes absents de Fira Math : redéfinis à partir de glyphes disponibles
\AtBeginDocument{%
  \def\setminus{\mathbin{\backslash}}\let\smallsetminus\setminus
  \def\vdots{\mathord{\vbox{\baselineskip3.2pt\lineskiplimit0pt\kern4pt\hbox{.}\hbox{.}\hbox{.}}}}%
  \def\ddots{\mathinner{\mkern1mu\raise6.5pt\hbox{.}\mkern2mu\raise3.6pt\hbox{.}\mkern2mu\raise0.7pt\hbox{.}\mkern1mu}}%
  \def\triangle{\mathord{\scalebox{0.9}{$\symup{\Delta}$}}}%
  \def\nabla{\mathord{\raisebox{\depth}{\scalebox{1}[-1]{$\symup{\Delta}$}}}}%
  \def\checkmark{\popcheck{popdark}}%
  \def\star{\mathbin{\ast}}\def\bigstar{\mathord{\ast}}%
  \def\diamond{\mathbin{\scalebox{0.6}{$\Diamond$}}}%
  \def\bigcup{\mathop{\mathchoice{\raisebox{-0.3em}{\scalebox{1.7}{$\cup$}}}{\scalebox{1.25}{$\cup$}}{\cup}{\cup}}\displaylimits}%
  \def\bigcap{\mathop{\mathchoice{\raisebox{-0.3em}{\scalebox{1.7}{$\cap$}}}{\scalebox{1.25}{$\cap$}}{\cap}{\cap}}\displaylimits}%
  \def\lesseqqgtr{\mathrel{\lessgtr}}\def\bigoplus{\mathop{\oplus}\displaylimits}%
}
\providecommand{\ssi}{si et seulement si} % abréviation fréquente
\AtBeginDocument{\providecommand{\wideparen}{\overparen}} % arc de cercle

\begin{document}

\pagegarde{Traduction et production de textes informatifs sur l'environnement, la santé et le bien-être}{Allemand}{Tle A}{Chapitre 3 : Environnement, santé et bien-être}{AL15}{2 h}{Cette leçon mixte du chapitre « Environnement, santé et bien-être » entraîne les élèves de Tle A à lire, produire et traduire des textes informatifs (articles, tracts) sur ces thèmes. Elle réinvestit le passif, l'irréel et les propositions relatives dans des tâches de thème et de version conformes aux exigences du Bac.
\vspace{4pt}
\begin{multicols}{2}\small
\begin{itemize}
\item À la fin de cette leçon, je sais identifier la structure d'un article informatif et d'un tract en allemand
\item À la fin de cette leçon, je sais utiliser le lexique thématique de l'environnement, de la santé et du bien-être dans un texte
\item À la fin de cette leçon, je sais rédiger un court article informatif ou un tract en allemand sur ces thèmes
\item À la fin de cette leçon, je sais réviser et employer correctement le passif (Vorgangspassiv) dans mes productions
\item À la fin de cette leçon, je sais employer les propositions relatives et l'irréel (Konjunktiv II) pour nuancer et argumenter
\item À la fin de cette leçon, je sais traduire de courtes phrases et un texte du français vers l'allemand (thème)
\end{itemize}\end{multicols}}

\begin{sommaire}
\begin{multicols}{2}
\begin{enumerate}
\item Texte support : article informatif et compréhension
\item Lexique thématique : environnement, santé, bien-être
\item Grammaire réinvestie I : le passif
\item Grammaire réinvestie II : relatives et irréel
\item Structure de l'article informatif et du tract
\item Traduction : thème et version guidés
\item Activité d'intégration
\item Méthodes et exercices résolus
\item Exercices
\item Corrigés des exercices
\item Fiche bilan
\end{enumerate}
\end{multicols}
\end{sommaire}

\begin{objectifs}
\begin{itemize}
\item À la fin de cette leçon, je sais identifier la structure d'un article informatif et d'un tract en allemand
\item À la fin de cette leçon, je sais utiliser le lexique thématique de l'environnement, de la santé et du bien-être dans un texte
\item À la fin de cette leçon, je sais rédiger un court article informatif ou un tract en allemand sur ces thèmes
\item À la fin de cette leçon, je sais réviser et employer correctement le passif (Vorgangspassiv) dans mes productions
\item À la fin de cette leçon, je sais employer les propositions relatives et l'irréel (Konjunktiv II) pour nuancer et argumenter
\item À la fin de cette leçon, je sais traduire de courtes phrases et un texte du français vers l'allemand (thème)
\item À la fin de cette leçon, je sais traduire un texte court de l'allemand vers le français (version)
\item À la fin de cette leçon, je sais relire et corriger ma production (orthographe, majuscules, ordre des mots, cas)
\end{itemize}
\end{objectifs}

\begin{prerequis}
\begin{itemize}
\item Le passif (Präsens, Präteritum, Perfekt) vu en Première
\item Les propositions relatives (der/die/das, déclinaison) vues en Première
\item Le Konjunktiv II (irréel, conseils, souhaits) vu en Première
\item Le vocabulaire de base de l'environnement et de la santé du chapitre 3
\item L'ordre des mots dans la phrase allemande (place du verbe, subordonnées)
\end{itemize}
\end{prerequis}

\newpage

\coursec{Texte support : article informatif et compréhension}

\courssub{Situation d'accroche et problématique}

À Yaoundé, un jeune volontaire d'une ONG écologique reçoit un e-mail de son partenaire allemand de Munich : il doit rédiger un court article informatif sur la gestion des déchets plastiques dans son quartier et traduire un tract de sensibilisation à l'hygiène destiné à une campagne commune Cameroun-Allemagne. En Allemagne, en Autriche et en Suisse, les journaux scolaires et les associations publient régulièrement des articles et des tracts sur l'environnement et la santé : la presse germanophone accorde une place centrale à ces sujets, et savoir lire, commenter, produire et traduire ce type de texte est une compétence exigée à l'examen.

\textbf{Problématique :} Comment structurer un texte informatif sur l'environnement, la santé et le bien-être ? Quelles formes grammaticales (passif, subordonnées relatives, irréel) faut-il réinvestir ? Cette leçon vous donne les outils pour produire et traduire des textes informatifs convaincants. Nous commençons par la lecture active d'un article modèle.

\courssub{Lecture du texte support}

Lisez d'abord le texte en silence, sans dictionnaire, pour saisir l'idée générale. Puis relisez-le en vous aidant du glossaire.

\begin{textebox}[Plastikmüll in der Stadt – ein wachsendes Problem]
\textbf{Plastikmüll in der Stadt – ein wachsendes Problem}

\emph{Ein Informationsartikel aus einer Schülerzeitung in Yaoundé}

In vielen Städten Kameruns wird zu viel Plastik weggeworfen. Auf den Märkten, in den Straßen und sogar in den Flüssen sieht man überall leere Plastikflaschen und alte Tüten. Die Flüsse werden verschmutzt, und die Gesundheit der Menschen wird gefährdet. Wenn das Regenwasser nicht mehr abfließen kann, weil die Abflüsse durch Plastik blockiert werden, entstehen Überschwemmungen. In diesen stehenden Gewässern vermehren sich Mücken, die Krankheiten wie Malaria übertragen.

Experten sagen, dass der Müll getrennt und recycelt werden muss. In Deutschland, Österreich und der Schweiz wird der Abfall seit vielen Jahren sortiert: Papier, Glas, Plastik und Bioabfall haben jeweils eigene Container. Diese Systeme, die von den Gemeinden finanziert werden, könnten auch in afrikanischen Städten eingeführt werden, wenn die Bevölkerung gut informiert würde.

Wenn jeder Bürger weniger Plastik benutzte, wäre die Umwelt besser geschützt. Man könnte zum Beispiel Stofftaschen statt Plastiktüten benutzen und Trinkwasser in wiederverwendbaren Flaschen kaufen. Eine Kampagne, die von Jugendlichen einer Umweltorganisation organisiert wurde, informiert jetzt die Bevölkerung von Yaoundé. Die Jugendlichen verteilen Flugblätter auf den Märkten und erklären den Händlern, wie der Müll richtig getrennt wird. „Wir wollen, dass unsere Stadt sauberer und gesünder wird“, sagt eine junge Aktivistin.

Das Problem ist groß, aber es ist nicht unlösbar. Jeder Einzelne kann etwas tun: weniger wegwerfen, mehr wiederverwenden, andere informieren. Die Zukunft unserer Umwelt liegt in unseren Händen.
\end{textebox}

\begin{definition}[Glossaire allemand-français du texte]
\begin{itemize}
\item \all{der Müll (Sg.)} : les déchets, les ordures
\item \all{wegwerfen (wirft weg, warf weg, hat weggeworfen)} : jeter (verbe à particule séparable)
\item \all{die Plastikflasche, -n} : la bouteille en plastique
\item \all{die Tüte, -n} : le sachet, le sac en plastique
\item \all{verschmutzen} : polluer, salir
\item \all{gefährden} : mettre en danger
\item \all{der Abfluss, -üe} : l'égout, l'évacuation des eaux
\item \all{die Überschwemmung, -en} : l'inondation
\item \all{die Mücke, -n} : le moustique
\item \all{übertragen (überträgt, übertrug, hat übertragen)} : transmettre (une maladie)
\item \all{trennen} : trier, séparer
\item \all{recyceln} : recycler
\item \all{der Abfall, -äe} : le déchet
\item \all{die Gemeinde, -n} : la commune
\item \all{die Stofftasche, -n} : le sac en tissu
\item \all{wiederverwendbar} : réutilisable
\item \all{die Kampagne, -n} : la campagne (de sensibilisation)
\item \all{das Flugblatt, -äer} : le tract, le prospectus
\item \all{der Händler, -} : le commerçant, le vendeur
\item \all{unlösbar} : insoluble
\end{itemize}
\end{definition}

\courssub{Compréhension globale du texte}

Avant d'entrer dans les détails, il faut répondre à quatre questions fondamentales qui valent pour tout texte informatif.

\begin{exoresolu}[Questions de compréhension globale]
Répondez en français, avec vos propres mots :
\begin{enumerate}
\item Quel est le thème (\all{das Thema}) du texte ?
\item Quelle est la source (\all{die Quelle}) du texte ?
\item À qui s'adresse ce texte (\all{der Empfänger}) ?
\item Quelle est l'intention (\all{die Absicht}) de l'auteur ?
\end{enumerate}
\tcblower
\textbf{Solution détaillée :}
\begin{enumerate}
\item \textbf{Thème :} le problème croissant des déchets plastiques dans les villes camerounaises et ses conséquences sur l'environnement et la santé, ainsi que les solutions possibles (tri, recyclage, sensibilisation).
\item \textbf{Source :} un article informatif publié dans un journal scolaire (\all{eine Schülerzeitung}) à Yaoundé. Le chapeau sous le titre l'indique explicitement.
\item \textbf{Destinataire :} les lecteurs du journal scolaire, c'est-à-dire les élèves, les jeunes et plus largement la population urbaine ; le texte vise aussi les citoyens appelés à agir (\all{jeder Bürger}, \all{jeder Einzelne}).
\item \textbf{Intention :} informer sur un problème concret (faits, mécanismes), sensibiliser aux dangers pour la santé (malaria, inondations) et inciter à l'action : la conclusion est un appel (\all{Appell}) : \all{„Jeder Einzelne kann etwas tun.“}
\end{enumerate}
\end{exoresolu}

\begin{attention}[Ne traduisez pas mot à mot !]
Dans la compréhension de texte, le piège classique du francophone est de traduire le texte mot à mot au lieu de le comprendre. Méthode correcte :
\begin{itemize}
\item répondez d'abord \textbf{avec vos propres mots}, en français ou en allemand simple ;
\item appuyez chaque réponse sur une \textbf{citation courte} du texte entre guillemets : \all{„Die Flüsse werden verschmutzt.“} ;
\item ne recopiez jamais trois lignes entières : une citation justificative fait quelques mots, pas un paragraphe.
\end{itemize}
Exemple de mauvaise réponse : recopier tout le deuxième paragraphe pour prouver que le tri est nécessaire. Bonne réponse : « Les experts exigent le tri des déchets : \all{„der Müll [muss] getrennt und recycelt werden“}. »
\end{attention}

\courssub{Compréhension détaillée : vrai ou faux, citations, questions ouvertes}

\begin{exercice}[Vrai ou faux ? Justifiez par une citation courte]
\begin{enumerate}
\item On ne trouve du plastique que sur les marchés.
\item Les déchets plastiques peuvent provoquer des inondations.
\item Le tri des déchets existe depuis longtemps dans les pays germanophones.
\item La campagne de sensibilisation est organisée par le gouvernement.
\item Selon le texte, le problème est sans solution.
\end{enumerate}
\end{exercice}

\begin{corrige}[Exercice « Vrai ou faux »]
\begin{enumerate}
\item \textbf{Faux.} \all{„Auf den Märkten, in den Straßen und sogar in den Flüssen“} : le plastique est partout, pas seulement sur les marchés.
\item \textbf{Vrai.} \all{„weil die Abflüsse durch Plastik blockiert werden, entstehen Überschwemmungen“}.
\item \textbf{Vrai.} \all{„In Deutschland, Österreich und der Schweiz wird der Abfall seit vielen Jahren sortiert.“}
\item \textbf{Faux.} \all{„eine Kampagne, die von Jugendlichen einer Umweltorganisation organisiert wurde“} : ce sont des jeunes d'une ONG, pas le gouvernement.
\item \textbf{Faux.} \all{„Das Problem ist groß, aber es ist nicht unlösbar.“}
\end{enumerate}
\end{corrige}

\begin{exercice}[Questions ouvertes]
\begin{enumerate}
\item Quel lien le texte établit-il entre le plastique et la malaria ?
\item Quelles solutions concrètes sont proposées aux citoyens ?
\item Pourquoi l'auteur cite-t-il l'Allemagne, l'Autriche et la Suisse ?
\end{enumerate}
\end{exercice}

\begin{corrige}[Questions ouvertes]
\begin{enumerate}
\item Le plastique bloque les égouts, ce qui provoque des inondations ; dans les eaux stagnantes, les moustiques se multiplient et transmettent des maladies comme la malaria (\all{„In diesen stehenden Gewässern vermehren sich Mücken, die Krankheiten wie Malaria übertragen.“}).
\item Utiliser des sacs en tissu au lieu de sachets plastiques, acheter de l'eau dans des bouteilles réutilisables, trier les déchets, jeter moins, réutiliser davantage et informer les autres.
\item Ces pays servent de modèle : leurs systèmes de tri, financés par les communes, montrent qu'une solution existe et pourrait être introduite au Cameroun si la population était bien informée.
\end{enumerate}
\end{corrige}

\courssub{Repérage de la structure de l'article informatif}

Un article informatif germanophone suit une architecture très régulière, qu'il faut savoir repérer avant de la reproduire. Observons notre texte :

\begin{itemize}
\item \textbf{Le titre} (\all{der Titel}) : \all{„Plastikmüll in der Stadt – ein wachsendes Problem“}. Court, il annonce le sujet et le problème.
\item \textbf{Le chapeau / l'introduction} (\all{die Einleitung}) : le premier paragraphe pose le constat : trop de plastique est jeté, les fleuves sont pollués, la santé est menacée. Il répond aux questions : de quoi parle-t-on ? où ? pourquoi est-ce grave ?
\item \textbf{Le corps} (\all{der Hauptteil}) : il développe les faits (\all{die Fakten}) : mécanisme des inondations et des moustiques, comparaison avec les pays germanophones, exemple de la campagne des jeunes avec une citation directe.
\item \textbf{La conclusion / l'appel à l'action} (\all{der Schluss / der Appell}) : \all{„Jeder Einzelne kann etwas tun … Die Zukunft unserer Umwelt liegt in unseren Händen.“} Le texte se termine par une invitation à agir.
\end{itemize}

\begin{popfigure}
\begin{center}
\begin{tikzpicture}[node distance=0.55cm]
\node[draw=popblue, fill=popblueL, rounded corners, minimum width=9.5cm, minimum height=0.9cm, align=center] (t) {\textbf{der Titel} : court, annonce le sujet};
\node[draw=popgreen, fill=popgreenL, rounded corners, minimum width=9.5cm, minimum height=0.9cm, align=center, below=of t] (e) {\textbf{die Einleitung} : le constat (quoi ? où ? pourquoi ?)};
\node[draw=poporange, fill=poporangeL, rounded corners, minimum width=9.5cm, minimum height=0.9cm, align=center, below=of e] (h) {\textbf{der Hauptteil} : faits, chiffres, exemples, citations};
\node[draw=poppurple, fill=poppurpleL, rounded corners, minimum width=9.5cm, minimum height=0.9cm, align=center, below=of h] (s) {\textbf{der Schluss / Appell} : conclusion et appel à l'action};
\draw[-{Stealth}, thick, popdark] (t) -- (e);
\draw[-{Stealth}, thick, popdark] (e) -- (h);
\draw[-{Stealth}, thick, popdark] (h) -- (s);
\end{tikzpicture}
\end{center}
\legende{Structure d'un article informatif : Titel, Einleitung, Hauptteil, Schluss/Appell}
\end{popfigure}

\courssub{Structure du tract (Flugblatt / Handzettel)}

Le tract (\all{das Flugblatt} ou \all{der Handzettel}) est un texte court, visuel, destiné à être distribué dans la rue ou affiché. Sa structure diffère de l'article : il va droit au but.

\begin{propriete}[Structure type d'un tract d'information / sensibilisation]
\begin{enumerate}
\item \textbf{Accroche visuelle et titre choc} (\all{der Eyecatcher / die Schlagzeile}) : grand, coloré, souvent une question ou un ordre. Ex. : \all{„Stopp dem Plastikmüll!“}
\item \textbf{Problème condensé} (\all{das Problem in Kürze}) : 2--3 phrases, faits clés, souvent au passif. Ex. : \all{„Unsere Flüsse werden verschmutzt.“}
\item \textbf{Solution concrète / Consignes d'action} (\all{Handlungsaufforderungen}) : verbes à l'impératif ou \all{man} + modal / infinitif. Ex. : \all{„Trennen Sie Ihren Müll! Nutzen Sie Stofftaschen!“}
\item \textbf{Appel final / Slogan} (\all{der Slogan / der Appell}) : court, mémorisable. Ex. : \all{„Saubere Stadt – gesunde Zukunft.“}
\item \textbf{Coordonnées / QR-Code / Logo} (\all{Kontakt / Veranstalter}) : qui organise ? où s'informer ?
\end{enumerate}
\end{propriete}

\begin{exemplebox}[Exemple de tract (modèle pour production)]
\begin{textebox}[FLUGBLATT: Sauberes Yaoundé – Wir handeln jetzt!]
\textbf{STOPP DEM PLASTIKMÜLL!}

\emph{Unsere Flüsse werden durch Plastik blockiert.}
\emph{Stehendes Wasser bringt Mücken und Malaria.}
\emph{Ihre Gesundheit ist in Gefahr!}

\textbf{Was SIE tun können:}
\begin{itemize}
\item Trennen Sie Müll: Plastik, Papier, Bio!
\item Nutzen Sie Stofftaschen statt Plastiktüten!
\item Trinken Sie Wasser aus Mehrwegflaschen!
\item Werden Sie aktiv: Informieren Sie Nachbarn und Händler!
\end{itemize}

\textbf{Saubere Stadt – Gesunde Zukunft.}

\emph{Organisiert von: Jugend-Umwelt-Netzwerk Yaoundé}
\emph{Kontakt: umwelt@yaounde.cm | QR-Code: [Scan mich]}
\end{textebox}
\end{exemplebox}

\begin{attention}[Article vs Tract : différences clés]
\begin{center}
\begin{tabularx}{\linewidth}{|X|X|X|}
\hline
\rowcolor{popblueL} \textbf{Critère} & \textbf{Artikel (Article)} & \textbf{Flugblatt (Tract)} \\
\hline
Longueur & Moyenne/Longue & Très courte (1 page) \\
\hline
Ton & Objectif, informatif & Direct, incitatif, visuel \\
\hline
Verbes & Passif, indicatif, Konj. II & \textbf{Impératif}, \all{man soll/muss}, infinitif \\
\hline
Structure & Titel – Einleitung – Hauptteil – Schluss & Schlagzeile – Problem – Lösungen – Slogan – Kontakt \\
\hline
Public & Lecteurs assis (journal) & Passants pressés (rue, marché) \\
\hline
\end{tabularx}
\end{center}
\end{attention}

\courssub{Repérage des formes grammaticales du chapitre dans le texte}

Le texte support a été construit pour réinvestir les trois points de grammaire du chapitre. Cherchons-les.

\textbf{1. Le passif} (\all{das Passiv}) : il est omniprésent dans les textes informatifs, car on s'intéresse à l'action, pas à l'auteur.

\begin{exemplebox}[Le passif dans le texte]
\begin{itemize}
\item \all{„In vielen Städten Kameruns wird zu viel Plastik weggeworfen.“} = Dans beaucoup de villes du Cameroun, on jette trop de plastique. (Präsens)
\item \all{„Die Flüsse werden verschmutzt.“} = Les fleuves sont pollués. (Präsens)
\item \all{„die Gesundheit der Menschen wird gefährdet“} = la santé des gens est mise en danger. (Präsens)
\item \all{„der Müll [muss] getrennt und recycelt werden“} = les déchets doivent être triés et recyclés. (Passif modal Präsens)
\item \all{„eine Kampagne, die von Jugendlichen organisiert wurde“} = une campagne qui a été organisée par des jeunes (Passif Präteritum dans relative).
\end{itemize}
Remarquez la forme : auxiliaire \all{werden} conjugué + participe II ; l'agent est introduit par \all{von} + datif. Le français préfère souvent « on » actif là où l'allemand utilise le passif.
\end{exemplebox}

\textbf{2. Les subordonnées relatives} (\all{die Relativsätze}) : elles enrichissent un nom et placent le verbe conjugué en fin de proposition.

\begin{exemplebox}[Les relatives dans le texte]
\begin{itemize}
\item \all{„Mücken, die Krankheiten wie Malaria übertragen“} = des moustiques qui transmettent des maladies comme la malaria (relatif nominatif pluriel \all{die}, antécédent \all{Mücken}).
\item \all{„Diese Systeme, die von den Gemeinden finanziert werden“} = ces systèmes, qui sont financés par les communes (relatif nominatif pluriel \all{die}, antécédent \all{Systeme}, passif dans la relative).
\item \all{„eine Kampagne, die von Jugendlichen organisiert wurde“} = une campagne qui a été organisée par des jeunes (relatif féminin singulier \all{die}, antécédent \all{Kampagne}).
\end{itemize}
\cle{Règle d'or} : Le pronom relatif s'accorde en \textbf{genre} et en \textbf{nombre} avec son antécédent, mais son \textbf{cas} (Nominatif, Accusatif, Datif, Génitif) dépend de sa fonction \textbf{dans la proposition relative}.
\end{exemplebox}

\textbf{3. L'irréel} (\all{der Konjunktiv II}) : il exprime l'hypothèse, la condition irréelle ou le souhait.

\begin{exemplebox}[Le Konjunktiv II dans le texte]
\begin{itemize}
\item \all{„Wenn jeder Bürger weniger Plastik benutzte, wäre die Umwelt besser geschützt.“} = Si chaque citoyen utilisait moins de plastique, l'environnement serait mieux protégé. (Condition + Conséquence, tous deux au KII).
\item \all{„[Diese Systeme] könnten auch in afrikanischen Städten eingeführt werden, wenn die Bevölkerung gut informiert würde.“} = Ces systèmes pourraient aussi être introduits dans les villes africaines si la population était bien informée. (\all{könnten} = KII de \all{können}, \all{würde} + infinitif pour \all{informieren}).
\item \all{„Man könnte zum Beispiel Stofftaschen benutzen.“} = On pourrait par exemple utiliser des sacs en tissu. (KII de \all{können} + infinitif).
\end{itemize}
Formes à retenir pour le thème/version : \all{wäre} (serait), \all{benutzte} (utilisait / Präteritum = KII forme faible), \all{könnte} (pourrait), \all{würde} + infinitif (substitut universel du KII).
\end{exemplebox}

\begin{center}
\begin{tabularx}{\linewidth}{|X|X|X|}
\hline
\rowcolor{popblueL} \textbf{Forme grammaticale} & \textbf{Exemple du texte} & \textbf{Fonction dans l'article} \\
\hline
Passif Präsens & \all{Die Flüsse werden verschmutzt.} & Décrire un fait sans nommer l'auteur \\
\hline
Passif + modal & \all{Der Müll muss getrennt werden.} & Exprimer une nécessité, une règle \\
\hline
Passif Präteritum & \all{die von Jugendlichen organisiert wurde} & Rapporter un événement passé (dans relative) \\
\hline
Relativsatz (Nom.) & \all{Mücken, die Krankheiten übertragen} & Préciser le sujet (précision définitoire) \\
\hline
Relativsatz (Passiv) & \all{Systeme, die finanziert werden} & Préciser l'objet d'une action passive \\
\hline
Konjunktiv II (wäre) & \all{wäre die Umwelt besser geschützt} & Conséquence hypothétique \\
\hline
Konjunktiv II (könnte/würde) & \all{könnten ... eingeführt werden / würde informiert} & Proposition de solution, hypothèse \\
\hline
\end{tabularx}
\end{center}

\courssub{Carte mentale des thèmes du texte}

Le texte relie trois grands champs lexicaux du chapitre : l'environnement, la santé et le bien-être. Cette carte mentale vous servira de réserve de vocabulaire pour vos propres productions.

\begin{popfigure}
\begin{center}
\begin{tikzpicture}[mindmap, grow cyclic, every node/.style={concept, align=center},
root concept/.style={concept color=popblue!40, minimum size=2.6cm},
level 1/.style={level distance=3.4cm, sibling angle=120},
level 2/.style={level distance=2.4cm, sibling angle=45, concept color=popgreen!35}]
\node[root concept]{Plastikmüll in der Stadt}
child[concept color=popgreen!45]{ node {die Umwelt}
  child { node[align=center]{der Müll\\trennen} }
  child { node[align=center]{die Flüsse\\verschmutzen} }
  child { node[align=center]{recyceln\\wiederverwenden} }
}
child[concept color=poporange!45]{ node {die Gesundheit}
  child { node[align=center]{die Mücken\\Malaria} }
  child { node[align=center]{die Gesundheit\\gefährden} }
}
child[concept color=poppurple!40]{ node {das Wohlbefinden}
  child { node[align=center]{saubere\\Stadt} }
  child { node[align=center]{die Kampagne\\informieren} }
};
\end{tikzpicture}
\end{center}
\legende{Carte mentale lexicale : les trois thèmes du texte et leur vocabulaire}
\end{popfigure}

\courssub{Méthode : commenter un texte informatif}

\begin{methode}[Les quatre étapes du commentaire de texte informatif]
\begin{enumerate}
\item \textbf{Présenter le texte} : indiquez le type de texte (article informatif, tract, reportage), le thème, la source et le destinataire. Exemple : « Il s'agit d'un article informatif paru dans un journal scolaire de Yaoundé, qui traite du problème des déchets plastiques. »
\item \textbf{Dégager la structure et l'intention} : montrez le plan (Titel, Einleitung, Hauptteil, Schluss/Appell) et dites ce que l'auteur veut obtenir : informer, alerter, convaincre, faire agir.
\item \textbf{Relever le lexique et les procédés} : repérez le champ lexical dominant (ici : environnement et santé), les faits et chiffres, les citations directes, le passif (faits objectifs), les impératifs ou le Konjunktiv II (appel à l'action), les relatives (précision).
\item \textbf{Donner un avis argumenté} : le texte atteint-il son but ? Les arguments sont-ils convaincants ? Comparez éventuellement avec la situation de votre propre ville. Justifiez toujours votre jugement par un exemple précis du texte.
\end{enumerate}
\end{methode}

\begin{exoresolu}[Mini-commentaire guidé]
Présentez le texte \all{„Plastikmüll in der Stadt“} en trois ou quatre phrases, comme à l'oral de l'examen.
\tcblower
\textbf{Solution détaillée (modèle de réponse) :}

\all{Der Text ist ein informativer Artikel aus einer Schülerzeitung in Yaoundé. Das Thema ist das wachsende Problem des Plastikmülls in kamerunischen Städten und seine Folgen für Umwelt und Gesundheit. Der Artikel ist klar strukturiert: Nach dem Titel stellt die Einleitung das Problem dar, der Hauptteil gibt Fakten und Beispiele (Überschwemmungen, Malaria, die Kampagne der Jugendlichen), und der Schluss ist ein Appell an jeden Bürger. Der Autor benutzt oft das Passiv („Die Flüsse werden verschmutzt“), um objektiv zu informieren, und den Konjunktiv II („wäre die Umwelt besser geschützt“), um Lösungen vorzuschlagen. Ich finde den Text überzeugend, weil er nicht nur Probleme zeigt, sondern auch konkrete Lösungen gibt.}

Traduction : Le texte est un article informatif tiré d'un journal scolaire de Yaoundé. Le thème est le problème croissant des déchets plastiques dans les villes camerounaises et ses conséquences pour l'environnement et la santé. L'article est clairement structuré : après le titre, l'introduction présente le problème, le corps donne des faits et des exemples (inondations, malaria, campagne des jeunes), et la conclusion est un appel à chaque citoyen. L'auteur utilise souvent le passif pour informer objectivement et le Konjunktiv II pour proposer des solutions. Je trouve le texte convaincant parce qu'il ne montre pas seulement les problèmes, mais donne aussi des solutions concrètes.
\end{exoresolu}

\begin{aretenir}[L'essentiel de la section]
\begin{itemize}
\item Un article informatif suit le schéma : \textbf{Titel → Einleitung → Hauptteil (Fakten) → Schluss/Appell}.
\item Un tract suit le schéma : \textbf{Schlagzeile → Problem → Handlungsaufforderungen (Impératif) → Slogan → Kontakt}.
\item En compréhension écrite : répondez \textbf{avec vos propres mots} et justifiez par des \textbf{citations courtes}.
\item Le \textbf{passif} (\all{werden} + participe II) rend le texte objectif ; les \textbf{relatives} précisent les noms (accord genre/nombre, cas selon fonction) ; le \textbf{Konjunktiv II} exprime l'hypothèse et les solutions possibles.
\item Le commentaire se fait en quatre étapes : présenter, analyser la structure, relever les procédés, donner un avis argumenté.
\end{itemize}
\end{aretenir}

\begin{exercice}[À vous de jouer : repérage grammatical]
Relevez dans le texte support : deux passifs au Präsens, un passif au Präteritum, deux propositions relatives et une phrase au Konjunktiv II. Traduisez chaque exemple en français.
\end{exercice}

\begin{corrige}[Exercice de repérage]
\begin{enumerate}
\item \textbf{Passifs Präsens :} \all{„wird ... weggeworfen“} (on jette / est jeté), \all{„werden verschmutzt“} (sont pollués), \all{„wird gefährdet“} (est mise en danger), \all{„muss ... getrennt und recycelt werden“} (doit être trié et recyclé).
\item \textbf{Passif Präteritum :} \all{„wurde ... organisiert“} (a été organisée).
\item \textbf{Relatives :} \all{„die Krankheiten ... übertragen“} (qui transmettent), \all{„die von den Gemeinden finanziert werden“} (qui sont financées par les communes), \all{„die von Jugendlichen ... organisiert wurde“} (qui a été organisée par des jeunes).
\item \textbf{Konjunktiv II :} \all{„Wenn jeder Bürger weniger Plastik benutzte, wäre die Umwelt besser geschützt.“} (Si chaque citoyen utilisait moins de plastique, l'environnement serait mieux protégé).
\end{enumerate}
\end{corrige}

% ============================================================
% NOUVELLE PARTIE : TRADUCTION THÈME / VERSION
% ============================================================

\coursec{Traduction : thème et version guidés}

\courssub{Méthodologie de la traduction (Thème / Version)}

\begin{methode}[Étapes de la traduction spécialisée (Environnement / Santé)]
\begin{enumerate}
\item \textbf{Analyse de la phrase source} : identifiez le sujet, le verbe, les compléments, le temps, le mode (indicatif / subjonctif / conditionnel), la voix (actif / passif).
\item \textbf{Transposition grammaticale} :
      \begin{itemize}
      \item \textbf{Passif FR $\to$ DE} : \all{werden} + Part. II (processus) ou \all{sein} + Part. II (état). \all{On jette} $\to$ \all{es wird ... weggeworfen} / \all{werden ... weggeworfen}.
      \item \textbf{Conditionnel FR (si + imparfait, conditionnel présent) $\to$ DE} : \all{Konjunktiv II} (\all{wäre, hätte, könnte, würde} + inf.) dans les deux propositions. \all{Si je pouvais, je le ferais} $\to$ \all{Wenn ich könnte, würde ich es tun}.
      \item \textbf{Relatives FR $\to$ DE} : pronom relatif (\all{der, die, das} selon genre/nombre antécédent, cas selon fonction dans la relative), verbe \textbf{à la fin}.
      \item \textbf{Impératif FR (tract) $\to$ DE} : \all{Sie}-Form (\all{Trennen Sie!}) ou \all{man} + modal (\all{Man muss trennen}) ou infinitif (\all{Müll trennen!}).
      \end{itemize}
\item \textbf{Recherche lexicale précise} : utilisez le vocabulaire du chapitre (\all{der Abfall, die Mülltrennung, die Überschwemmung, die Mücke, die Malaria, das Wohlbefinden, die Kampagne, das Flugblatt}). Attention aux \cle{faux amis} : \all{das Gift} = le poison (pas le don), \all{die Gift} n'existe pas ; \all{das Giftmüll} = déchets toxiques.
\item \textbf{Relecture / Contrôle} : majuscules aux noms ? Umlaute (ä, ö, ü) et ß corrects ? Ordre des mots (verbe 2ème position / fin de subordonnée) ? Accords (article, adjectif, participe) ? Ponctuation allemande (virgule avant subordonnée).
\end{enumerate}
\end{methode}

\courssub{Thème guidé : Français $\to$ Allemand}

Traduisez les phrases suivantes en allemand en réinvestissant le passif, les relatives et le Konjunktiv II. Le vocabulaire nécessaire figure dans le glossaire ou la carte mentale.

\begin{exercice}[Thème : Environnement et Santé]
\begin{enumerate}
\item Dans notre ville, trop de déchets plastiques sont jetés dans les rues.
\item Les égouts sont bloqués par les sachets plastiques, ce qui cause des inondations.
\item Les moustiques qui transmettent la malaria se multiplient dans l'eau stagnante.
\item Si la population triait ses déchets, l'environnement serait mieux protégé.
\item Une campagne qui a été organisée par des jeunes informe les commerçants sur le tri.
\item On pourrait utiliser des sacs en tissu au lieu de sacs en plastique.
\item Il faut que les déchets soient recyclés. (Indice : impératif ou passif modal)
\end{enumerate}
\end{exercice}

\begin{corrige}[Thème : Solutions commentées]
\begin{enumerate}
\item \all{In unserer Stadt wird zu viel Plastikmüll auf den Straßen weggeworfen.} (Passif Präsens, \all{werden} + Part.II, sujet \all{Plastikmüll}).
\item \all{Die Abflüsse werden durch Plastiktüten blockiert, was Überschwemmungen verursacht.} (Passif Präsens \all{werden blockiert} ; relative continue avec \all{was} = chose toute la proposition antérieure, verbe \all{verursacht} à la fin).
\item \all{Die Mücken, die Malaria übertragen, vermehren sich in stehendem Wasser.} (Relative \all{die} nominatif pluriel, antécédent \all{Mücken}, verbe \all{übertragen} à la fin ; \all{sich vermehren} réfléchi).
\item \all{Wenn die Bevölkerung ihren Müll trennte, wäre die Umwelt besser geschützt.} (Konjunktiv II : \all{trennte} / \all{wäre} + Part.II \all{geschützt}. \all{ihren Müll} accusatif).
\item \all{Eine Kampagne, die von Jugendlichen organisiert wurde, informiert die Händler über die Mülltrennung.} (Relative \all{die} fém. sg., passif Präteritum \all{wurde organisiert} à la fin ; \all{informieren über} + accusatif).
\item \all{Man könnte statt Plastiktüten Stofftaschen benutzen.} (Konjunktiv II \all{könnte} + infinitif \all{benutzen} à la fin ; \all{statt} + génitif/datif \all{Plastiktüten}).
\item \all{Der Müll muss recycelt werden.} (Passif modal \all{muss ... werden} + Part.II). Ou version tract : \all{Recyceln Sie Ihren Müll!} (Impératif formel).
\end{enumerate}
\end{corrige}

\courssub{Version guidée : Allemand $\to$ Français}

Traduisez le court texte suivant en français naturel (style journalistique ou tract). Notez les structures grammaticales ciblées.

\begin{exercice}[Version : Un tract pour la santé]
\textbf{Textgrundlage:}
\all{„Gesunde Ernährung in der Schule – Ein Projekt für unser Wohlbefinden“}

\all{In vielen Kantinen wird zu viel Fast Food angeboten. Das Essen, das viel Fett und Zucker enthält, gefährdet die Gesundheit der Schüler. Wenn mehr frisches Obst und Gemüse gekocht würde, wären die Kinder leistungsfähiger. Ein Projekt, das von Eltern und Lehrern ins Leben gerufen wurde, fordert jetzt: „Gesundes Essen für alle!“ Die Kantinen sollen verpflichtet werden, regionale Produkte zu kaufen. Machen Sie mit!}
\end{exercice}

\begin{corrige}[Version : Proposition de traduction]
\textbf{Titre :} « Alimentation saine à l'école – Un projet pour notre bien-être »

\textbf{Traduction :}
« Dans de nombreuses cantines, on propose trop de restauration rapide (\all{Fast Food}). La nourriture, qui contient beaucoup de graisse et de sucre, met en danger la santé des élèves. Si l'on cuisinait plus de fruits et légumes frais, les enfants seraient plus performants. Un projet, qui a été lancé par des parents et des professeurs, exige maintenant : "Une alimentation saine pour tous !" Les cantines devraient être tenues (\all{sollen verpflichtet werden} = obligation passive) d'acheter des produits régionaux. Participez ! / Faites partie du mouvement ! »

\textbf{Analyse grammaticale du texte source :}
\begin{itemize}
\item \all{wird ... angeboten} : Passif Präsens (on propose / est proposé).
\item \all{das viel Fett und Zucker enthält} : Relative (antécédent \all{das Essen} neutre sg. $\to$ \all{das}, sujet dans la relative $\to$ nominatif, verbe \all{enthält} à la fin).
\item \all{würde ... gekocht} / \all{wären} : Konjunktiv II (hypothèse irréelle présente). \all{würde kochen} (forme faible) / \all{wären} (sein).
\item \all{das ... ins Leben gerufen wurde} : Relative (antécédent \all{Ein Projekt} neutre sg. $\to$ \all{das}), Passif Präteritum \all{wurde gerufen} (particule \all{ins Leben rufen} = lancer), verbe à la fin.
\item \all{sollen ... verpflichtet werden} : Passif modal Präsens (obligation administrative).
\item \all{Machen Sie mit!} : Impératif formel (\all{Sie}-Form).
\end{itemize}
\end{corrige}

% ============================================================
% NOUVELLE PARTIE : PRODUCTION ÉCRITE (RÉDACTION)
% ============================================================

\coursec{Production écrite : rédiger un article et un tract}

\courssub{Consignes d'examen (Type Bac)}

\begin{attention}[Attendus du Baccalauréat - Production Écrite]
\begin{itemize}
\item \textbf{Sujet imposé :} Environnement, Santé ou Bien-être (contexte Cameroun ou pays germanophones).
\item \textbf{Format imposé :} Article informatif (\all{Artikel}) OU Tract (\all{Flugblatt}).
\item \textbf{Longueur :} Environ 120--150 mots (Article) / 80--100 mots (Tract).
\item \textbf{Critères de notation :} Respect du format (structure), richesse lexicale (champs lexicaux du chapitre), correction grammaticale (passif, relatives, Konj. II, ordre des mots, majuscules), cohérence et cohésion (connecteurs logiques).
\end{itemize}
\end{attention}

\courssub{Boîte à outils : Connecteurs logiques pour texte informatif}

\begin{center}
\begin{tabularx}{\linewidth}{|X|X|X|}
\hline
\rowcolor{popblueL} \textbf{Fonction} & \textbf{Connecteurs allemands} & \textbf{Français} \\
\hline
Énumérer / Ajouter & \all{zunächst, außerdem, ferner, zudem, auch} & premièrement, de plus, en outre, aussi \\
\hline
Expliquer / Cause & \all{weil, da, denn, deshalb, daher, deswegen} & parce que, car, donc, c'est pourquoi \\
\hline
Conséquence / But & \all{damit, um ... zu, sodass} & pour que, afin de, de sorte que \\
\hline
Contraste / Concession & \all{aber, jedoch, zwar ... aber, obwohl} & mais, cependant, bien que \\
\hline
Exemple & \all{zum Beispiel, beispielsweise, etwa} & par exemple \\
\hline
Conclusion / Résumé & \all{zusammenfassend, kurz gesagt, insgesamt} & en résumé, en bref, au total \\
\hline
\end{tabularx}
\end{center}

\courssub{Sujet d'entraînement 1 : Article informatif}

\begin{exercice}[Rédaction d'un article (120--150 mots)]
\textbf{Sujet :} Vous êtes rédacteur au journal de votre lycée à Douala. Écrivez un article informatif sur le thème : \all{„Lärmbelästigung in der Stadt – eine unterschätzte Gefahr für die Gesundheit“} (Pollution sonore en ville – un danger sous-estimé pour la santé).

\textbf{Contraintes :}
\begin{itemize}
\item Structure : Titel, Einleitung, Hauptteil (faits : trafic, générateurs, bars ; conséquences : stress, surdité, troubles du sommeil ; solutions : zones calmes, lois, sensibilisation), Schluss/Appell.
\item Utilisez au moins : 2 passifs (différents temps/modaux), 1 relative, 1 Konjunktiv II.
\item Vocabulaire imposé : \all{der Lärm, die Lärmbelästigung, der Verkehr, der Generator, die Bar, der Schlaf, das Gehör, die Zone, das Gesetz, die Aufklärung}.
\end{itemize}
\end{exercice}

\begin{corrige}[Proposition de correction - Article]
\textbf{Lärmbelästigung in der Stadt – eine unterschätzte Gefahr für die Gesundheit}

\emph{Ein Artikel der Schülerzeitung „Lycee Info“, Douala}

In Douala wird der Lärmpegel immer höher. Der Verkehr, laute Generatoren und Musik aus Bars belasten die Nerven der Einwohner. Die Gesundheit wird ernsthaft gefährdet: Ständiger Lärm führt zu Stress, Schlafstörungen und dauerhaften Hörschäden. Besonders Kinder und alte Menschen leiden unter der Lärmbelästigung.

Experten fordern, dass strenge Gesetze erlassen und kontrolliert werden. In Deutschland gibt es Ruhezeiten, die von allen eingehalten werden müssen. Solche Regeln könnten auch in kamerunischen Städten eingeführt werden, wenn die Politik entschlossen handeln würde. Eine Kampagne, die von Schülern organisiert wurde, misst derzeit den Lärm in verschiedenen Vierteln und informiert die Bevölkerung über die Risiken. „Wir brauchen ruhige Zonen für unser Wohlbefinden“, sagt ein Schüler.

Lärm ist nicht nur lästig, er macht krank. Jeder kann helfen: Musik leiser stellen, Generatoren warten, Ruhezeiten respektieren. Eine gesunde Stadt ist eine leise Stadt.
\end{corrige}

\courssub{Sujet d'entraînement 2 : Tract de sensibilisation}

\begin{exercice}[Rédaction d'un tract (80--100 mots)]
\textbf{Sujet :} Votre association \all{„Gesundes Yaoundé“} organise une journée « Zéro Déchet Plastique » au Marché Central. Rédigez le tract (\all{Flugblatt}) qui sera distribué la veille.

\textbf{Contraintes :}
\begin{itemize}
\item Structure : Schlagzeile, Problem (2 lignes), Handlungsaufforderungen (3--4 puces à l'impératif ou \all{man}), Slogan final, Kontakt/QR-Code.
\item Utilisez l'impératif formel (\all{Sie}-Form) ou \all{Man muss / soll}.
\item Vocabulaire : \all{der Markt, die Plastiktüte, die Stofftasche, der Müll, trennen, wiederverwenden, die Umwelt, schützen}.
\end{itemize}
\end{exercice}

\begin{corrige}[Proposition de correction - Tract]
\begin{textebox}[FLUGBLATT: Tag „Null Plastikmüll“ – Markttag für die Umwelt!]
\textbf{STOPP DEN PLASTIKMÜLL AUF UNSEREM MARKT!}

\emph{Jeden Tag landen tausende Plastiktüten im Müll und in der Natur.}
\emph{Die Umwelt wird verschmutzt, unsere Gesundheit gefährdet.}

\textbf{Was SIE morgen tun können:}
\begin{itemize}
\item \textbf{Bringen Sie Ihre eigenen Stofftaschen mit!}
\item \textbf{Verzichten Sie auf Plastiktüten bei jedem Kauf!}
\item \textbf{Trennen Sie Ihren Müll direkt am Stand!}
\item \textbf{Informieren Sie andere Händler und Kunden!}
\end{itemize}

\textbf{Sauberer Markt – Gesündere Zukunft.}

\emph{Veranstalter: Verein „Gesundes Yaoundé“}
\emph{Kontakt: gesund@yaounde.cm | QR-Code: [Aktion scannen]}
\end{textebox}
\end{corrige}

\courssub{Grille d'auto-évaluation avant de rendre votre copie}

\begin{center}
\begin{tabularx}{\linewidth}{|X|X|}
\hline
\rowcolor{popblueL} \textbf{Critère} & \textbf{Vérification (cochez)} \\
\hline
Structure respectée (Titel/Einleitung/Hauptteil/Schluss OU Schlagzeile/Problem/Aktion/Slogan) & $\square$ \\
\hline
Vocabulaire du chapitre utilisé (min. 8 mots précis : Müll, trennen, recyceln, Gesundheit, Gefahr, Kampagne, Flugblatt, Lärm, Wohlbefinden...) & $\square$ \\
\hline
Grammaire : Au moins 1 Passif (Präsens/Präteritum/Modal) & $\square$ \\
\hline
Grammaire : Au moins 1 Relative correcte (genre/nombre/cas/verbe fin) & $\square$ \\
\hline
Grammaire : Au moins 1 Konjunktiv II (hypothèse/solution) & $\square$ \\
\hline
Ordre des mots : Verbe position 2 (principale) / Verbe fin (subordonnée) & $\square$ \\
\hline
Majuscules à TOUS les noms (Umlaute ä/ö/ü, ß) & $\square$ \\
\hline
Orthographe / Ponctuation (virgule avant subordonnée) & $\square$ \\
\hline
Connecteurs logiques (zunächst, deshalb, aber, zum Beispiel...) & $\square$ \\
\hline
Longueur respectée (mots comptés) & $\square$ \\
\hline
\end{tabularx}
\end{center}

\begin{savaistu}[Le saviez-vous ? L'Allemagne et la gestion des déchets]
L'Allemagne est souvent citée en modèle pour le tri des déchets (\all{Mülltrennung}). Le système du \all{Pfand} (consigne) sur les bouteilles (verre, PET, canettes) existe depuis 2003 : on paie un supplément (0,25 €) qu'on récupère en rapportant le contenant vide dans un automate de supermarché (\all{Pfandautomat}). Cela finance le recyclage et évite les déchets sauvages. Au Cameroun, des initiatives locales (comme \all{Plastic Waste Free Cameroon} ou les ramasseurs de bouteilles au Marché Central) s'en inspirent. Le mot \all{der Müll} vient du moyen haut allemand \all{mülle} « déchet, pourriture », apparenté à \all{malen} « moudre, broyer » (d'où \all{die Mühle} = le moulin).
\end{savaistu}

\begin{aretenir}[Synthèse finale de la leçon AL15]
\begin{itemize}
\item \textbf{Deux formats maîtrisés :} Article (Titel–Einleitung–Hauptteil–Schluss/Appell) et Tract (Schlagzeile–Problem–Aufforderungen–Slogan–Kontakt).
\item \textbf{Grammaire réinvestie :} Passif (objectivité, règles), Relatives (précision, accord genre/nombre, cas fonctionnel, verbe final), Konjunktiv II (hypothèses, solutions, politesse).
\item \textbf{Traduction :} Thème (FR$\to$DE) : attention passif \all{werden}, KII \all{wäre/könnte/würde}, relatives \all{der/die/das}. Version (DE$\to$FR) : repérer structures, rendre le sens naturel (pas mot à mot), \all{man} $\to$ « on », passif $\to$ « on » ou passif FR.
\item \textbf{Production :} Planifier $\to$ Rédiger $\to$ Relire (Grille d'auto-évaluation). Le respect des majuscules aux noms et de l'ordre des mots (verbe) est la première marque de compétence au Bac.
\end{itemize}
\end{aretenir}

\coursec{Lexique thématique : environnement, santé, bien-être}

Dans la section précédente, nous avons lu et commenté un article informatif consacré à une campagne de sensibilisation au tri des déchets à Yaoundé. Ce texte nous a permis de repérer en contexte de nombreux mots nouveaux : \all{die Umwelt}, \all{der Müll}, \all{die Gesundheit}, \all{sich erholen}\ldots{} Pour pouvoir à notre tour \emph{produire} et \emph{traduire} des textes sur l'environnement, la santé et le bien-être, il faut maintenant organiser ce vocabulaire, le fixer avec ses articles et ses pluriels, et apprendre à l'employer dans des constructions correctes. C'est l'objet de cette section : constituer une « boîte à outils lexicale » fiable, sans laquelle aucune rédaction ni aucune traduction de qualité n'est possible au Baccalauréat.

\courssub{Le champ lexical de l'environnement : die Umwelt}

\begin{definition}[die Umwelt]
\all{die Umwelt} (sans pluriel) désigne \emph{l'environnement}, c'est-à-dire le milieu naturel dans lequel vivent les êtres humains, les animaux et les plantes. Attention : c'est un nom \emph{féminin}, alors que le français « environnement » est masculin. On dit \all{die Umwelt schützen} (protéger l'environnement), \all{die Umwelt verschmutzen} (polluer l'environnement), \all{umweltfreundlich} (écologique, respectueux de l'environnement).
\end{definition}

Observons d'abord quelques phrases tirées de notre texte support et du type de phrases que l'on rencontre dans la presse germanophone :

\all{Man sollte den Müll trennen und recyceln.} « On devrait trier et recycler les déchets. »

\all{Die Verschmutzung der Flüsse ist ein großes Problem.} « La pollution des rivières est un grand problème. »

\all{Erneuerbare Energien werden immer wichtiger.} « Les énergies renouvelables deviennent de plus en plus importantes. »

\begin{propriete}[Vocabulaire essentiel : l'environnement]
Apprends chaque nom \emph{avec son article et son pluriel} — c'est la règle d'or du vocabulaire allemand :
\begin{itemize}
\item \all{die Umwelt} (---) : l'environnement ; \all{der Umweltschutz} (---) : la protection de l'environnement
\item \all{der Müll} (---) : les déchets, les ordures ; \all{die Mülltonne} (-n) : la poubelle
\item \all{die Verschmutzung} (-en) : la pollution ; \all{verschmutzen} : polluer
\item \all{die Energie} (-n) : l'énergie ; \all{erneuerbar} : renouvelable
\item \all{der Abfall} (Abfälle) : le déchet ; \all{die Flasche} (-n) : la bouteille
\item \all{recyceln} : recycler ; \all{trennen} : trier, séparer ; \all{schützen} : protéger ; \all{sparen} : économiser
\end{itemize}
\end{propriete}

\begin{attention}[der Müll ou die Abfälle ?]
\all{der Müll} est un nom \emph{singulier indénombrable} : on ne dit jamais \all{die Mülle}. Pour parler de « déchets » au pluriel, utilise \all{die Abfälle} (pluriel de \all{der Abfall}). De même, « trier les déchets » se dit \all{den Müll trennen} ou \all{die Abfälle trennen}, jamais avec un article indéfini pluriel.
\end{attention}

\courssub{Le champ lexical de la santé : die Gesundheit}

\begin{definition}[die Gesundheit]
\all{die Gesundheit} (---) signifie \emph{la santé}. C'est un nom abstrait féminin formé sur l'adjectif \all{gesund} (en bonne santé) avec le suffixe \all{-heit}, très productif en allemand (\all{die Krankheit}, \all{die Freiheit}, \all{die Sicherheit}). Tous les noms en \all{-heit} sont féminins : \emph{die}.
\end{definition}

\begin{propriete}[Vocabulaire essentiel : la santé]
\begin{itemize}
\item \all{die Gesundheit} (---) : la santé ; \all{gesund} : en bonne santé ; \all{krank} : malade
\item \all{die Krankheit} (-en) : la maladie ; \all{die Erkältung} (-en) : le rhume
\item \all{der Arzt} (Ärzte) / \all{die Ärztin} (-nen) : le médecin / la femme médecin
\item \all{die Behandlung} (-en) : le traitement (médical) ; \all{behandeln} : soigner, traiter
\item \all{die Impfung} (-en) : la vaccination ; \all{impfen} : vacciner
\item \all{vorbeugen} (+ Dat.) : prévenir (une maladie) ; \all{sich erholen} : se rétablir, se reposer
\item \all{das Krankenhaus} (Krankenhäuser) : l'hôpital ; \all{das Medikament} (-e) : le médicament
\end{itemize}
\end{propriete}

\begin{exemplebox}[Exemples traduits]
\all{Nach der Krankheit hat er sich schnell erholt.} « Après la maladie, il s'est vite rétabli. »

\all{Die Ärztin empfiehlt eine Impfung gegen die Grippe.} « Le médecin (femme) recommande une vaccination contre la grippe. »

\all{Regelmäßige Bewegung beugt vielen Krankheiten vor.} « L'exercice régulier prévient de nombreuses maladies. » (remarque : \all{vorbeugen} + datif, et la particule \all{vor-} se place en fin de phrase)
\end{exemplebox}

\courssub{Le champ lexical du bien-être : das Wohlbefinden}

\begin{definition}[das Wohlbefinden]
\all{das Wohlbefinden} (---) désigne \emph{le bien-être}, c'est-à-dire l'état de bien-être physique et mental d'une personne. C'est un mot composé typique de l'allemand : \all{wohl} (bien) + \all{das Befinden} (l'état, du verbe \all{sich befinden}, « se trouver, se sentir »). L'allemand adore ces composés ; le français doit souvent les traduire par plusieurs mots. On dit : \all{Sport und gesunde Ernährung sind wichtig für das Wohlbefinden.} « Le sport et une alimentation saine sont importants pour le bien-être. »
\end{definition}

\begin{propriete}[Vocabulaire essentiel : le bien-être]
\begin{itemize}
\item \all{das Wohlbefinden} (---) : le bien-être ; \all{sich wohlfühlen} : se sentir bien
\item \all{der Stress} (---) : le stress ; \all{gestresst} : stressé
\item \all{sich entspannen} : se détendre ; \all{die Entspannung} (-en) : la détente
\item \all{die Ernährung} (---) : l'alimentation ; \all{sich gesund ernähren} : manger sainement
\item \all{Sport treiben} : faire du sport ; \all{die Bewegung} (-en) : le mouvement, l'exercice physique
\item \all{gesund leben} : vivre sainement ; \all{der Schlaf} (---) : le sommeil
\end{itemize}
\end{propriete}

\begin{exemplebox}[Exemples traduits]
\all{Sport treiben ist gut für die Gesundheit.} « Faire du sport est bon pour la santé. » (remarque : l'infinitif substantivé \all{Sport treiben} est sujet, comme « faire du sport » en français.)

\all{Nach der Arbeit entspanne ich mich mit Musik.} « Après le travail, je me détends avec de la musique. »

\all{Zu viel Stress schadet dem Wohlbefinden.} « Trop de stress nuit au bien-être. » (remarque : \all{schaden} + datif : \all{dem Wohlbefinden}.)
\end{exemplebox}

\courssub{Carte mentale : trois champs lexicaux}

\begin{popfigure}
\begin{center}
\begin{tikzpicture}[
  col/.style={rectangle, rounded corners=4pt, draw=#1, fill=#1!12, align=center, font=\small, inner sep=4pt},
  head/.style={rectangle, rounded corners=4pt, draw=#1, fill=#1!35, align=center, font=\small\bfseries, inner sep=5pt}]
\node[head=popblue] (h1) at (0,0) {die Umwelt};
\node[col=popblue, align=center] at (0,-1.1){der Müll\\ die Verschmutzung};
\node[col=popblue, align=center] at (0,-2.3){recyceln, trennen\\ schützen, sparen};
\node[col=popblue, align=center] at (0,-3.5){die Energie\\ erneuerbar};
\node[head=popgreen] (h2) at (5.2,0) {die Gesundheit};
\node[col=popgreen, align=center] at (5.2,-1.1){die Krankheit\\ der Arzt / die Ärztin};
\node[col=popgreen, align=center] at (5.2,-2.3){die Behandlung\\ die Impfung};
\node[col=popgreen, align=center] at (5.2,-3.5){sich erholen\\ vorbeugen};
\node[head=poporange] (h3) at (10.4,0) {das Wohlbefinden};
\node[col=poporange, align=center] at (10.4,-1.1){der Stress\\ sich entspannen};
\node[col=poporange, align=center] at (10.4,-2.3){die Ernährung\\ gesund leben};
\node[col=poporange, align=center] at (10.4,-3.5){Sport treiben\\ der Schlaf};
\end{tikzpicture}
\end{center}
\legende{Carte mentale lexicale : trois champs, dix-huit mots clés à mémoriser avec leurs articles.}
\end{popfigure}

\courssub{Verbes à particule et constructions fréquentes}

Les textes sur l'environnement et la santé emploient massivement des \cle{verbes à particule séparable} et des verbes à \cle{complément prépositionnel}. Ce sont deux pièges classiques pour les francophones.

\begin{propriete}[Quatre constructions à maîtriser]
\begin{enumerate}
\item \all{wegwerfen} (wirft weg, warf weg, hat weggeworfen) : jeter. Verbe fort à particule séparable. \all{Man darf Plastik nicht einfach wegwerfen.} « On ne doit pas simplement jeter le plastique. »
\item \all{aufklären} (über + Akk.) : sensibiliser, informer. \all{Die Kampagne klärt die Jugendlichen über die Gefahren auf.} « La campagne sensibilise les jeunes aux dangers. »
\item \all{sich kümmern um} (+ Akk.) : s'occuper de. \all{Wir müssen uns um unsere Gesundheit kümmern.} « Nous devons nous occuper de notre santé. »
\item \all{achten auf} (+ Akk.) : faire attention à. \all{Achte auf eine gesunde Ernährung!} « Fais attention à ton alimentation ! »
\end{enumerate}
\end{propriete}

\begin{attention}[La particule à la fin !]
À l'indicatif présent ou prétérit, la particule séparable va \emph{en toute fin de phrase}, même si le complément est long : \all{Er wirft die leeren Flaschen in die gelbe Tonne weg.} « Il jette les bouteilles vides dans la poubelle jaune. » Ne colle jamais la particule au verbe conjugué : *\all{Er wegwirft\ldots} est une erreur grave. En revanche, à l'infinitif avec \all{zu}, le \all{zu} s'insère \emph{entre} la particule et le verbe : \all{wegzuwerfen}, \all{aufzuklären}.
\end{attention}

\courssub{Expressions utiles pour l'article et le tract}

Pour rédiger un article informatif ou un tract, quelques formules toutes faites donnent immédiatement un ton authentique :

\begin{center}
\begin{tabularx}{\linewidth}{|X|X|}
\hline
\rowcolor{popblueL} \textbf{Deutsch} & \textbf{Français} \\
\hline
\all{Es ist wichtig, dass man den Müll trennt.} & Il est important que l'on trie les déchets. \\
\hline
\all{Man sollte mehr Sport treiben.} & On devrait faire plus de sport. \\
\hline
\all{Ziel dieser Kampagne ist es, die Bevölkerung aufzuklären.} & Le but de cette campagne est de sensibiliser la population. \\
\hline
\all{Jeder kann etwas für die Umwelt tun.} & Chacun peut faire quelque chose pour l'environnement. \\
\hline
\all{Gesundheit ist unser größtes Gut.} & La santé est notre bien le plus précieux. \\
\hline
\all{Handeln Sie jetzt! / Machen Sie mit!} & Agissez maintenant ! / Participez ! \\
\hline
\end{tabularx}
\end{center}

\begin{attention}[Es ist wichtig, dass\ldots]
Après \all{dass}, le verbe conjugué va \emph{en dernière position} : \all{Es ist wichtig, dass alle Kinder geimpft \textbf{werden}.} « Il est important que tous les enfants soient vaccinés. » Le français garde l'ordre normal, l'allemand non : c'est la subordonnée, règle absolue. Remarque aussi le passif dans la subordonnée : participe \all{geimpft} + \all{werden} à la fin.
\end{attention}

\courssub{Faux amis et choix lexicaux}

\begin{attention}[Les faux amis du chapitre]
\begin{itemize}
\item \all{bekommen} signifie \emph{recevoir}, jamais « devenir » : « devenir » se dit \all{werden}. \all{Ich habe ein Geschenk bekommen.} « J'ai reçu un cadeau. » / \all{Er wird Arzt.} « Il devient médecin. »
\item \all{die Diät} (-en) désigne \emph{le régime alimentaire} au sens large (ce qu'on mange habituellement), pas seulement une « diète » d'amaigrissement : \all{eine gesunde Diät} = une alimentation saine.
\item \all{die Krankheit} est le nom (« la maladie ») ; \all{krank} est l'adjectif (« malade »). Ne traduis jamais « il est malade » par *\all{Er ist Krankheit} : dis \all{Er ist krank.}
\item \all{die Kur} (-en) = la cure (thermale), alors que « court » se dit \all{kurz}.
\end{itemize}
\end{attention}

\begin{savaistu}[Le « Gelbe Sack », star allemande]
L'Allemagne est championne du tri sélectif : le \all{Gelbe Sack} (le « sac jaune »), destiné aux emballages en plastique et en métal, est connu dans le monde entier. Chaque commune distribue ces sacs jaunes, et les habitants trient aussi le papier (conteneur bleu) et le verre (par couleur : blanc, vert, brun). Le mot \all{die Mülltrennung} (le tri des déchets) fait partie du quotidien — un beau mot composé à réutiliser dans tes productions !
\end{savaistu}

\courssub{Tableau récapitulatif et entraînement}

\begin{center}
\begin{tabularx}{\linewidth}{|l|X|X|}
\hline
\rowcolor{popblueL} \textbf{Thème} & \textbf{Deutsch} & \textbf{Français} \\
\hline
Environnement & \all{die Umwelt, der Müll, die Verschmutzung, recyceln, trennen, erneuerbar} & l'environnement, les déchets, la pollution, recycler, trier, renouvelable \\
\hline
Santé & \all{die Gesundheit, die Krankheit, der Arzt, die Behandlung, die Impfung, vorbeugen} & la santé, la maladie, le médecin, le traitement, la vaccination, prévenir \\
\hline
Bien-être & \all{das Wohlbefinden, der Stress, sich entspannen, die Ernährung, Sport treiben, gesund leben} & le bien-être, le stress, se détendre, l'alimentation, faire du sport, vivre sainement \\
\hline
Constructions & \all{wegwerfen, aufklären, sich kümmern um, achten auf + Akk.} & jeter, sensibiliser, s'occuper de, faire attention à \\
\hline
\end{tabularx}
\end{center}

\begin{exoresolu}[Traduction guidée : du français vers l'allemand]
Énoncé : traduis en allemand, en utilisant le lexique de cette section. 1) On devrait trier les déchets et protéger l'environnement. 2) Faire du sport est bon pour la santé. 3) Le médecin m'a conseillé de me détendre davantage. 4) Chacun doit s'occuper de son alimentation.
\tcblower
Solution détaillée :
\begin{enumerate}
\item \all{Man sollte den Müll trennen und die Umwelt schützen.} — « on » = \all{man} ; \all{sollte} (Konjunktiv II de \all{sollen}) exprime le conseil ; \all{den Müll} : accusatif masculin.
\item \all{Sport treiben ist gut für die Gesundheit.} — l'infinitif substantivé est sujet ; \all{gut für} + accusatif.
\item \all{Der Arzt hat mir geraten, mich mehr zu entspannen.} — verbe réfléchi \all{sich entspannen} ; infinitive avec \all{zu} : \all{mich mehr zu entspannen} (le \all{zu} s'insère dans le verbe à particule).
\item \all{Jeder muss sich um seine Ernährung kümmern.} — \all{sich kümmern um} + accusatif ; le réfléchi \all{sich} reste près du verbe conjugué, la préposition \all{um} introduit le complément.
\end{enumerate}
\end{exoresolu}

\begin{exercice}[À toi de jouer]
1) Complète avec l'article correct : \all{\ldots Umwelt, \ldots Müll, \ldots Gesundheit, \ldots Wohlbefinden, \ldots Impfung, \ldots Arzt}.
2) Traduis en français : \all{Die Kampagne klärt die Schüler über die Mülltrennung auf.}
3) Traduis en allemand : « Il est important que nous fassions attention à notre santé. »
\end{exercice}

\begin{corrige}[Exercice : lexique]
1) \all{die Umwelt, der Müll, die Gesundheit, das Wohlbefinden, die Impfung, der Arzt}.
2) « La campagne sensibilise les élèves au tri des déchets. » (particule \all{auf} en fin de phrase ; \all{über} + accusatif).
3) \all{Es ist wichtig, dass wir auf unsere Gesundheit achten.} — subordonnée avec \all{dass} : verbe \all{achten} en dernière position ; \all{achten auf} + accusatif.
\end{corrige}

Ce lexique, désormais structuré et enrichi de constructions typiques, va nous servir immédiatement : dans la section suivante, nous réinvestirons ces mots dans la grammaire du \cle{passif}, indispensable pour décrire les actions d'une campagne de santé ou d'environnement (\all{Der Müll wird getrennt}, \all{Die Kinder werden geimpft}).

\coursec{Grammaire réinvestie I : le passif}

Dans la section précédente, nous avons constitué un lexique riche sur l'environnement, la santé et le bien-être : \all{der Müll}, \all{die Umweltverschmutzung}, \all{die Gesundheit}, \all{die Impfung}... Mais un vocabulaire ne suffit pas pour rédiger un article informatif ou un tract. Il faut encore savoir \textbf{comment} les journalistes et les rédacteurs de campagnes formulent les faits. Or, en relisant le texte support de la section 1, une forme verbale revient sans cesse : le \cle{passif}. C'est \textbf{elle} que nous allons maintenant observer, comprendre et réinvestir.

\courssub{Observation : le passif dans le texte support}

Reprenons quatre phrases tirées de l'article informatif étudié en lecture :

\begin{exemplebox}[Phrases relevées dans le texte]
\begin{itemize}
\item \all{Jedes Jahr wird in Kamerun viel Plastikmüll weggeworfen.} — Chaque année, beaucoup de déchets plastiques sont jetés au Cameroun.
\item \all{Viele Flüsse werden durch Chemikalien verschmutzt.} — De nombreux fleuves sont pollués par des produits chimiques.
\item \all{Letzte Woche wurde in Yaoundé eine Gesundheitskampagne organisiert.} — La semaine dernière, une campagne de santé a été organisée à Yaoundé.
\item \all{Der Müll muss getrennt werden.} — Les déchets doivent être triés.
\end{itemize}
\end{exemplebox}

Observons attentivement. Dans chaque phrase, on retrouve deux éléments :

\begin{enumerate}
\item une forme conjuguée du verbe \all{werden} : \all{wird}, \all{werden}, \all{wurde}, \all{muss ... werden} ;
\item un \cle{participe passé} (Partizip II) placé \textbf{en fin de phrase} : \all{weggeworfen}, \all{verschmutzt}, \all{organisiert}, \all{getrennt}.
\end{enumerate}

Autre observation essentielle : dans ces phrases, on ne dit pas \emph{qui} jette, \emph{qui} pollue, \emph{qui} a organisé. L'agent de l'action est absent ou secondaire : ce qui compte, c'est le \textbf{fait}, le résultat, la procédure. C'est exactement le \textbf{registre} de l'article informatif et du tract.

\courssub{Formation du passif}

\begin{propriete}[Règle fondamentale du passif allemand]
Le passif allemand (\emph{Vorgangspassiv}) se forme avec :
\begin{center}
\textbf{werden (conjugué) + Partizip II en fin de phrase}
\end{center}
\begin{itemize}
\item \all{werden} se conjugue au temps voulu et occupe la place du verbe conjugué (2\up{e} position dans la phrase déclarative) ;
\item le Partizip II reste invariable et se place \textbf{à la fin} du groupe verbal ;
\item l'agent, s'il est exprimé, est introduit par \all{von} + \textbf{datif} : \all{Der Wald wird von der Regierung geschützt.} (La forêt est protégée par le gouvernement.)
\end{itemize}
\textbf{Exception capitale au Perfekt} : le participe passé de \all{werden} au passif est \all{worden} et non \all{geworden} : \all{Die Kampagne ist organisiert worden.} (La campagne a été organisée.) On dit \all{geworden} uniquement quand \all{werden} signifie « devenir » : \all{Er ist Arzt geworden.} (Il est devenu médecin.)
\end{propriete}

\begin{attention}[Ne pas confondre passif et futur]
Le verbe \all{werden} a deux emplois radicalement différents :
\begin{itemize}
\item \all{werden} + \textbf{infinitif} = \textbf{futur} : \all{Er wird kommen.} = Il viendra.
\item \all{werden} + \textbf{Partizip II} = \textbf{passif} : \all{Er wird gelobt.} = Il est félicité.
\end{itemize}
Le sens change totalement ! Comparez : \all{Der Patient wird operieren.} (Le patient va opérer — futur) et \all{Der Patient wird operiert.} (Le patient est opéré — passif). Pour un francophone, le réflexe de sécurité est simple : regardez toujours la \textbf{fin de la phrase} — infinitif ou participe passé ?
\end{attention}

\courssub{Tableau des formes du passif}

Voici les formes du passif du verbe \all{machen} (faire) à la 3\up{e} personne du singulier, puis avec un modal :

\begin{center}
\begin{tabular}{|l|l|l|}
\hline
\rowcolor{popblueL} \textbf{Temps} & \textbf{Forme} & \textbf{Exemple} \\
\hline
Präsens & wird gemacht & \all{Der Müll wird getrennt.} \\
\hline
Präteritum & wurde gemacht & \all{Die Bevölkerung wurde informiert.} \\
\hline
Perfekt & ist gemacht worden & \all{Die Kampagne ist organisiert worden.} \\
\hline
Futur & wird gemacht werden & \all{Der Wald wird geschützt werden.} \\
\hline
Modal + passif & muss gemacht werden & \all{Der Müll muss getrennt werden.} \\
\hline
\end{tabular}
\end{center}

\begin{popfigure}
\begin{tikzpicture}[font=\small]
\node[draw=popdark, fill=popblueL, rounded corners, align=center, minimum width=3.4cm, minimum height=1cm] (base) at (0,0) {\textbf{PASSIF}\\ werden + Partizip II};
\node[draw=popblue, fill=popblue!15, rounded corners, align=center] (pr) at (-5.2,2.2) {Präsens\\ \all{wird gemacht}};
\node[draw=poporange, fill=poporangeL, rounded corners, align=center] (pt) at (-2.6,2.2) {Präteritum\\ \all{wurde gemacht}};
\node[draw=popgreen, fill=popgreenL, rounded corners, align=center] (pf) at (0,2.2) {Perfekt\\ \all{ist gemacht worden}};
\node[draw=poppurple, fill=poppurpleL, rounded corners, align=center] (fu) at (2.6,2.2) {Futur\\ \all{wird gemacht werden}};
\node[draw=poppink, fill=poppinkL, rounded corners, align=center] (mo) at (5.2,2.2) {Modal\\ \all{muss gemacht werden}};
\draw[-{Stealth}, thick, popblue] (base) -- (pr);
\draw[-{Stealth}, thick, poporange] (base) -- (pt);
\draw[-{Stealth}, thick, popgreen] (base) -- (pf);
\draw[-{Stealth}, thick, poppurple] (base) -- (fu);
\draw[-{Stealth}, thick, poppink] (base) -- (mo);
\node[draw=popgold, fill=popgoldL, rounded corners, align=center] (ag) at (0,-2) {Agent : \all{von} + Datif\\ \all{Der Fluss wird von Fabriken verschmutzt.}};
\draw[-{Stealth}, thick, popgold] (base) -- (ag);
\end{tikzpicture}
\legende{Carte mentale : les formes du passif allemand et l'expression de l'agent}
\end{popfigure}

\begin{exemplebox}[Exemples traduits à tous les temps]
\begin{itemize}
\item \all{Der Müll wird getrennt.} — Les déchets sont triés. (Präsens)
\item \all{Die Bevölkerung wurde informiert.} — La population a été informée. (Präteritum)
\item \all{Viele Bäume sind gefällt worden.} — De nombreux arbres ont été abattus. (Perfekt)
\item \all{Neue Krankenhäuser werden gebaut werden.} — De nouveaux hôpitaux seront construits. (Futur)
\item \all{Der Wald muss geschützt werden.} — La forêt doit être protégée. (Modal)
\item \all{Kinder sollen gegen Polio geimpft werden.} — Les enfants doivent être vaccinés contre la polio. (Modal \all{sollen})
\end{itemize}
\end{exemplebox}

\courssub{Comparaison français / allemand : le passif et la traduction de « on »}

Le passif est \textbf{beaucoup plus fréquent en allemand qu'en français}. Là où le français préfère le pronom indéfini « on » à la voix active, l'allemand choisit souvent soit \all{man} à l'actif, soit le passif, surtout dans les textes informatifs, les notices, les consignes et les tracts.

\begin{center}
\begin{tabularx}{\linewidth}{|X|X|X|}
\hline
\rowcolor{popblueL} \textbf{Français} & \textbf{Deutsch (actif)} & \textbf{Deutsch (passif)} \\
\hline
On trie les déchets. & \all{Man trennt den Müll.} & \all{Der Müll wird getrennt.} \\
\hline
On a informé la population. & \all{Man hat die Bevölkerung informiert.} & \all{Die Bevölkerung wurde informiert.} \\
\hline
On doit protéger la forêt. & \all{Man muss den Wald schützen.} & \all{Der Wald muss geschützt werden.} \\
\hline
Ici, on parle allemand. & \all{Hier spricht man Deutsch.} & \all{Hier wird Deutsch gesprochen.} \\
\hline
On vaccine les enfants. & \all{Man impft die Kinder.} & \all{Die Kinder werden geimpft.} \\
\hline
\end{tabularx}
\end{center}

\begin{methode}[Traduire « on » au thème (français $\rightarrow$ allemand)]
\begin{enumerate}
\item \textbf{Identifier} le complément d'objet direct de la phrase française : dans « On trie les déchets », le COD est « les déchets ».
\item \textbf{Choisir la stratégie} :
  \begin{itemize}
  \item solution active (toujours possible) : \all{man} + verbe actif $\rightarrow$ \all{Man trennt den Müll.}
  \item solution passive (registre informatif) : le COD français devient le \textbf{sujet nominatif} allemand $\rightarrow$ \all{Der Müll wird getrennt.}
  \end{itemize}
\item \textbf{Vérifier l'accord} : \all{werden} s'accorde avec le nouveau sujet : \all{Die Flüsse werden verschmutzt.} (pluriel !)
\item \textbf{Placer le Partizip II en fin de phrase}, même si des compléments s'intercalent : \all{Der Müll wird in vielen Städten getrennt.}
\item \textbf{Traduire « par »} de l'agent par \all{von} + datif : « Les fleuves sont pollués par les usines » $\rightarrow$ \all{Die Flüsse werden von den Fabriken verschmutzt.}
\end{enumerate}
\end{methode}

\courssub{Le passif sans agent : la marque du texte informatif}

Dans les articles, les tracts et les campagnes d'information, l'agent est le plus souvent \textbf{omis}, car il est inconnu, évident ou sans importance. L'allemand pousse même la logique jusqu'au \cle{passif impersonnel}, sans aucun sujet :

\begin{exemplebox}[Passif impersonnel]
\begin{itemize}
\item \all{Es wird informiert.} — On informe. / Une information est donnée.
\item \all{Es wird geimpft.} — On vaccine. (séance de vaccination en cours)
\item \all{In der Schule wird über Umweltschutz gesprochen.} — À l'école, on parle de protection de l'environnement.
\item \all{Hier wird nicht geraucht.} — Il est interdit de fumer ici. (litt. « ici, on ne fume pas »)
\end{itemize}
Notez que les verbes \textbf{intransitifs} (sans COD) peuvent aussi se mettre au passif impersonnel en allemand, ce qui est impossible en français : \all{Es wird getanzt.} (On danse.)
\end{exemplebox}

\begin{aretenir}[Emplois typiques du passif dans notre chapitre]
Le passif est le temps-roi des textes sur l'environnement, la santé et le bien-être, car il permet de présenter :
\begin{itemize}
\item des \textbf{faits objectifs} : \all{Jedes Jahr werden Millionen Bäume gefällt.} (Chaque année, des millions d'arbres sont abattus.)
\item des \textbf{procédures et consignes} : \all{Der Abfall muss getrennt werden.} (Les ordures doivent être triées.)
\item des \textbf{campagnes d'information} : \all{Die Bevölkerung wurde über die Gefahren des Rauchens informiert.} (La population a été informée des dangers du tabac.)
\item des \textbf{interdictions} : \all{Plastiktüten dürfen nicht mehr verkauft werden.} (Les sachets plastiques ne doivent plus être vendus.)
\end{itemize}
\end{aretenir}

\courssub{Entraînement guidé}

\begin{exoresolu}[Thème : du français à l'allemand]
Traduisez en allemand, en utilisant le passif : « Chaque année, beaucoup de déchets sont jetés dans la mer. La population doit être informée. Une campagne a été organisée par la mairie. »
\tcblower
\textbf{Solution détaillée.}
\begin{enumerate}
\item « beaucoup de déchets sont jetés » : le COD français « déchets » devient sujet nominatif : \all{viel Müll} (singulier en allemand) ; Präsens du passif de \all{wegwerfen} (verbe à particule séparable, Partizip II : \all{weggeworfen}) : \all{Jedes Jahr wird viel Müll ins Meer weggeworfen.}
\item « La population doit être informée » : modal \all{müssen} + passif à l'infinitif (\all{informiert werden}) : \all{Die Bevölkerung muss informiert werden.}
\item « Une campagne a été organisée par la mairie » : passé $\rightarrow$ Perfekt du passif (\all{ist ... worden}) ; « par la mairie » = \all{von der Stadtverwaltung} (datif) : \all{Eine Kampagne ist von der Stadtverwaltung organisiert worden.}
\end{enumerate}
Points de vigilance : \all{worden} et non \all{geworden} ; participe en fin de phrase ; \all{von} + datif (\all{der Stadtverwaltung}).
\end{exoresolu}

\begin{exercice}[À vous de jouer]
Mettez les phrases suivantes au passif, au temps indiqué entre parenthèses :
\begin{enumerate}
\item \all{Man verschmutzt die Flüsse.} (Präsens)
\item \all{Man hat die Kinder geimpft.} (Perfekt)
\item \all{Man muss den Regenwald schützen.} (Modal \all{müssen})
\item \all{Man organisierte eine Kampagne in Douala.} (Präteritum)
\end{enumerate}
Puis traduisez en allemand : « On doit protéger la santé des enfants. » (deux solutions possibles).
\end{exercice}

\begin{corrige}[Exercice « À vous de jouer »]
\begin{enumerate}
\item \all{Die Flüsse werden verschmutzt.} — sujet pluriel, donc \all{werden}.
\item \all{Die Kinder sind geimpft worden.} — Perfekt du passif : \all{sein} + Partizip II + \all{worden}.
\item \all{Der Regenwald muss geschützt werden.} — le modal se conjugue, le passif reste à l'infinitif \all{geschützt werden} en fin de phrase.
\item \all{Eine Kampagne wurde in Douala organisiert.} — Präteritum : \all{wurde}.
\end{enumerate}
Traduction : solution active \all{Man muss die Gesundheit der Kinder schützen.} ; solution passive \all{Die Gesundheit der Kinder muss geschützt werden.} Attention au génitif \all{der Kinder} (« des enfants »).
\end{corrige}

\begin{attention}[Les quatre pièges du passif pour les francophones]
\begin{enumerate}
\item \textbf{Confondre futur et passif} : \all{Er wird operieren.} (il va opérer) $\neq$ \all{Er wird operiert.} (il est opéré). Regardez la fin de la phrase !
\item \textbf{Écrire \all{geworden} au lieu de \all{worden}} au Perfekt du passif : on dit \all{Die Kampagne ist organisiert worden.} et jamais \emph{geworden}.
\item \textbf{Oublier la place du Partizip II} : il se place en toute fin de phrase, après tous les compléments : \all{Der Müll wird jeden Tag in den Schulen getrennt.}
\item \textbf{Traduire « par » par \all{durch}} : l'agent du passif s'introduit par \all{von} + datif ; \all{durch} n'exprime que le moyen ou la cause : \all{Der Fluss wird durch Chemikalien verschmutzt.} (cause) mais \all{Der Fluss wird von den Fabriken verschmutzt.} (agent).
\end{enumerate}
\end{attention}

\begin{savaistu}[Le passif, langue des affiches et des tracts]
Dans les pays germanophones, les panneaux officiels adorent le passif impersonnel : \all{Hier wird nicht geparkt.} (stationnement interdit), \all{Es wird gebeten, den Müll zu trennen.} (vous êtes priés de trier vos déchets). En Suisse, en Autriche comme en Allemagne, les campagnes de santé publique — vaccination, lutte contre le tabac, tri des déchets — sont rédigées presque entièrement au passif, car cette forme donne au message un ton objectif, neutre et impersonnel : exactement le style que l'on attend de vous dans un article informatif ou un tract au baccalauréat.
\end{savaistu}

Ainsi armés pour manier le passif, nous pourrons, dans la section suivante, enrichir nos phrases grâce aux propositions relatives et exprimer l'hypothèse et le souhait grâce à l'irréel — deux outils indispensables pour nuancer un article sur l'environnement ou la santé.

\coursec{Grammaire réinvestie II : relatives et irréel}

Dans la section précédente, nous avons révisé le passif, structure typique des textes informatifs (\all{Die Kampagne wurde von Jugendlichen organisiert.}). Mais un article de presse ne se contente pas de phrases courtes : il enchaîne les informations grâce aux \cle{propositions relatives}, et il formule conseils et hypothèses grâce au \cle{Konjunktiv II}. Ces deux outils, déjà rencontrés en classe de Première, sont indispensables pour produire et traduire des textes sur l'environnement, la santé et le bien-être au niveau Terminale. C'est l'objet de cette deuxième révision grammaticale.

\courssub{Observation : la relative dans le texte informatif}

\begin{textebox}[Extrait d'article — à observer]
\all{Eine Kampagne, die von Jugendlichen organisiert wurde, hat in Yaoundé großen Erfolg gehabt. Das Projekt, an dem die Schüler seit Monaten arbeiten, soll die Bevölkerung über die Gefahren des Plastiks informieren. Viele Einwohner, denen die Umwelt am Herzen liegt, haben an den Aktionen teilgenommen. Der Bürgermeister, dessen Stadt mit dem Müllproblem kämpft, hat die Initiative gelobt.}

\emph{Une campagne, qui a été organisée par des jeunes, a connu un grand succès à Yaoundé. Le projet, sur lequel les élèves travaillent depuis des mois, doit informer la population sur les dangers du plastique. Beaucoup d'habitants, à qui l'environnement tient à cœur, ont participé aux actions. Le maire, dont la ville lutte contre le problème des déchets, a salué l'initiative.}
\end{textebox}

\textbf{Observons.} Soulignons les pronoms relatifs : \all{die}, \all{dem}, \all{denen}, \all{dessen}. Trois constats s'imposent :

\begin{enumerate}
\item Le pronom relatif reprend un nom qui le précède (l'\cle{antécédent}) : \all{Kampagne} $\rightarrow$ \all{die} ; \all{Projekt} $\rightarrow$ \all{dem} ; \all{Einwohner} $\rightarrow$ \all{denen} ; \all{Bürgermeister} $\rightarrow$ \all{dessen}.
\item Le pronom relatif est toujours précédé d'une \textbf{virgule}.
\item Le verbe conjugué de la subordonnée se trouve \textbf{en toute fin} : \all{die von Jugendlichen organisiert \textbf{wurde}} ; \all{an dem die Schüler seit Monaten \textbf{arbeiten}}.
\end{enumerate}

La proposition relative est donc une subordonnée qui \emph{décrit ou précise un nom} et qui s'insère, entre virgules, au milieu de la phrase principale — ou se place après elle. C'est l'équivalent des propositions en « qui / que / dont / lequel » du français, mais avec une mécanique beaucoup plus stricte.

\courssub{Le tableau des pronoms relatifs}

\begin{propriete}[La règle d'or du pronom relatif]
Le pronom relatif obéit à \textbf{deux logiques différentes} :
\begin{itemize}
\item son \textbf{genre} et son \textbf{nombre} sont ceux de l'\cle{antécédent} (le nom auquel il se rapporte) ;
\item son \textbf{cas} dépend de sa \textbf{fonction dans la subordonnée} (sujet $\rightarrow$ nominatif ; COD $\rightarrow$ accusatif ; COI ou verbe datif $\rightarrow$ datif ; possession $\rightarrow$ génitif), ou de la \textbf{préposition} qui le gouverne.
\end{itemize}
De plus : virgule obligatoire devant le relatif, et verbe conjugué en fin de proposition.
\end{propriete}

\begin{center}
\begin{tabular}{|l|l|l|l|l|}
\hline
\rowcolor{popblueL} \textbf{Cas} & \textbf{Masculin} & \textbf{Féminin} & \textbf{Neutre} & \textbf{Pluriel} \\
\hline
Nominatif & der & die & das & die \\
\hline
Accusatif & den & die & das & die \\
\hline
Datif & dem & der & dem & denen \\
\hline
Génitif & dessen & deren & dessen & deren \\
\hline
\end{tabular}
\end{center}

\begin{attention}[Deux formes à mémoriser spécialement]
Au datif pluriel, on dit \all{denen} (et non « den ») : \all{die Jugendlichen, \textbf{denen} wir helfen}. Au génitif, \all{dessen} (masculin et neutre) et \all{deren} (féminin et pluriel) sont invariables et correspondent au français « dont / duquel » : \all{der Arzt, \textbf{dessen} Rat wichtig ist} (le médecin dont le conseil est important).
\end{attention}

\begin{aretenir}[Algorithme de choix du pronom relatif]
Pour ne jamais se tromper, pose-toi toujours ces deux questions dans l'ordre :
\begin{enumerate}
\item \textbf{Quel est l'antécédent ?} $\rightarrow$ Détermine le \textbf{genre} (m/f/n) et le \textbf{nombre} (sg/pl) du pronom.
\item \textbf{Quelle est la fonction du pronom DANS la relative ?} (Sujet ? Objet ? Après préposition \all{an/für/mit} ? Possession ?) $\rightarrow$ Détermine le \textbf{cas} (Nom/Akk/Dat/Gen).
\end{enumerate}
\emph{Exemple :} « La ville \textbf{dont} le maire… » $\rightarrow$ Antécédent = \all{die Stadt} (féminin sg) ; Fonction = possession (le maire \textbf{de} la ville) $\rightarrow$ Génitif féminin sg = \all{deren}. $\rightarrow$ \all{die Stadt, \textbf{deren} Bürgermeister…}
\end{aretenir}

\begin{popfigure}
\begin{center}
\begin{tikzpicture}[font=\small]
\node[draw=popdark, fill=popblueL, rounded corners, align=center, minimum width=4.2cm] (ante) at (0,0) {\textbf{Antécédent}\\ \all{der Artikel / die Kampagne}\\ \all{das Projekt / die Jugendlichen}};
\node[draw=popdark, fill=poporangeL, rounded corners, align=center, minimum width=4.2cm] (fkt) at (8.2,0) {\textbf{Fonction dans la relative}\\ sujet ? COD ? COI ?\\ possession ? préposition ?};
\node[draw=popdark, fill=popgreenL, rounded corners, align=center, minimum width=4.6cm] (form) at (4.1,-3.2) {\textbf{Forme du relatif}\\ der / die / das / den / dem /\\ der / denen / dessen / deren};
\draw[-{Stealth[length=3mm]}, very thick, popblue] (ante.south) -- (form.north west);
\draw[-{Stealth[length=3mm]}, very thick, poporange] (fkt.south) -- (form.north east);
\node[align=center, popmuted] at (4.1,1.3) {genre + nombre};
\node[align=center, popmuted] at (10.9,1.3) {cas};
\end{tikzpicture}
\end{center}
\legende{Les deux questions à se poser pour choisir le pronom relatif : le genre et le nombre viennent de l'antécédent, le cas vient de la fonction dans la subordonnée.}
\end{popfigure}

\begin{exemplebox}[Un exemple par cas, traduit]
\begin{itemize}
\item \textbf{Nominatif (sujet)} : \all{Der Artikel, \textbf{der} in der Zeitung steht, ist interessant.} « L'article qui paraît dans le journal est intéressant. »
\item \textbf{Accusatif (COD)} : \all{Der Artikel, \textbf{den} ich gelesen habe, ist interessant.} « L'article que j'ai lu est intéressant. »
\item \textbf{Datif (verbe datif)} : \all{Der Arzt, \textbf{dem} ich vertraue, arbeitet im Krankenhaus.} « Le médecin à qui je fais confiance travaille à l'hôpital. »
\item \textbf{Génitif (possession)} : \all{Die Kampagne, \textbf{deren} Erfolg alle überrascht, kostet wenig Geld.} « La campagne, dont le succès surprend tout le monde, coûte peu d'argent. »
\end{itemize}
\end{exemplebox}

\courssub{La relative avec préposition}

Quand le nom est introduit par une préposition dans la subordonnée, la préposition \textbf{précède} le pronom relatif et détermine son cas. C'est l'équivalent du français « pour lequel, sur lequel, à qui ».

\begin{exemplebox}[Préposition + relatif]
\begin{itemize}
\item \all{Das Projekt, \textbf{an dem} wir arbeiten, ist wichtig.} « Le projet sur lequel nous travaillons est important. » (\all{arbeiten an} + datif)
\item \all{Die Kampagne, \textbf{für die} wir uns engagieren, hat Erfolg.} « La campagne pour laquelle nous nous engageons a du succès. » (\all{sich engagieren für} + accusatif)
\item \all{Der Freund, \textbf{mit dem} ich Sport treibe, wohnt in Douala.} « L'ami avec qui je fais du sport habite à Douala. » (\all{mit} + datif)
\item \all{Die Krankheit, \textbf{gegen die} es noch keine Impfung gibt, ist gefährlich.} « La maladie contre laquelle il n'existe pas encore de vaccin est dangereuse. »
\end{itemize}
\end{exemplebox}

\begin{attention}[Pas de contraction, pas d'oubli]
Contrairement à l'allemand parlé avec les choses (\all{wo(r)-} : \all{woran, wofür}), dans un texte écrit soigné — et au Bac — on préfère la forme complète : \all{an dem}, \all{für die}. Et surtout, ne traduis jamais « dont » par \all{von dem} quand il s'agit d'une possession : « la ville dont le maire… » = \all{die Stadt, \textbf{deren} Bürgermeister…}.
\end{attention}

\courssub{Le Konjunktiv II : conseil, hypothèse, souhait}

\begin{textebox}[Conseils santé et environnement — à observer]
\all{Man sollte weniger Fleisch essen und mehr Sport treiben. Wenn wir weniger Plastik benutzten, wäre die Umwelt sauberer. Hätte ich mehr Zeit, würde ich jeden Tag zu Fuß zur Schule gehen. Ich wünschte, alle Menschen könnten in einer gesunden Umwelt leben!}

\emph{On devrait manger moins de viande et faire plus de sport. Si nous utilisions moins de plastique, l'environnement serait plus propre. Si j'avais plus de temps, j'irais chaque jour à l'école à pied. Je voudrais que tous les hommes puissent vivre dans un environnement sain !}
\end{textebox}

\begin{definition}[Le Konjunktiv II]
Le \cle{Konjunktiv II} est le mode de l'\textbf{irréel} : il exprime le conseil, l'hypothèse, le souhait, le rêve, la politesse atténuée. Il correspond globalement au \emph{conditionnel} français.
\end{definition}

\begin{propriete}[Formation du Konjunktiv II]
Deux systèmes coexistent :
\begin{itemize}
\item \textbf{La forme composée} : \all{würde} + infinitif, pour la grande majorité des verbes : \all{ich würde gehen, wir würden helfen}.
\item \textbf{Les formes simples} (obligatoires ou fortement préférées pour les verbes fréquents) : \all{sein $\rightarrow$ wäre} ; \all{haben $\rightarrow$ hätte} ; \all{können $\rightarrow$ könnte} ; \all{sollen $\rightarrow$ sollte} ; \all{müssen $\rightarrow$ müsste} ; \all{geben $\rightarrow$ es gäbe}.
\end{itemize}
Formation des formes simples : radical du \textbf{Präteritum} + Umlaut (pour les verbes forts) + terminaisons : \all{ich wäre, du wärst/wärest, er wäre, wir wären, ihr wärt/wäret, sie wären}.
\end{propriete}

\begin{center}
\begin{tabularx}{\linewidth}{|l|X|X|}
\hline
\rowcolor{popblueL} \textbf{Emploi} & \textbf{Deutsch} & \textbf{Français} \\
\hline
Conseil & \all{Man sollte mehr Sport treiben.} & On devrait faire plus de sport. \\
\hline
Hypothèse & \all{Wenn wir weniger Plastik benutzten, wäre die Luft sauberer.} & Si nous utilisions moins de plastique, l'air serait plus pur. \\
\hline
Souhait & \all{Wenn ich doch mehr Zeit hätte!} & Si seulement j'avais plus de temps ! \\
\hline
Rêve & \all{Ich würde gern in einer sauberen Stadt leben.} & J'aimerais vivre dans une ville propre. \\
\hline
Politesse & \all{Könnten Sie mir bitte helfen?} & Pourriez-vous m'aider, s'il vous plaît ? \\
\hline
\end{tabularx}
\end{center}

\begin{exemplebox}[Exemples traduits à réinvestir]
\begin{itemize}
\item \all{Man sollte weniger Fleisch essen.} « On devrait manger moins de viande. »
\item \all{Wenn ich Zeit hätte, würde ich mehr Sport treiben.} « Si j'avais le temps, je ferais plus de sport. »
\item \all{An deiner Stelle würde ich zum Arzt gehen.} « À ta place, j'irais chez le médecin. »
\item \all{Ohne Sport wäre das Leben langweilig.} « Sans sport, la vie serait ennuyeuse. »
\item \all{Wenn alle Menschen Fahrrad \textbf{führ}en, gäbe es weniger Abgase.} « Si tout le monde faisait du vélo, il y aurait moins de gaz d'échappement. »
\end{itemize}
\end{exemplebox}

\begin{attention}[Le conditionnel français n'est pas toujours « würde »]
Les francophones traduisent mécaniquement tout conditionnel par \all{würde} + infinitif. Or l'allemand préfère les \textbf{formes simples} pour les verbes fréquents : « je serais » = \all{ich \textbf{wäre}} (et non « ich würde sein ») ; « j'aurais » = \all{ich \textbf{hätte}} ; « je pourrais » = \all{ich \textbf{könnte}} ; « on devrait » = \all{man \textbf{sollte}}. À l'écrit, « ich würde sein » est considéré comme lourd, voire fautif dans un texte soigné. Réserve \all{würde} + infinitif aux autres verbes : \all{ich würde kommen, wir würden recyceln}.
\end{attention}

\begin{methode}[Construire une phrase d'hypothèse en trois étapes]
\begin{enumerate}
\item \textbf{Identifier la condition} : « si + imparfait » en français $\rightarrow$ \all{Wenn} + Konjunktiv II en allemand, verbe en fin : \all{Wenn wir weniger Plastik benutzten, …}
\item \textbf{Construire la conséquence} : la principale commence par le verbe (inversion !) : \all{…, \textbf{wäre} die Umwelt sauberer.} ou \all{…, \textbf{würden} wir gesünder leben.}
\item \textbf{Vérifier} : virgule entre les deux propositions, formes simples (\all{wäre, hätte, könnte}) privilégiées, verbe conjugué en fin de la subordonnée.
\end{enumerate}
\end{methode}

\begin{exoresolu}[Thème guidé : relatives et Konjunktiv II]
\textbf{Énoncé.} Traduisez en allemand : 1. « La campagne que les élèves ont organisée est un succès. » 2. « Le projet sur lequel nous travaillons concerne la santé. » 3. « Si j'avais le temps, je ferais plus de sport. » 4. « On devrait moins jeter de plastique. »
\tcblower
\textbf{Solution détaillée.}
\begin{enumerate}
\item Antécédent \all{die Kampagne} (féminin) ; dans la relative, « que » est COD $\rightarrow$ accusatif \all{die} ; verbe en fin : \all{Die Kampagne, \textbf{die} die Schüler organisiert \textbf{haben}, ist ein Erfolg.}
\item \all{arbeiten an} + datif ; \all{das Projekt} neutre $\rightarrow$ \all{an dem} : \all{Das Projekt, \textbf{an dem} wir \textbf{arbeiten}, betrifft die Gesundheit.}
\item Condition $\rightarrow$ Konjunktiv II ; « j'avais » = forme simple \all{hätte} ; conséquence avec inversion : \all{Wenn ich Zeit \textbf{hätte}, \textbf{würde} ich mehr Sport \textbf{treiben}.}
\item Conseil $\rightarrow$ \all{sollte} : \all{Man \textbf{sollte} weniger Plastik wegwerfen.}
\end{enumerate}
\end{exoresolu}

\begin{attention}[Les trois pièges classiques du francophone]
\begin{enumerate}
\item \textbf{La virgule} : en allemand, elle est \emph{obligatoire} devant toute relative, même quand le français n'en met pas : \all{Der Artikel, den ich lese, …} (jamais « Der Artikel den ich lese »).
\item \textbf{Le verbe en fin} : ne jamais calquer l'ordre français. « l'article que j'ai lu » = \all{der Artikel, den ich gelesen \textbf{habe}} — l'auxiliaire ferme la proposition.
\item \textbf{L'accord} : le relatif prend le genre de l'antécédent allemand, pas du mot français : « la campagne » est féminine, et \all{die Kampagne} aussi — mais « le problème » donne \all{das Problem, \textbf{das}…} (neutre !). Pense toujours à l'article allemand du nom.
\end{enumerate}
\end{attention}

\begin{savaistu}[Relatives et Konjunktiv II : des points au Bac]
Dans les productions écrites du baccalauréat, les correcteurs attribuent des points de « richesse et correction de la langue ». Une proposition relative bien construite (bon cas, virgule, verbe en fin) et un Konjunktiv II correctement employé (\all{Man sollte…}, \all{Wenn ich Zeit hätte, …}) font immédiatement la différence entre une copie moyenne et une excellente copie. Entraîne-toi à glisser au moins une relative et une phrase d'hypothèse ou de conseil dans chaque rédaction sur l'environnement ou la santé !
\end{savaistu}

\begin{exercice}[À toi de jouer]
1. Complète avec le pronom relatif correct : \all{a) Der Arzt, … ich vertraue, ist sehr kompetent. b) Die Stadt, … Bürgermeister die Kampagne unterstützt, heißt Yaoundé. c) Die Jugendlichen, … wir geholfen haben, danken uns. d) Das Buch, … ich lese, handelt von der Umwelt.}

2. Mets au Konjunktiv II : \all{a) Man treibt mehr Sport. (conseil avec \all{sollen}) b) Wir benutzen weniger Plastik. (hypothèse avec \all{wenn}) c) Ich habe mehr Zeit.}

3. Traduis : « Si nous recyclions plus, notre ville serait plus propre. »
\end{exercice}

\begin{corrige}[Exercice « À toi de jouer »]
1. a) \all{dem} (datif, \all{vertrauen} + datif) ; b) \all{deren} (génitif féminin, possession) ; c) \all{denen} (datif pluriel, \all{helfen} + datif) ; d) \all{das} (accusatif neutre).

2. a) \all{Man sollte mehr Sport treiben.} b) \all{Wenn wir weniger Plastik benutzten, …} c) \all{Wenn ich mehr Zeit hätte, …}

3. \all{Wenn wir mehr recycelten, wäre unsere Stadt sauberer.} (forme simple \all{wäre} préférée à « würde sein ».)
\end{corrige}

\coursec{Structure de l'article informatif et du tract}

Dans la section précédente, nous avons révisé les propositions relatives et l'irréel du présent (\all{Konjunktiv II}) : autant d'outils qui permettent de nuancer, d'expliquer et d'argumenter. Avant d'aborder la rédaction, rappelons les trois points de grammaire indispensables pour réussir cette leçon.

\courssub{Rappels grammaticaux indispensables}

\begin{propriete}[Le passif (\all{das Passiv}) : forme et emploi]
Le passif met l'accent sur l'action ou l'objet, pas sur l'agent. Il se forme avec \all{werden} + \all{Partizip II}.
\begin{center}
\begin{tabular}{|l|l|l|l|}
\hline
\rowcolor{popblueL} Temps & Actif & Passif (Vorgangspassiv) & Exemple (environnement) \\
\hline
Présent & \all{man sammelt} & \all{es wird gesammelt} & \all{Der Müll wird gesammelt.} \\
\hline
Prétérit & \all{man sammelte} & \all{es wurde gesammelt} & \all{Die Märkte wurden gereinigt.} \\
\hline
Parfait & \all{man hat gesammelt} & \all{es ist gesammelt worden} & \all{Mehrere Tonnen Müll sind gesammelt worden.} \\
\hline
Futur & \all{man wird sammeln} & \all{es wird gesammelt werden} & \all{Die Aktion wird wiederholt werden.} \\
\hline
\end{tabular}
\end{center}
\emph{Agent} : introduit par \all{von} + datif (\all{von der Stadtverwaltung}). \emph{Passif statal (Zustandspassif)} : \all{sein} + \all{Partizip II} (\all{Die Stadt ist sauber}).
\end{propriete}

\begin{propriete}[Les propositions relatives (\all{Relativsätze})]
Elles précisent un antécédent. Le pronom relatif (\all{der, die, das}, \all{welcher}, \all{dessen}) s'accorde en \textbf{genre} et \textbf{nombre} avec l'antécédent ; son \textbf{cas} dépend de sa fonction dans la relative. Le verbe conjugué va à la \textbf{fin}.
\begin{center}
\begin{tabular}{|l|l|l|l|}
\hline
\rowcolor{popblueL} Antécédent & Cas & Pronom & Exemple \\
\hline
\all{der Freiwillige} (Nom.) & Nominatif & \all{der} & \all{der Freiwillige, \textbf{der} Müll sammelt} \\
\hline
\all{den Müll} (Akk.) & Accusatif & \all{den} & \all{den Müll, \textbf{den} er sammelt} \\
\hline
\all{dem Viertel} (Dat.) & Datif & \all{dem} & \all{dem Viertel, \textbf{in dem} er wohnt} \\
\hline
\all{der Ärztin} (Gen.) & Génitif & \all{deren} & \all{die Ärztin, \textbf{deren} Rat hilft} \\
\hline
\end{tabular}
\end{center}
\emph{Attention :} préposition + pronom relatif $\rightarrow$ \all{wor-} + préposition si chose (\all{worüber}), pronom + préposition si personne (\all{bei dem}).
\end{propriete}

\begin{propriete}[L'irréel du présent (\all{Konjunktiv II} présent) : forme et emploi]
Exprime un souhait, une hypothèse irréelle ou une politesse. Forme : \all{würde} + infinitif (verbes faibles/mixtes) ou forme en \all{-e} + umlaut si possible (verbes forts : \all{ginge, käme, wäre, hätte}).
\begin{center}
\begin{tabular}{|l|l|l|}
\hline
\rowcolor{popblueL} Verbe & \all{Konjunktiv II} (ich/er/sie/es) & Exemple phrase \\
\hline
\all{sein} & \all{wäre} & \all{Wenn die Stadt sauberer \textbf{wäre}...} \\
\hline
\all{haben} & \all{hätte} & \all{Ich \textbf{hätte} gern mehr Grünflächen.} \\
\hline
\all{werden} & \all{würde} & \all{Er \textbf{würde} mitmachen.} \\
\hline
\all{können} & \all{könnte} & \all{Wir \textbf{könnten} die Umwelt schützen.} \\
\hline
\all{müssen} & \all{müsste} & \all{Man \textbf{müsste} mehr tun.} \\
\hline
\all{helfen} & \all{hülfe} / \all{würde helfen} & \all{Wenn er \textbf{hülfe}...} \\
\hline
\end{tabular}
\end{center}
Structure conditionnelle : \all{Wenn} + \all{Konjunktiv II} (si), proposition principale \all{Konjunktiv II} (alors). \all{Wenn jeder mithelfen \textbf{würde}, sähe die Stadt anders \textbf{aus}.}
\end{propriete}

\courssub{Observation : deux textes, un même sujet, deux registres}

\begin{textebox}[Article informatif — \all{„Yaoundé startet eine Woche für Gesundheit und Umwelt“}]
\textbf{Yaoundé startet eine Woche für Gesundheit und Umwelt}

\emph{Seit Montag findet in der kamerunischen Hauptstadt eine Aktionswoche statt, die Gesundheit und Umweltschutz verbindet. Organisiert wird sie von der Stadtverwaltung und mehreren Vereinen.}

Yaoundé. Wer am Montagmorgen durch das Stadtzentrum ging, konnte es nicht übersehen: Überall standen junge Freiwillige mit Handschuhen und Müllsäcken. Sie sammelten Plastikflaschen, die sonst die Straßen und Flüsse verschmutzen. „Unsere Stadt wird immer sauberer, aber es bleibt noch viel zu tun“, erklärte eine Ärztin, die an der Aktion teilnahm.

Die Aktionswoche wurde von der Stadtverwaltung und von mehreren Umweltvereinen organisiert. Zuerst wurden die Märkte gereinigt, außerdem fanden kostenlose Gesundheitsuntersuchungen statt. Viele Einwohner konnten ihren Blutdruck messen lassen und bekamen Ratschläge für eine gesunde Ernährung. Ein Experte erinnerte daran, dass regelmäßige Bewegung und wenig Zucker das Wohlbefinden deutlich verbessern.

Jedoch gibt es auch Probleme: Das Sammelsystem für den Müll funktioniert noch nicht in allen Stadtteilen. Deshalb fordern die Vereine mehr Container und eine bessere Information der Bevölkerung. Wenn jeder Einwohner mithelfen würde, sähe die Stadt schon in einem Jahr ganz anders aus.

Schließlich endete die Woche mit einem großen Fest auf dem Marktplatz, bei dem Schüler Theaterstücke zum Thema Umweltschutz zeigten. Die Organisatoren zogen eine positive Bilanz: Mehrere Tonnen Müll wurden gesammelt, und Hunderte Menschen wurden untersucht. Nächstes Jahr soll die Aktion wiederholt werden — vielleicht sogar in anderen Städten des Landes.
\end{textebox}

\textbf{Glossaire :}
\begin{itemize}
\item \all{die Aktionswoche, die Aktionswochen} (la semaine d'action)
\item \all{die Stadtverwaltung, die Stadtverwaltungen} (l'administration municipale)
\item \all{der Verein, die Vereine} (l'association)
\item \all{der Freiwillige, die Freiwilligen} (le bénévole — N-Deklination)
\item \all{der Müllsack, die Müllsäcke} (le sac poubelle)
\item \all{verschmutzen} (polluer)
\item \all{die Gesundheitsuntersuchung, die Gesundheitsuntersuchungen} (l'examen médical)
\item \all{der Blutdruck} (la tension artérielle)
\item \all{die Ernährung} (l'alimentation)
\item \all{das Wohlbefinden} (le bien-être)
\item \all{der Stadtteil, die Stadtteile} (le quartier)
\item \all{die Bilanz, die Bilanzen} (le bilan)
\item \all{untersuchen} (examiner)
\end{itemize}

\textbf{Questions de compréhension :}
\begin{enumerate}
\item \all{Wer organisiert die Aktionswoche?} (Qui organise la semaine d'action ?)
\item \all{Was fanden außer der Reinigung der Märkte noch statt?} (Qu'est-ce qui a eu lieu en plus du nettoyage des marchés ?)
\item \all{Welches Problem wird im Text genannt?} (Quel problème est mentionné ?)
\item Relevez dans le texte deux phrases au passif (un présent, un prétérit), une proposition relative avec préposition, et une phrase à l'irréel (\all{Konjunktiv II}). Analysez la structure de chacune.
\end{enumerate}

\begin{textebox}[Tract — \all{„Sauberes Yaoundé — mach mit!“}]
\textbf{SAUBERES YAOUNDÉ — MACH MIT!}

Willst du in einer sauberen und gesunden Stadt leben?

Dann komm am \textbf{Samstag, den 14. Juni, um 9 Uhr} auf den \textbf{Marktplatz}!

\begin{itemize}
\item Hilf mit! Sammle Müll in deinem Viertel!
\item Trenne deinen Müll! Plastik gehört nicht in den Fluss!
\item Bringe Handschuhe und gute Laune mit!
\end{itemize}

Kostenlose Gesundheitsuntersuchung für alle Teilnehmer!

\textbf{Kontakt: Verein „Grüne Stadt“, Tel. 690 00 00 00}

\textbf{GEMEINSAM SIND WIR STARK!}
\end{textebox}

\textbf{Glossaire :}
\begin{itemize}
\item \all{das Viertel, die Viertel} (le quartier)
\item \all{trennen, trennte, hat getrennt} (trier, séparer — verbe fort, particule séparable \all{trennen} n'est pas séparable, mais \all{wegwerfen} l'est)
\item \all{der Teilnehmer, die Teilnehmer} (le participant)
\item \all{mitbringen, brachte mit, hat mitgebracht} (apporter — verbe à particule séparable)
\end{itemize}

Comparons maintenant ces deux documents. L'article \emph{informe} : il rapporte des faits, cite des témoins, reste objectif. Le tract \emph{mobilise} : il interpelle le lecteur, donne des ordres amicaux, annonce un événement. C'est toute la différence entre le registre \cle{factuel} et le registre \cle{incitatif}.

\courssub{L'article informatif : structure et registre}

\begin{definition}[L'article informatif]
L'\cle{article informatif} (\all{der Zeitungsartikel}) est un texte journalistique qui rapporte des faits récents de manière objective. Il se compose de quatre parties : le \cle{titre} (\all{die Überschrift}), le \cle{chapeau} (\all{die Einleitung} / \all{der Vorspann}), le \cle{corps} (\all{der Hauptteil}) et la \cle{conclusion} (\all{der Schluss}).
\end{definition}

\begin{propriete}[Les quatre parties de l'article]
\begin{enumerate}
\item \textbf{Le titre} (\all{die Überschrift}) : court, accrocheur, souvent sans article. \all{„Yaoundé startet eine Woche für Gesundheit und Umwelt“}.
\item \textbf{Le chapeau} (\all{die Einleitung}) : un ou deux phrases en italique ou en gras qui répondent aux questions \all{Wer? Was? Wo? Wann?} (qui ? quoi ? où ? quand ?). C'est la règle des \cle{W-Fragen}.
\item \textbf{Le corps} (\all{der Hauptteil}) : deux à trois paragraphes de faits, avec des chiffres, des détails et des citations d'experts ou de témoins au discours direct : \all{„Unsere Stadt wird immer sauberer“, erklärte eine Ärztin.}
\item \textbf{La conclusion} (\all{der Schluss}) : un bilan (\all{die Bilanz}) ou une perspective, parfois un appel : \all{Nächstes Jahr soll die Aktion wiederholt werden.}
\end{enumerate}
\end{propriete}

\begin{attention}[Le registre de l'article : objectif et factuel]
Dans un article informatif, on n'écrit jamais \all{ich finde} ou \all{je pense}. Le journaliste rapporte : il utilise le passif (\all{Die Woche wurde organisiert}), des citations (\all{erklärte eine Ärztin}) et des formulations neutres. L'opinion, si elle existe, passe par le \all{Konjunktiv II} ou par les propos rapportés : \all{Wenn jeder mithelfen würde, sähe die Stadt anders aus.} (Si tout le monde participait, la ville aurait un autre visage.)
\end{attention}

\courssub{Le tract : structure et registre}

\begin{definition}[Das Flugblatt]
\all{Das Flugblatt, die Flugblätter} (aussi \all{der Flugzettel, die Flugzettel}) désigne le \cle{tract} : un feuillet distribué dans la rue, sur les marchés ou dans les écoles pour informer rapidement la population ou la mobiliser pour une action (manifestation, campagne de nettoyage, vaccination).
\end{definition}

\begin{propriete}[Les quatre éléments du tract]
\begin{enumerate}
\item \textbf{Le slogan} : une phrase-choc en grosses lettres, souvent avec un tiret ou un point d'exclamation : \all{„Sauberes Yaoundé — mach mit!“}
\item \textbf{Les impératifs} : phrases courtes à l'impératif pour interpeller directement : \all{Hilf mit!} \all{Trenne deinen Müll!} \all{Schütze die Umwelt!}
\item \textbf{Les informations pratiques} : date, heure, lieu, contact : \all{Samstag, den 14. Juni, um 9 Uhr, Marktplatz. Kontakt: Verein „Grüne Stadt“.}
\item \textbf{L'appel final} : une dernière phrase mobilisatrice : \all{Gemeinsam sind wir stark!} (Ensemble, nous sommes forts !)
\end{enumerate}
Les \cle{questions rhétoriques} (\all{Willst du in einer sauberen Stadt leben?}) et les \cle{chiffres marquants} (\all{Mehrere Tonnen Müll!}) renforcent l'effet.
\end{propriete}

\begin{exemplebox}[Impératifs de tract : exemples traduits]
\all{Hilf mit! Trenne deinen Müll!} = Participe ! Trie tes déchets !

\all{Komm am Samstag um 10 Uhr zum Marktplatz!} = Viens samedi à 10 h sur la place du marché !

\all{Schütze die Umwelt! Wirf keinen Müll auf die Straße!} = Protège l'environnement ! Ne jette pas de déchets dans la rue !

\all{Geh zum Arzt! Lass dich untersuchen!} = Va chez le médecin ! Fais-toi examiner !

Rappel : à l'impératif \all{du}, le verbe fort prend parfois un \emph{e} $\rightarrow$ \emph{i} : \all{helfen} $\rightarrow$ \all{hilf!} ; \all{werfen} $\rightarrow$ \all{wirf!} ; \all{geben} $\rightarrow$ \all{gib!} ; mais \all{schützen} $\rightarrow$ \all{schütze!} (verbe faible, sans changement), \all{trennen} $\rightarrow$ \all{trenne!}.
\end{exemplebox}

\begin{center}
\begin{tabularx}{\linewidth}{|X|X|X|}
\hline
\rowcolor{popblueL} Élément & Article informatif & Tract \\
\hline
But & informer (objectif) & mobiliser (incitatif) \\
\hline
Ouverture & titre + chapeau (W-Fragen) & slogan + question rhétorique \\
\hline
Phrases & longues, passif, citations & courtes, impératifs \\
\hline
Personne & 3\up{e} personne (il/on) & 2\up{e} personne (du/ihr) \\
\hline
Fin & bilan ou perspective & appel + infos pratiques \\
\hline
\end{tabularx}
\end{center}

\begin{popfigure}
\begin{tikzpicture}[
  box/.style={draw=popblue, thick, rounded corners, fill=popblueL, align=center, minimum width=4.6cm, minimum height=0.85cm, font=\small},
  box2/.style={draw=poporange, thick, rounded corners, fill=poporangeL, align=center, minimum width=4.6cm, minimum height=0.85cm, font=\small},
  title/.style={font=\bfseries}
]
\node[title, popblue] at (-2.8,4.6) {Der Zeitungsartikel};
\node[title, poporange] at (2.8,4.6) {Das Flugblatt};
\node[box, align=center] (a1) at (-2.8,3.6){Überschrift\\ (titre accrocheur)};
\node[box, align=center] (a2) at (-2.8,2.5){Einleitung\\ (Wer? Was? Wo? Wann?)};
\node[box, align=center] (a3) at (-2.8,1.4){Hauptteil\\ (Fakten, Zahlen, Zitate)};
\node[box, align=center] (a4) at (-2.8,0.3){Schluss\\ (Bilanz oder Ausblick)};
\node[box2, align=center] (b1) at (2.8,3.6){Slogan\\ (phrase-choc)};
\node[box2, align=center] (b2) at (2.8,2.5){Imperative\\ (Hilf mit! Trenne!)};
\node[box2, align=center] (b3) at (2.8,1.4){Praktische Infos\\ (Datum, Ort, Kontakt)};
\node[box2, align=center] (b4) at (2.8,0.3){Appell\\ (Gemeinsam sind wir stark!)};
\draw[-{Stealth}, thick, popblue] (a1) -- (a2);
\draw[-{Stealth}, thick, popblue] (a2) -- (a3);
\draw[-{Stealth}, thick, popblue] (a3) -- (a4);
\draw[-{Stealth}, thick, poporange] (b1) -- (b2);
\draw[-{Stealth}, thick, poporange] (b2) -- (b3);
\draw[-{Stealth}, thick, poporange] (b3) -- (b4);
\end{tikzpicture}
\legende{Structure comparée de l'article informatif et du tract}
\end{popfigure}

\courssub{Les connecteurs : la colonne vertébrale du texte}

Pour que le corps de l'article se lise facilement, les paragraphes doivent être reliés par des \cle{connecteurs} (\all{Konnektoren}). Attention à la place du verbe : après un connecteur en tête de phrase (position 1), le verbe conjugué vient en \textbf{deuxième position}, puis le sujet (inversion).

\begin{center}
\begin{tabularx}{\linewidth}{|l|X|X|}
\hline
\rowcolor{popblueL} Connecteur & Sens & Exemple \\
\hline
\all{zuerst} & d'abord & \all{Zuerst wurden die Märkte gereinigt.} (D'abord, les marchés ont été nettoyés.) \\
\hline
\all{außerdem} & de plus & \all{Außerdem fanden Untersuchungen statt.} (De plus, des examens ont eu lieu.) \\
\hline
\all{deshalb} & c'est pourquoi & \all{Deshalb fordern die Vereine mehr Container.} (C'est pourquoi les associations réclament plus de conteneurs.) \\
\hline
\all{jedoch} & cependant & \all{Jedoch gibt es auch Probleme.} (Cependant, il y a aussi des problèmes.) \\
\hline
\all{schließlich} & finalement & \all{Schließlich endete die Woche mit einem Fest.} (Finalement, la semaine s'est terminée par une fête.) \\
\hline
\all{einerseits / andererseits} & d'une part / d'autre part & \all{Einerseits ist Sport gesund, andererseits braucht man auch Ruhe.} (D'une part le sport est sain, d'autre part il faut aussi du repos.) \\
\hline
\end{tabularx}
\end{center}

\begin{attention}[Pièges des francophones avec les connecteurs]
\begin{itemize}
\item \all{Zuerst die Märkte wurden gereinigt.} est \textbf{FAUX} : après \all{zuerst}, inversion obligatoire $\rightarrow$ \all{Zuerst \textbf{wurden} die Märkte gereinigt.}
\item \all{Aber} (conjonction, position 0) et \all{jedoch} (adverbe, position 1) : \all{Es gibt Probleme, aber...} vs \all{Es gibt jedoch Probleme.} / \all{Jedoch gibt es Probleme.}
\item Ne confondez pas \all{deshalb} (conséquence, verbe 2\up{e} pos.) et \all{weil} (cause, verbe en fin) : \all{Es regnet, deshalb bleibe ich zu Hause.} / \all{Ich bleibe zu Hause, weil es regnet.}
\end{itemize}
\end{attention}

\courssub{Méthode de production guidée et grille de relecture}

\begin{methode}[Rédiger un article informatif en quatre étapes]
\begin{enumerate}
\item \textbf{Planifier} : notez les réponses aux \all{W-Fragen} (qui ? quoi ? où ? quand ?) et choisissez deux ou trois faits avec un chiffre ou une citation d'expert.
\item \textbf{Rédiger le titre et le chapeau} : le chapeau résume l'essentiel en une ou deux phrases (passif souvent utile : \all{Organisiert wird...}).
\item \textbf{Rédiger le corps} : deux à trois paragraphes reliés par \all{zuerst}, \all{außerdem}, \all{jedoch} ; intégrez une relative (\all{die/der/das...}) et un irréel si hypothèse (\all{Wenn... würde...}). Terminez par une conclusion avec bilan (\all{eine positive Bilanz ziehen}) ou appel.
\item \textbf{Relire} avec la grille de correction ci-dessous : majuscules, ordre des mots, cas, orthographe, ponctuation.
\end{enumerate}
Pour un \textbf{tract}, la démarche est la même, mais le plan devient : slogan $\rightarrow$ question rhétorique $\rightarrow$ impératifs (3-4) $\rightarrow$ date, lieu, contact $\rightarrow$ appel final.
\end{methode}

\begin{aretenir}[Grille de relecture — toujours vérifier]
\begin{enumerate}
\item \textbf{Majuscules} à \textbf{tous} les noms communs : \all{der \textbf{M}üll, die \textbf{G}esundheit, das \textbf{W}ohlbefinden, die \textbf{U}mwelt, die \textbf{S}tadt}.
\item \textbf{Ordre des mots} : verbe conjugué en 2\up{e} position ; inversion après connecteur en tête ; verbe en fin dans les subordonnées (\all{weil}, \all{dass}, \all{wenn}, relatives).
\item \textbf{Cas après les prépositions} : \all{nach dem Fest} (dat.), \all{mit einem Fest} (dat.), \all{für die Gesundheit} (akk.), \all{in der Stadt} (dat.), \all{in die Stadt} (akk. mouvement).
\item \textbf{Orthographe} : Umlaute (\all{ä, ö, ü}), \all{ß} (après voyelle longue/diphtongue) vs \all{ss} (après voyelle brève) ; particules séparables collées à l'infinitif (\all{mitzumachen}, \all{wegzuwerfen}).
\item \textbf{Ponctuation} : virgule avant toute subordonnée et avant l'infinitif avec \all{zu} ; point d'exclamation après l'impératif du tract.
\end{enumerate}
\end{aretenir}

\begin{attention}[La faute n° 1 des francophones]
En français, les noms communs prennent la minuscule ; en allemand, \textbf{TOUS} les noms prennent la majuscule, partout dans la phrase. À la relecture, surlignez chaque nom et vérifiez sa majuscule : \all{die \textbf{G}esundheit der \textbf{M}enschen und die \textbf{U}mwelt der \textbf{S}tadt}. C'est la faute la plus fréquente — et la plus sanctionnée — dans les productions des élèves francophones.
\end{attention}

\courssub{Exercices d'application}

\begin{exoresolu}[Transformer un article en tract]
Énoncé : À partir de cette phrase d'article — \all{Am Samstag findet um 10 Uhr eine Müllsammelaktion auf dem Marktplatz statt. Die Bürger sollen ihren Müll trennen und Handschuhe mitbringen.} — rédigez trois lignes de tract (une question rhétorique, deux impératifs, une information pratique).
\tcblower
Solution détaillée :
\begin{itemize}
\item Question rhétorique : \all{Willst du in einer sauberen Stadt leben?}
\item Impératifs : \all{Trenne deinen Müll! Bringe Handschuhe mit!} (attention : \all{mitbringen} est séparable $\rightarrow$ \all{bringe ... mit}).
\item Information pratique : \all{Komm am Samstag um 10 Uhr zum Marktplatz!}
\end{itemize}
On passe ainsi du registre factuel (3\up{e} personne, passif/verbe modal : \all{findet ... statt}, \all{sollen ... trennen}) au registre incitatif (2\up{e} personne, impératif), avec majuscules à tous les noms : \all{Müll, Handschuhe, Samstag, Marktplatz}.
\end{exoresolu}

\begin{exercice}[À vous de produire]
\begin{enumerate}
\item Rédigez le chapeau (\all{Einleitung}) d'un article sur une campagne de vaccination dans votre lycée : répondez à \all{Wer? Was? Wo? Wann?} (2 phrases max, passif recommandé).
\item Écrivez trois impératifs pour un tract contre la pollution plastique à Douala, en utilisant les verbes \all{helfen}, \all{schützen} et \all{wegwerfen} (attention à la particule séparable).
\item Relisez cette phrase et corrigez les cinq fautes (majuscules, ordre des mots, Umlaut, ponctuation) : \all{zuerst die schüler sammelten den müll, weil sie die umwelt schützen wollten.}
\end{enumerate}
\end{exercice}

\begin{corrige}[Exercice « À vous de produire »]
\begin{enumerate}
\item Exemple : \all{Seit Montag findet an unserem Gymnasium eine Impfkampagne statt. Organisiert wird sie vom Schulleiter und einem Team von Ärzten.} (Depuis lundi, une campagne de vaccination a lieu dans notre lycée. Elle est organisée par le proviseur et une équipe de médecins.)
\item \all{Hilf mit! Schütze die Umwelt! Wirf kein Plastik weg!} (Participe ! Protège l'environnement ! Ne jette pas de plastique ! — \all{wegwerfen} : impératif \all{wirf ... weg}).
\item Correction : \all{Zuerst sammelten die Schüler den Müll, weil sie die Umwelt schützen wollten.}
      \begin{itemize}
      \item Inversion après \all{zuerst} : \all{Zuerst sammelten...} (verbe 2\up{e}, sujet 3\up{e}).
      \item Majuscules : \all{Schüler}, \all{Müll} (Umlaut \all{ü}), \all{Umwelt}.
      \item Virgule avant \all{weil} (subordonnée).
      \item Verbe \all{wollten} en fin de subordonnée (correct ici).
      \end{itemize}
\end{enumerate}
\end{corrige}

\begin{savaistu}[Le tract, une tradition germanophone]
En Allemagne, en Autriche et en Suisse, le \all{Flugblatt} est un genre très ancien : dès le XVI\up{e} siècle, on distribuait des feuillets imprimés pour informer la population (réforme, révoltes paysannes). Aujourd'hui encore, les associations écologiques comme le \all{BUND} (Bund für Umwelt und Naturschutz Deutschland) ou les groupements locaux de protection de la nature distribuent des tracts sur les marchés — exactement comme les associations de quartier de Yaoundé ou de Douala lors des journées de nettoyage (\all{„Aktion Saubere Stadt“}). Savoir rédiger un tract efficace en allemand est donc une compétence réellement utile, pas seulement scolaire !
\end{savaistu}

\coursec{Traduction : thème et version guidés}

Après avoir maîtrisé la structure de l'article informatif et du tract, nous abordons maintenant l'épreuve reine du baccalauréat : la traduction. Que ce soit de l'allemand vers le français (\emph{Version}) ou du français vers l'allemand (\emph{Thème}), la traduction n'est pas un simple remplacement mot à mot. C'est un processus de \textbf{compréhension}, d'\textbf{analyse structurelle} et de \textbf{reformulation} dans la langue cible. Les thèmes de l'environnement, de la santé et du bien-être offrent un terrain idéal pour réinvestir le passif, les relatives et l'irréel (Konjunktiv II) que nous venons de réviser.

\courssub{Méthode de la version (allemand → français)}

\begin{methode}[Les 5 étapes de la version réussie]
\begin{enumerate}
\item \textbf{Lire globalement} : Lisez le texte une première fois sans stylo pour saisir le sujet, le ton (journalistique, informatif) et le message global.
\item \textbf{Segmenter et analyser} : Découpez en phrases. Pour chaque phrase, identifiez : le \cle{sujet} (qui ? quoi ?), le \cle{verbe conjugué} (temps, voix active/passive, mode), les \cle{compléments} (cas : nominatif, accusatif, datif, génitif) et les \cle{subordonnées} (relatives avec \all{der/die/das}, conjonctives avec \all{dass/weil/wenn}).
\item \textbf{Traduire par unités de sens} : Ne traduisez pas mot à mot. Traitez la proposition principale, puis les subordonnées. Attention : en allemand, le verbe de la subordonnée est à la fin.
\item \textbf{Reformuler en français naturel} : Adaptez la structure. Le passif allemand (\all{werden} + Partizip II) devient souvent un passif français (\all{être} + participe passé) ou \all{on} + actif. Le \all{Konjunktiv II} devient un conditionnel ou un imparfait de l'indicatif selon le contexte. Les relatives allemandes (\all{der/die/das}) deviennent \all{qui/que/dont/où}.
\item \textbf{Relire et polir} : Vérifiez l'orthographe, la grammaire française, la ponctuation et la fluidité. Le français doit sonner « français », pas « allemand traduit ».
\end{enumerate}
\end{methode}

\begin{exemplebox}[Exemple d'analyse de version]
\all{Die Jugendlichen, die die Kampagne organisiert haben, wurden von der Zeitung gelobt.}
\par\medskip
\textbf{Analyse :}
\begin{itemize}
\item Sujet principal : \all{Die Jugendlichen} (Nominatif pluriel) $\rightarrow$ \emph{Les jeunes}.
\item Relative : \all{die die Kampagne organisiert haben} (sujet \all{die} = les jeunes, verbe \all{haben organisiert} au Perfekt à la fin) $\rightarrow$ \emph{qui ont organisé la campagne}.
\item Verbe principal : \all{wurden ... gelobt} (Passif Präterit : \all{werden} + Partizip II) $\rightarrow$ \emph{ont été félicités / ont été loués}.
\item Complément d'agent : \all{von der Zeitung} $\rightarrow$ \emph{par le journal}.
\end{itemize}
\textbf{Traduction :} « Les jeunes qui ont organisé la campagne ont été félicités par le journal. »
\end{exemplebox}

\begin{attention}[Piège n°1 : La traduction mot à mot]
Ne traduisez jamais \all{wurden gelobt} par « sont devenus loués » (calque sur \all{werden} = devenir). Ne traduisez pas \all{die ... haben} par « les quels ... ont ». Identifiez la \textbf{fonction} (sujet, COD, agent) et le \textbf{temps/voix} avant de choisir le mot français.
\end{attention}

\courssub{Méthode du thème (français → allemand)}

\begin{methode}[Les 5 étapes du thème réussi]
\begin{enumerate}
\item \textbf{Analyser la phrase française} : Identifiez le temps (présent, passé composé, imparfait, conditionnel), la voix (actif/passif/« on »), le sujet et les compléments.
\item \textbf{Choisir la structure allemande} :
      \begin{itemize}
      \item \textbf{Passif} : Si le sujet français subit l'action $\rightarrow$ \all{werden} + Partizip II (sujet = ancien COD). Si agent $\rightarrow$ \all{von} + Datif.
      \item \textbf{« On » indéfini} : Souvent \all{man} + actif (conjugaison 3e pers. sg) OU passif \all{werden} (style plus formel).
      \item \textbf{Conditionnel / Hypothèse} : $\rightarrow$ \all{Konjunktiv II} (\all{würde} + infinitif ou forme simple \all{triebe/wäre/hätte}).
      \item \textbf{Relatives} : Antécédent $\rightarrow$ genre/nombre/cas du pronom relatif \all{der/die/das} \textbf{dans la relative}.
      \end{itemize}
\item \textbf{Construire la phrase allemande} : Placez le verbe conjugué (position 2 en principale, position finale en subordonnée). Déclinez articles, adjectifs, pronoms (Cas !).
\item \textbf{Vérifier l'orthographe} : Majuscules aux \cle{Noms}, Umlaute (ä/ö/ü), ß (après voyelle longue/diphtongue), guillemets „ “.
\item \textbf{Relire à voix haute} : Vérifiez l'ordre des mots (TeKaMoLo : Temporel, Kausal, Modal, Local) et la concordance des temps.
\end{enumerate}
\end{methode}

\begin{exemplebox}[Exemples de thème : structures clés]
\begin{center}
\begin{tabularx}{\linewidth}{|X|X|X|}
\hline
\rowcolor{popblueL} \textbf{Français} & \textbf{Allemand} & \textbf{Analyse grammaticale} \\
\hline
Les déchets doivent être triés. & \all{Der Müll muss getrennt werden.} & Passif modal : \all{müssen} + Partizip II \all{getrennt} + \all{werden}. Sujet \all{Der Müll} (masc. sg). \\
\hline
Si chacun faisait du sport, on serait en meilleure santé. & \all{Wenn jeder Sport triebe, wäre man gesünder.} & Irréel présent : \all{Wenn} + Konjunktiv II (\all{triebe}, \all{wäre}). \all{man} = « on ». Verbe en fin de subordonnée (\all{triebe}). \\
\hline
L'eau que nous buvons est précieuse. & \all{Das Wasser, das wir trinken, ist kostbar.} & Relative : antécédent \all{Das Wasser} (neutre) $\rightarrow$ pronom \all{das} (Nominatif dans la relative, sujet de \all{trinken}). Verbe \all{trinken} à la fin. \\
\hline
On a planté des arbres dans le parc. & \all{Man hat Bäume in den Park gepflanzt.} / \all{Im Park wurden Bäume gepflanzt.} & Alternative \all{man} (actif) ou Passif (focus sur \all{Bäume}). \all{in den Park} (Akkusativ direction) vs \all{im Park} (Datif lieu). \\
\hline
Il faut protéger l'environnement. & \all{Man muss die Umwelt schützen.} / \all{Die Umwelt muss geschützt werden.} & \all{Il faut} $\rightarrow$ \all{man muss} + infinitif OU passif modal. \all{Umwelt} (fém.) $\rightarrow$ \all{die Umwelt}. \\
\hline
\end{tabularx}
\end{center}
\end{exemplebox}

\begin{attention}[Piège n°2 : L'ordre des mots en subordonnée (\all{dass, weil, wenn})]
En français, le verbe suit le sujet : « ... \textbf{parce que l'environnement \underline{doit être} protégé} ».
En allemand, le verbe conjugué \textbf{va à la fin} de la proposition :
\all{..., weil die Umwelt geschützt werden \textbf{muss}.} (pas *\all{weil die Umwelt muss geschützt werden})
\all{..., dass wir mehr Sport \textbf{treiben}.}
\all{..., wenn jeder Müll \textbf{trennt}.}
\emph{Règle d'or :} Dès que vous voyez \all{dass, weil, wenn, als, ob}..., attendez la fin de la phrase pour poser le verbe.
\end{attention}

\courssub{Entraînement au thème : 5 phrases mobilisant passif, relatives, Konjunktiv II}

\begin{exoresolu}[Thème guidé : Environnement et santé]
\textbf{Consigne :} Traduisez en allemand. Justifiez chaque choix grammatical (temps, voix, cas, ordre des mots).

\begin{enumerate}
\item \textbf{Français :} « Les arbres qui ont été plantés l'année dernière poussent bien. »
\item \textbf{Français :} « Si le gouvernement interdisait les voitures en ville, l'air serait plus pur. »
\item \textbf{Français :} « On a distribué des prospectus aux élèves pour la campagne de vaccination. »
\item \textbf{Français :} « Le médecin recommande que l'on fasse attention à son alimentation. »
\item \textbf{Français :} « Les jeunes qui participent à l'action nettoient la rivière. »
\end{enumerate}

\tcblower
\textbf{Solutions détaillées :}

\begin{enumerate}
\item \all{Die Bäume, die letztes Jahr gepflanzt wurden, wachsen gut.}
      \par\emph{Analyse :} « Les arbres » = Sujet principal $\rightarrow$ \all{Die Bäume} (Nominatif Pluriel). Relative : antécédent \all{Bäume} $\rightarrow$ pronom \all{die} (Nominatif Pluriel, sujet de la relative). Passif dans la relative : \all{wurden gepflanzt} (Präterit passif, verbe à la fin). Principal : \all{wachsen} (Présent actif), position 2.
\item \all{Wenn die Regierung Autos in der Stadt verbieten würde, wäre die Luft reiner.}
      \par\emph{Analyse :} Hypothèse irréelle présente $\rightarrow$ \all{Konjunktiv II}. Subordonnée \all{Wenn} : verbe \all{würde verbieten} (ou \all{verböte}) à la \textbf{fin}. Principal : \all{wäre} (Konjunktiv II de \all{sein}) position 2. \all{die Regierung} (Nominatif), \all{Autos} (Akkusativ), \all{in der Stadt} (Datif lieu), \all{die Luft} (Nominatif), \all{reiner} (comparatif).
\item \all{Man hat den Schülern Prospekte für die Impfkampagne verteilt.} \quad \textbf{OU} \quad \all{Den Schülern wurden Prospekte für die Impfkampagne verteilt.}
      \par\emph{Analyse :} « On » $\rightarrow$ \all{Man} + Actif (\all{hat verteilt}, Perfekt) OU Passif (\all{wurden verteilt}, Präterit passif). « Aux élèves » = Datif (\all{den Schülern}). « Des prospectus » = Accusatif (\all{Prospekte}) sujet du passif. « Pour la campagne » = \all{für} + Akkusatif (\all{die Impfkampagne}).
\item \all{Der Arzt empfiehlt, dass man auf seine Ernährung achte.} \quad (\text{ou } \all{achten sollte})
      \par\emph{Analyse :} \all{empfiehlt} (Présent) + \all{dass} $\rightarrow$ Subordonnée. « Que l'on fasse » = Subjonctif français $\rightarrow$ \all{Konjunktiv I} (\all{achte}, 3e pers. sg) style rapporté formel, ou \all{Konjunktiv II} (\all{achten sollte}) plus courant. \all{auf etwas achten} : verbe à préposition fixe (\all{auf} + Akkusatif), \textbf{pas} un verbe à particule séparable. La préposition \all{auf} reste devant son objet \all{seine Ernährung}. Le verbe conjugué (\all{achte} ou \all{sollte}) va en position finale. Structure : \all{dass man auf seine Ernährung achte / achten sollte}.
\item \all{Die Jugendlichen, die an der Aktion teilnehmen, säubern den Fluss.}
      \par\emph{Analyse :} « Les jeunes » = \all{Die Jugendlichen} (Nom. Pl.). Relative : antécédent \all{Jugendlichen} $\rightarrow$ pronom \all{die} (Nom. Pl., sujet). Verbe à particule séparable \all{teilnehmen} : en subordonnée, la particule \all{teil} se colle au verbe conjugué final $\rightarrow$ \all{teilnimmt}. Structure relative : \all{die an der Aktion teilnehmen}. Principal : \all{säubern} (Présent, position 2), \all{den Fluss} (Akkusativ).
\end{enumerate}
\end{exoresolu}

\courssub{Entraînement à la version : texte court sur une campagne de santé}

\begin{textebox}[Campagne „Gesund leben“ — Allemagne, 2024]
In Deutschland hat das Bundesministerium für Gesundheit eine neue Kampagne gestartet. „Gesund leben – selbstbestimmt“ heißt das Motto. Plakate, die in allen U-Bahnhöfen und Schulen hängen, zeigen junge Menschen beim Sport und beim gesunden Kochen. Die Botschaft ist klar: Wer sich regelmäßig bewegt und ausgewogen isst, fühlt sich besser und lebt länger.

Die Kampagne wurde von Experten entwickelt, die vor allem Jugendliche erreichen wollen. „Wir müssen früh anfangen“, sagte die Ministerin in einem Interview. „Wenn Kinder lernen, auf ihren Körper zu hören, vermeiden sie später viele Krankheiten.“ Auch soziale Medien spielen eine große Rolle. Influencer, die viele Follower haben, wurden gebeten, die Botschaft zu teilen.

In den ersten Wochen wurde die Webseite der Kampagne über eine Million Mal besucht. Man hofft, dass die Aktion langfristig das Bewusstsein für Gesundheit verändert.
\end{textebox}

\begin{definition}[Glossaire du texte]
\begin{itemize}
\item \all{das Bundesministerium für Gesundheit, die -en} : le Ministère fédéral de la Santé.
\item \all{starten} : lancer, démarrer.
\item \all{das Motto, die -s} : la devise, le slogan.
\item \all{der U-Bahnhof, die U-Bahnhöfe} : la station de métro.
\item \all{ausgewogen} : équilibré (alimentation).
\item \all{sich bewegen} : bouger, faire de l'exercice.
\item \all{der Experte, die -n} : l'expert.
\item \all{erreichen} : atteindre (un public).
\item \all{vermeiden} : éviter.
\item \all{die Krankheit, die -en} : la maladie.
\item \all{der Influencer, die -} : l'influenceur.
\item \all{der Follower, die -} : l'abonné (réseaux sociaux).
\item \all{bitten, bat, gebeten} : prier, demander.
\item \all{das Bewusstsein} : la conscience, la sensibilisation.
\item \all{langfristig} : à long terme.
\item \all{verändern} : changer, modifier.
\end{itemize}
\end{definition}

\begin{exoresolu}[Version guidée : Étapes de traduction]
\textbf{Consigne :} Traduisez le texte en français naturel. Suivez la méthode en 5 étapes.

\tcblower
\textbf{Traduction proposée :}
« En Allemagne, le Ministère fédéral de la Santé a lancé une nouvelle campagne. « Vivre sainement – de manière autonome » est la devise. Des affiches, qui sont placardées dans toutes les stations de métro et les écoles, montrent de jeunes gens faisant du sport et cuisinant sainement. Le message est clair : qui bouge régulièrement et mange équilibré se sent mieux et vit plus longtemps.

La campagne a été développée par des experts qui veulent avant tout atteindre les jeunes. « Nous devons commencer tôt », a déclaré la ministre dans une interview. « Si les enfants apprennent à écouter leur corps, ils évitent plus tard beaucoup de maladies. » Les réseaux sociaux jouent aussi un grand rôle. Des influenceurs, qui ont beaucoup d'abonnés, ont été priés de partager le message.

Dans les premières semaines, le site web de la campagne a été visité plus d'un million de fois. On espère que cette action changera durablement la conscience de la santé. »

\textbf{Justifications grammaticales clés :}
\begin{itemize}
\item \all{hat ... gestartet} (Perfekt actif) $\rightarrow$ « a lancé » (Passé composé).
\item \all{die ... hängen} (Relative, présent) $\rightarrow$ « qui sont placardées / qui pendent » $\rightarrow$ choix du passif/état en français pour \all{hängen} (accrochées). \all{in allen U-Bahnhöfen} (Datif pluriel lieu).
\item \all{beim Sport / beim gesunden Kochen} (Préposition \all{bei} + Datif + nominalisation) $\rightarrow$ « en faisant du sport / en cuisinant sainement » (gérondif/participe présent).
\item \all{Wer ... isst, fühlt ... lebt} (Relative libre \all{Wer} = sujet générique) $\rightarrow$ « Qui / Celui qui ... ».
\item \all{wurde ... entwickelt} (Passif Präterit) $\rightarrow$ « a été développée ».
\item \all{die ... wollen} (Relative, sujet \all{die} = experts) $\rightarrow$ « qui veulent ».
\item \all{Wenn Kinder lernen ... vermeiden} (Subordonnée \all{Wenn} + Principal) $\rightarrow$ Verbe \all{lernen} à la fin de la subordonnée, \all{vermeiden} position 2 du principal.
\item \all{auf ihren Körper zu hören} (Infinitif avec \all{zu}) $\rightarrow$ « à écouter leur corps ».
\item \all{wurden gebeten} (Passif Präterit) $\rightarrow$ « ont été priés / invités ».
\item \all{wurde ... besucht} (Passif Präterit impersonnel / sujet \all{die Webseite}) $\rightarrow$ « a été visitée ».
\item \all{dass ... verändert} (Subordonnée \all{dass}, verbe \all{verändert} à la fin) $\rightarrow$ « que ... change / changera ».
\end{itemize}
\end{exoresolu}

\courssub{Pièges de traduction et correction collective}

\begin{attention}[Les 3 pièges majeurs pour les francophones]
\begin{enumerate}
\item \textbf{« On » $\neq$ \all{man} systématique.}
      \begin{itemize}
      \item Sujet indéfini, actif : \all{Man sagt...} (« On dit... »).
      \item Passif français (« Les arbres ont été plantés ») $\rightarrow$ Passif allemand \all{Die Bäume wurden gepflanzt} (meilleur style écrit) OU \all{Man hat Bäume gepflanzt} (oral).
      \item « On » = « Nous » (familier) $\rightarrow$ \all{Wir}.
      \end{itemize}
\item \textbf{Conditionnel français $\rightarrow$ Konjunktiv II allemand.}
      \begin{itemize}
      \item Hypothèse (\all{Si} + imparfait / \all{Wenn} + Konj II) : \all{Wenn ich Zeit hätte, käme ich.}
      \item Politesse / Conseil : \all{Ich würde Ihnen raten...} (\all{würde} + infinitif).
      \item Discours rapporté (Konjunktiv I) : \all{Er sagt, er \textbf{komme} morgen.} (Pas Konj II ici !).
      \end{itemize}
\item \textbf{Ordre des mots (SOV) en subordonnée.}
      \all{..., weil / dass / wenn / als / ob} $\rightarrow$ \textbf{Verbe conjugué tout à la fin.}
      \all{..., weil die Umwelt \textbf{geschützt werden muss}.} (Pas *\all{muss geschützt werden}).
      Particule séparable : \all{..., weil er \textbf{aufräumt}.} (\all{aufräumen} $\rightarrow$ \all{räumt auf} $\rightarrow$ \all{aufräumt}).
\end{enumerate}
\end{attention}

\begin{experience}[Correction collective : Comparer pour progresser]
\textbf{Activité (20 min) :}
\begin{enumerate}
\item L'enseignant projette une phrase thème difficile (ex. phrase 2 de l'exercice résolu : « Si le gouvernement interdisait les voitures en ville, l'air serait plus pur. »).
\item Chaque élève écrit sa proposition au tableau / sur papier (anonyme).
\item On sélectionne 3 propositions variées (ex. \all{Wenn die Regierung Autos verbietet...} / \all{Falls die Regierung Autos verbieten würde...} / \all{Würde die Regierung Autos verbieten...}).
\item \textbf{Débat guidé :} « Laquelle est correcte ? Pourquoi les autres le sont-elles moins ? » (Identification du \all{Konjunktiv II} obligatoire pour l'irréel, place du verbe en \all{Wenn}-Satz, accord \all{Luft} fém. $\rightarrow$ \all{wäre}).
\item On généralise : \textbf{Tableau de choix} (voir ci-dessous).
\end{enumerate}
\end{experience}

\begin{center}
\begin{tabularx}{\linewidth}{|X|X|X|}
\hline
\rowcolor{popblueL} \textbf{Proposition élève} & \textbf{Correct ?} & \textbf{Justification grammaticale} \\
\hline
\all{Wenn die Regierung Autos in der Stadt verbietet, ist die Luft reiner.} & \textbf{Non} & Indicatif présent (\all{verbietet, ist}) = réel. Le français « interdirait / serait » marque l'irréel $\rightarrow$ \all{Konjunktiv II} (\all{verbiete/würde verbieten}, \all{wäre}). \\
\hline
\all{Wenn die Regierung Autos in der Stadt verbieten würde, wäre die Luft reiner.} & \textbf{Oui} & \all{Wenn} + \all{Konjunktiv II} (\all{würde verbieten} à la fin). Principal : \all{wäre} (Konj II de \all{sein}) position 2. Comparatif \all{reiner}. Structure parfaite. \\
\hline
\all{Falls die Regierung Autos verböte, würde die Luft reiner sein.} & \textbf{Oui (formel)} & \all{Falls} = \all{Wenn}. \all{verböte} = Konj II forme simple (rare, soutenu). Principal : \all{würde ... sein} (Konj II perifrastique). Correct mais style très administratif. \\
\hline
\all{Wenn die Regierung die Autos verbieten würde, die Luft wäre reiner.} & \textbf{Oui} & Principal : \all{die Luft} (Pos 1), \all{wäre} (Pos 2, V2) — ordre correct. Subordonnée : \all{verbieten würde} (fin). Nuance : \all{die Autos} (spécifiques) vs \all{Autos} (général), les deux acceptables. \\
\hline
\end{tabularx}
\end{center}

\courssub{Organigramme : Les étapes de la traduction}

\begin{popfigure}
\begin{tikzpicture}[
    node distance=1.2cm and 2cm,
    every node/.style={font=\small, align=center},
    stepS/.style={rectangle, rounded corners, draw=popblue, thick, fill=popblueL, minimum width=2.8cm, minimum height=1.2cm},
    arrow/.style={-Stealth, thick, popdark},
    decision/.style={diamond, aspect=2, draw=poporange, thick, fill=poporangeL, minimum width=2.5cm, minimum height=1.2cm}
]
% Nodes
\node[stepS, align=center] (start){Lire le texte\\globalement};
\node[stepS, below=of start, align=center] (analyze){Analyser les structures\\(Verbe, Cas, Subordonnées,\\Passif, Konjunktiv)};
\node[stepS, below=of analyze, align=center] (translate){Traduire par unités\\de sens (Phrase par phrase)};
\node[decision, below=of translate, align=center] (check){Français naturel ?\\Allemand correct ?\\(Ordre, Cas, Temps)};
\node[stepS, below=of check, align=center] (polish){Relire \& Polir\\(Orthographe, Majuscules,\\Ponctuation, Fluidité)};

% Arrows
\draw[arrow] (start) -- (analyze);
\draw[arrow] (analyze) -- (translate);
\draw[arrow] (translate) -- (check);
\draw[arrow] (check) -- node[left, xshift=-0.3cm] {Non} ++(-3,0) |- (analyze);
\draw[arrow] (check) -- node[right, xshift=0.3cm] {Oui} (polish);

% Loop back annotation
\node[left=0.5cm of check, text=poporange, font=\tiny\itshape, align=center]{Retour analyse\\si erreur};
\end{tikzpicture}
\legende{Processus itératif de la traduction (Version et Thème). La flèche de retour symbolise l'autocorrection.}
\end{popfigure}

\courssub{Expression écrite : Mini-production pour le Bac}

\begin{exercice}[Sujet d'entraînement type Bac]
\textbf{Consigne :} Vous participez à un échange scolaire à Berlin. Écrivez un article court (environ 120 mots) pour le journal du lycée partenaire sur le thème : \all{„Gesundheit und Umwelt an unserer Schule“}. Utilisez \textbf{au moins} : 1 passif, 1 relative, 1 \all{Konjunktiv II} (hypothèse ou conseil), 1 subordonnée \all{weil/dass/wenn}.
\end{exercice}

\begin{corrige}[Exercice : Article „Gesundheit und Umwelt“]
\textbf{Proposition de correction :}
\all{An unserer Schule wird die Umwelt großgeschrieben. Letztes Jahr wurde ein Schulgarten angelegt, der den Schülern frisches Gemüse liefert. Die AG „Gesunde Pause“, die von Schülern der Oberstufe geleitet wird, verkauft jeden Morgen Vollkornbrötchen und Obst. Wenn jeder Schüler sein Pausenbrot selbst mitbrächte, gäbe es weniger Verpackungsmüll. Unser Direktor sagt, dass die Schule nächstes Jahr „Klimaschule“ werden soll. Wir hoffen, dass alle mitmachen, weil Gesundheit und Umweltschutz zusammengehören.}

\textbf{Vérification des contraintes :}
\begin{itemize}
\item Passif : \all{wird ... großgeschrieben}, \all{wurde ... angelegt}, \all{wird ... geleitet}, \all{soll ... werden}.
\item Relative : \all{der den Schülern ... liefert} (antécédent \all{Schulgarten}, masc, Nom), \all{die von Schülern ... geleitet wird} (antécédent \all{AG}, fém, Nom, passif dans relative).
\item Konjunktiv II : \all{mitbrächte} (forme simple), \all{gäbe} (forme simple) — Hypothèse irréelle présente.
\item Subordonnée \all{dass} : \all{dass die Schule ... werden soll} (verbe \all{soll} à la fin).
\item Subordonnée \all{weil} : \all{weil Gesundheit und Umweltschutz zusammengehören} (verbe \all{zusammengehören} à la fin).
\item Vocabulaire thématique : \all{Schulgarten, Vollkornbrötchen, Verpackungsmüll, Klimaschule, Umweltschutz}.
\end{itemize}
\end{corrige}

\begin{savaistu}[Le saviez-vous ? La notation de la version au Baccalauréat]
Au Bac camerounais (Allemand LV2), l'épreuve de version est notée sur \textbf{deux critères indissociables} :
\begin{enumerate}
\item \textbf{La fidélité au sens (Respect du contenu) :} Avez-vous rendu l'information exacte ? Les chiffres, les noms propres, la négation, le mode (hypothèse vs réel) sont-ils conservés ?
\item \textbf{La qualité de la langue française (Expression) :} Le français est-il correct (orthographe, grammaire, syntaxe) et \textbf{idiomatique} ? Une phrase comme „Les jeunes qui ont organisé la campagne sont devenus loués par le journal“ (calque de \all{werden}) comprend le sens mais rate l'expression $\rightarrow$ \textbf{pénalité}. Une phrase comme „Les jeunes organisateurs de la campagne ont reçu les félicitations du journal“ (reformulation élégante) $\rightarrow$ \textbf{bonus}.
\end{enumerate}
\cle{Conseil de pro :} Ne sacrifiez jamais le français naturel pour coller à la structure allemande. Si vous devez changer la voix (passif $\rightarrow$ actif), le mode (subjonctif $\rightarrow$ indicatif/conditionnel) ou la structure de phrase pour que ce soit du bon français, \textbf{faites-le}. C'est la marque du traducteur compétent.
\end{savaistu}

\newpage

\coursec{Activité d'intégration}

\coursec{Leçon AL15 : Traduction et production de textes informatifs sur l'environnement, la santé et le bien-être}

\courssub{Objectifs de l'activité}

À l'issue de cette activité d'intégration, tu seras capable de :
\begin{itemize}
\item rédiger en allemand un \cle{tract} informatif et persuasif (der Flugzettel / das Flugblatt) pour une campagne de sensibilisation sur la santé et l'environnement ;
\item réinvestir dans une production authentique les structures grammaticales du chapitre : le \cle{passif} (Präsens et Futur), les \cle{propositions relatives}, le \cle{Konjunktiv II} de l'irréel et du conseil, ainsi que les \cle{impératifs} ;
\item mobiliser le lexique thématique de l'environnement, de la santé et du bien-être ;
\item traduire un texte court de l'allemand vers le français (\cle{version}) en respectant le sens, le registre et les contraintes des deux langues ;
\item comparer, évaluer et corriger une traduction en la confrontant à l'original (relecture : orthographe, majuscules, ordre des mots, cas).
\end{itemize}

\courssub{Situation}

Le club environnement de ton lycée à Yaoundé, en partenariat avec le club d'allemand, organise la semaine verte du lycée. À l'occasion de la journée portes ouvertes, à laquelle participeront des élèves germanophones en échange scolaire avec un Gymnasium de Hambourg, le club veut distribuer un \cle{tract bilingue} (allemand-français) intitulé \all{„Gesund und umweltfreundlich leben“} (« Vivre sainement et respecter l'environnement »). Le proviseur a demandé que le tract soit rédigé d'abord en allemand par les élèves de Terminale A, puis traduit en français pour les familles. Le meilleur tract sera imprimé et affiché à l'entrée du lycée, près du marché de Mokolo où la campagne de sensibilisation aura lieu.

Voici la note de service du club, qui fixe le cadre de la tâche :

\begin{textebox}[Note du club environnement — Aufgabenstellung]
Liebe Mitglieder,\par
für unsere Kampagne „Gesund und umweltfreundlich leben“ brauchen wir ein Flugblatt auf Deutsch. Es soll einen kurzen Slogan, vier bis fünf klare Sätze mit Imperativen und praktische Informationen (Ort, Datum, Uhrzeit) enthalten. Mindestens ein Satz muss im Passiv stehen, einer muss einen Relativsatz enthalten, und ein Ratschlag soll im Konjunktiv II formuliert werden. Danach wird das Flugblatt ins Französische übersetzt, damit alle Eltern es verstehen können.\par
Viel Erfolg!\par
Euer Umweltclub
\end{textebox}

\begin{definition}[Glossaire de la note]
der Ratschlag, die Ratschläge : le conseil ; das Flugblatt, die Flugblätter : le tract ; der Relativsatz, die Relativsätze : la proposition relative ; übersetzen : traduire ; damit : afin que ; die Eltern (pl.) : les parents.
\end{definition}

\courssub{Consignes}

Travail en groupes de 3 ou 4 élèves. Suis les étapes dans l'ordre.

\begin{enumerate}
\item \textbf{Compréhension et repérage.} Lis la note du club ci-dessus. Souligne les quatre contraintes grammaticales exigées (impératifs, passif, relative, Konjunktiv II) et les informations pratiques à mentionner. Reformule oralement la tâche en français en deux phrases.
\item \textbf{Remue-méninges lexical.} En 5 minutes, constitue avec ton groupe une liste de 12 mots ou expressions allemands utiles pour le thème (par exemple : \all{die Umwelt schützen}, \all{Müll trennen}, \all{gesund essen}, \all{Sport treiben}, \all{Wasser sparen}, \all{Plastik vermeiden}, \all{der Wald}, \all{die Gesundheit}, \all{recyceln}, \all{zu Fuß gehen}, \all{Obst und Gemüse}, \all{die Luft}). Classe-les en deux colonnes : \emph{Umwelt} / \emph{Gesundheit}.
\item \textbf{Choix du slogan.} Rédigez un slogan court et rythmé (maximum 8 mots), par exemple avec un impératif ou une comparaison : \all{„Schütze die Natur — sie schützt dich!“}
\item \textbf{Rédaction du tract.} Écrivez le corps du tract : 4 à 5 phrases contenant obligatoirement :
\begin{itemize}
\item deux \cle{impératifs} (2\textsuperscript{e} personne du pluriel : \all{Trennt euren Müll!}) ;
\item une phrase au \cle{passif} (Präsens ou Futur : \all{Es werden Bäume gepflanzt.}) ;
\item une \cle{proposition relative} correctement accordée (\all{...das Wasser, das wir trinken,...}) ;
\item un conseil au \cle{Konjunktiv II} (\all{Du solltest mehr Sport treiben.} / \all{Wir könnten...}).
\end{itemize}
Ajoutez les informations pratiques : lieu, date, heure de la journée de sensibilisation.
\item \textbf{Relecture grammaticale.} Échangez le tract avec le groupe voisin pour une première correction : vérifiez la majuscule à tous les noms, la place du verbe conjugué en 2\textsuperscript{e} position, la place du participe II et de \all{werden} au passif, l'accord du pronom relatif en genre, nombre et cas.
\item \textbf{Traduction (version).} Chaque groupe traduit en français le tract allemand de l'autre groupe. La traduction doit être fidèle mais naturelle : un impératif allemand se rend par un impératif français, le passif peut être rendu par « on » si c'est plus idiomatique.
\item \textbf{Comparaison et validation.} Les deux groupes comparent la traduction française avec l'original allemand : relevez deux différences de formulation et discutez-en. Corrigez ensemble les éventuelles erreurs.
\item \textbf{Affichage et vote.} Les tracts bilingues sont affichés au mur. La classe vote le meilleur tract selon trois critères : correction grammaticale, force du slogan, clarté des informations pratiques.
\end{enumerate}

\courssub{Ressources à mobiliser}

\begin{propriete}[Rappel : le passif au Präsens et au Futur]
\cle{Passif Präsens} : \all{werden} (conjugué) + Partizip II en fin de phrase. \par
\all{Der Müll wird getrennt.} « Les déchets sont triés. »\par
\cle{Passif Futur} : \all{werden} (conjugué) + Partizip II + \all{werden} à la fin. \par
\all{Es werden Bäume gepflanzt werden.} « Des arbres seront plantés. »\par
L'agent est introduit par \all{von} + datif : \all{von den Schülern} « par les élèves ».
\end{propriete}

\begin{propriete}[Rappel : relatives et Konjunktiv II]
Le \cle{pronom relatif} s'accorde en genre et en nombre avec son antécédent, et prend le cas exigé par sa fonction dans la relative ; le verbe conjugué va en fin de proposition. \par
\all{Die Bäume, die wir pflanzen, schützen die Luft.} « Les arbres que nous plantons protègent l'air. »\par
\cle{Konjunktiv II} du conseil : \all{solltest}, \all{könntest}, \all{würdest} + infinitif. \par
\all{Du solltest weniger Plastik benutzen.} « Tu devrais utiliser moins de plastique. »
\end{propriete}

\begin{propriete}[Rappel : l'impératif à la 2\textsuperscript{e} personne du pluriel]
Identique à la forme \all{ihr} du présent, sans pronom, verbe en première position : \par
\all{Esst mehr Obst! Trinkt viel Wasser! Spart Energie!} « Mangez plus de fruits ! Buvez beaucoup d'eau ! Économisez l'énergie ! »
\end{propriete}

\begin{center}
\begin{tabularx}{\linewidth}{|X|X|X|}
\hline
\rowcolor{popblueL} Français & Deutsch & Remarque \\
\hline
l'environnement & die Umwelt (sing.) & toujours avec article \\
\hline
protéger & schützen & régulier \\
\hline
trier les déchets & den Müll trennen & verbe à particule séparable \\
\hline
la santé & die Gesundheit (sing.) & féminin \\
\hline
faire du sport & Sport treiben & fort : treiben, trieb, getrieben \\
\hline
économiser l'eau & Wasser sparen & — \\
\hline
éviter le plastique & Plastik vermeiden & fort : vermeiden, vermied, vermieden \\
\hline
planter des arbres & Bäume pflanzen & régulier \\
\hline
le bien-être & das Wohlbefinden (sing.) & neutre \\
\hline
\end{tabularx}
\end{center}

\begin{attention}[Pièges fréquents]
\begin{itemize}
\item Ne dis pas \all{*die Bäume werden gepflanzt von den Schülern werden} : au passif Präsens, un seul \all{werden}, conjugué ; le participe II va en fin de phrase.
\item Le pronom relatif n'est pas toujours \all{das} : vérifie le genre de l'antécédent (\all{die Umwelt, die...} ; \all{der Wald, der...}).
\item Attention au faux ami : \all{die Rente} n'est pas « la rente » mais « la retraite » ; pour « campagne » au sens d'action de sensibilisation, dis \all{die Kampagne}, pas \all{*die Landschaft}.
\item À l'impératif pluriel, pas de pronom : \all{Trennt den Müll!} et non \all{*Ihr trennt den Müll!}.
\end{itemize}
\end{attention}

\courssub{Proposition de production et corrigé commenté}

\begin{exoresolu}[Modèle de tract bilingue — production attendue]
Énoncé : rédige le tract demandé par le club (slogan, 4-5 phrases, informations pratiques), puis propose sa traduction française.\par
\tcblower
\textbf{Production modèle en allemand :}\par
\medskip
\all{„Gesund leben — die Umwelt schützen!“}\par
\medskip
\all{Liebe Schülerinnen und Schüler!}\par
\all{Unsere Gesundheit und unsere Umwelt, die wir schützen müssen, sind unser größter Reichtum. Trennt euren Müll und vermeidet Plastik! Esst mehr Obst und Gemüse und treibt regelmäßig Sport! Am Samstag werden im Schulhof Bäume gepflanzt. Du solltest auch zu Fuß oder mit dem Fahrrad zur Schule kommen — das spart Energie und schützt die Luft.}\par
\medskip
\all{Wann? Samstag, 15. Juni, ab 9 Uhr. Wo? Im Schulhof unserer Schule. Kommt alle!}\par
\medskip
\textbf{Traduction française (version) :}\par
« Vivre sainement — protéger l'environnement ! »\par
Chers élèves, notre santé et notre environnement, que nous devons protéger, sont notre plus grande richesse. Triez vos déchets et évitez le plastique ! Mangez plus de fruits et de légumes et faites régulièrement du sport ! Samedi, des arbres seront plantés dans la cour de l'école. Tu devrais aussi venir à l'école à pied ou à vélo — cela économise l'énergie et protège l'air.\par
Quand ? Samedi 15 juin, à partir de 9 heures. Où ? Dans la cour de notre école. Venez nombreux !\par
\medskip
\textbf{Commentaire des choix de langue :}
\begin{itemize}
\item \emph{Relative} : \all{die wir schützen müssen} — le pronom \all{die} reprend l'antécédent \all{unsere Gesundheit und unsere Umwelt} (pluriel, accusatif dans la relative) ; avec un verbe modal, le modal conjugué (\all{müssen}) occupe la dernière place, l'infinitif (\all{schützen}) le précède immédiatement.
\item \emph{Impératifs pluriels} : \all{Trennt}, \all{vermeidet}, \all{Esst}, \all{treibt} — forme \all{ihr} sans pronom, verbe en tête de phrase ; note l'impératif fort \all{Esst} (de \all{essen}, irrégulier : du isst, ihr esst).
\item \emph{Passif Präsens} : \all{werden ... Bäume gepflanzt} — \all{werden} conjugué en 2\textsuperscript{e} position (sujet \all{Bäume} au pluriel), participe II \all{gepflanzt} en fin de phrase ; la traduction française « des arbres seront plantés » (passif) ou « on plantera des arbres » (indéfini) sont toutes deux correctes.
\item \emph{Konjunktiv II} : \all{Du solltest ... kommen} — conseil atténué, rendu par « tu devrais » ; l'infinitif \all{kommen} est renvoyé en fin de phrase.
\item \emph{Lexique} : \all{der Reichtum} (la richesse), \all{der Schulhof} (la cour de récréation), \all{regelmäßig} (régulièrement), \all{zu Fuß} (à pied), \all{das Fahrrad} (le vélo).
\item \emph{Version} : le passif allemand est rendu par un vrai passif français (« seront plantés ») ; le slogan infinitif allemand se rend naturellement par deux infinitifs français.
\end{itemize}
\end{exoresolu}

\courssub{Bilan de l'activité}

Cette activité t'a permis de mobiliser simultanément l'ensemble des savoir-faire du chapitre : comprendre une consigne authentique en allemand, sélectionner et organiser le lexique de l'environnement et de la santé, rédiger un texte informatif et persuasif de type tract, y intégrer volontairement des structures grammaticales complexes (passif, relative, Konjunktiv II, impératifs), puis passer d'une langue à l'autre par la traduction. La phase de comparaison entre l'original et la version t'a entraîné à la relecture critique : vérifier la fidélité du sens, mais aussi l'orthographe allemande, la majuscule des noms, la place du verbe et l'accord des pronoms relatifs. Ces compétences de production et de médiation sont exactement celles exigées à l'épreuve d'allemand du baccalauréat A : savoir écrire un texte court, correct et adapté à une situation de communication réelle, et savoir traduire en respectant les deux langues. Garde en tête la démarche en quatre temps : \cle{lire} la consigne, \cle{planifier} (lexique + structures obligatoires), \cle{rédiger}, \cle{relire} — elle est transférable à toute production écrite d'examen.

\newpage

\coursec{Méthodes et exercices résolus}

\courssub{Méthode 1 : Rédiger un article informatif sur l'environnement, la santé ou le bien-être}

\begin{methode}[Rédiger un article informatif]
\begin{enumerate}
\item \textbf{Analyser le sujet} : identifier le thème (environnement, santé, bien-être), le destinataire (lecteurs d'un journal scolaire, d'un magazine) et la longueur demandée.
\item \textbf{Construire le plan en trois parties} :
\begin{itemize}
\item \cle{Titre} accrocheur + \cle{chapeau} (1--2 phrases qui résument l'essentiel) ;
\item \cle{corps de l'article} : faits, causes, conséquences, exemples concrets (Cameroun, Allemagne), éventuellement une citation d'un acteur ;
\item \cle{conclusion} : appel à l'action ou ouverture.
\end{itemize}
\item \textbf{Choisir le registre} : ton objectif, verbes d'information (\all{berichten, erklären, warnen, empfehlen}), passif pour les faits (\all{Es wird berichtet, dass...}).
\item \textbf{Rédiger} en réinvestissant le lexique du chapitre et les structures révisées : passif, propositions relatives, Konjunktiv II pour les hypothèses.
\item \textbf{Relire} : majuscules aux noms, verbe en 2\textsuperscript{e} position, participes en fin de phrase, accords des cas.
\end{enumerate}
\textbf{Structures à retenir :} \all{Nach Angaben des Gesundheitsministeriums...} (selon le ministère de la Santé) ; \all{Immer mehr Menschen leiden unter...} (de plus en plus de personnes souffrent de) ; \all{Experten warnen vor...} (les experts mettent en garde contre) ; \all{Es ist wichtig, dass...} (il est important que).
\end{methode}

\begin{exoresolu}[Rédaction d'un article informatif]
\textbf{Énoncé (type Bac)} : Vous êtes rédacteur du journal de votre lycée à Yaoundé. Rédigez un article informatif (12--15 lignes) sur le thème : \emph{« La pollution plastique dans notre ville : un danger pour la santé »}. Vous présenterez les faits, les risques et une solution.
\tcblower
\textbf{Raisonnement} :
\begin{enumerate}
\item Sujet : environnement + santé ; destinataire : lycéens ; registre informatif.
\item Plan : titre + chapeau ; faits (sachets plastiques, caniveaux bouchés) ; risques sanitaires (paludisme, maladies) ; solution (recyclage, sensibilisation).
\item Grammaire à réinvestir : passif (\all{werden weggeworfen}), relatives (\all{die Plastiktüten, die...}), Konjunktiv II (\all{Man sollte...}).
\end{enumerate}
\textbf{Article rédigé} :

\all{\textbf{Plastikmüll in Yaoundé -- eine Gefahr für unsere Gesundheit}}

\all{In den Straßen von Yaoundé sieht man überall Plastiktüten und Flaschen, die achtlos weggeworfen werden. Dieser Müll, der die Abflüsse verstopft, ist nicht nur hässlich, sondern auch gefährlich.}

\all{Nach Angaben von Ärzten sammelt sich in den Plastikabfällen Regenwasser, in dem sich Mücken vermehren. So wird die Verbreitung von Malaria begünstigt. Außerdem werden viele Tiere krank, weil sie Plastik fressen.}

\all{Was könnte man tun? Zuerst sollten die Schüler über die Folgen der Verschmutzung informiert werden. Dann müssten mehr Mülleimer aufgestellt und der Müll regelmäßig recycelt werden. Wenn jeder mitmachte, wäre unsere Stadt sauberer und gesünder.}

\all{\emph{Schützen wir gemeinsam unsere Umwelt -- sie ist unsere Zukunft!}}

\textbf{Justifications grammaticales} : \all{weggeworfen werden} = passif présent, participe en fin de proposition ; \all{die achtlos weggeworfen werden} = relative avec verbe en fin ; \all{in dem sich Mücken vermehren} = relative au datif après la préposition \all{in} ; \all{Man sollte / müsste / mitmachte, wäre} = Konjunktiv II pour exprimer des mesures hypothétiques ; verbe conjugué toujours en 2\textsuperscript{e} position dans les principales (\all{So wird die Verbreitung begünstigt}).
\textbf{Conclusion} : l'article respecte la structure demandée (titre, chapeau, faits, risques, solution, appel) et réinvestit les trois points de grammaire du chapitre.
\end{exoresolu}

\courssub{Méthode 2 : Traduire un texte du français vers l'allemand (thème)}

\begin{methode}[Réussir le thème]
\begin{enumerate}
\item \textbf{Lire tout le texte} avant de traduire : repérer le temps dominant, le registre, les liens logiques.
\item \textbf{Traduire phrase par phrase}, en commençant par placer le \cle{sujet} et le \cle{verbe conjugué} (2\textsuperscript{e} position), puis les compléments, puis les éléments verbaux non conjugués \cle{en fin de phrase}.
\item \textbf{Choisir les temps correctement} : passé composé français de narration $\rightarrow$ \cle{Perfekt} (ou Präteritum pour \all{sein, haben, modalverben}) ; conditionnel $\rightarrow$ \cle{Konjunktiv II} ; « est + participe » $\rightarrow$ \cle{Passiv}.
\item \textbf{Attention aux cas} : vérifier la fonction de chaque nom (sujet = nominatif, COD = accusatif, COI = datif) et l'article qui l'accompagne.
\item \textbf{Relire} : majuscules aux noms, particules séparables en fin, orthographe (\all{ß}, Umlaute), aucun mot français laissé.
\end{enumerate}
\textbf{Réflexe anti-calque} : ne jamais traduire mot à mot. « Il y a de plus en plus de déchets » ne se traduit pas par \all{*Es gibt mehr und mehr Müll} mais par \all{Es gibt immer mehr Müll}.
\end{methode}

\begin{exoresolu}[Thème guidé : la santé au quotidien]
\textbf{Énoncé (type Bac)} : Traduisez en allemand :

« De nos jours, beaucoup de jeunes Camerounais ne font pas assez de sport. La plupart de leur temps libre est passée devant l'écran du téléphone portable. Pourtant, les médecins recommandent au moins trente minutes d'activité physique par jour. Si les jeunes bougeaient davantage, ils seraient en meilleure santé et apprendraient mieux à l'école. »
\tcblower
\textbf{Raisonnement phrase par phrase} :

\emph{Phrase 1} : « De nos jours » = \all{Heutzutage} ; « faire du sport » = \all{Sport treiben/machen} ; négation avec \all{nicht} ; « assez » = \all{genug}. Temps : présent.

\emph{Phrase 2} : « est passée » = passif au présent $\rightarrow$ \all{wird verbracht} (le verbe français « passer du temps » se dit \all{verbringen} en allemand ; au passif : \all{wird ... verbracht}). Attention : « la plupart de » = \all{die meiste Zeit} (féminin singulier, donc \all{wird}, pas \all{werden}).

\emph{Phrase 3} : « recommander » = \all{empfehlen} (verbe fort : \all{empfiehlt}) ; « au moins » = \all{mindestens} ; « par jour » = \all{pro Tag} (et non \all{*per Tag} dans ce registre courant).

\emph{Phrase 4} : conditionnel $\rightarrow$ Konjunktiv II : \all{bewegen} $\rightarrow$ \all{bewegten} (ou \all{würden sich bewegen}) ; \all{sein} $\rightarrow$ \all{wären} ; « davantage » = \all{mehr} ; « en meilleure santé » = \all{gesünder} (comparatif de \all{gesund}).

\textbf{Traduction rédigée} :

\all{Heutzutage treiben viele junge Kameruner nicht genug Sport. Die meiste Freizeit wird vor dem Bildschirm des Handys verbracht. Dabei empfehlen die Ärzte mindestens dreißig Minuten körperliche Aktivität pro Tag. Wenn sich die Jugendlichen mehr bewegen würden, wären sie gesünder und würden in der Schule besser lernen.}

\textbf{Justifications} : \all{treiben} en 2\textsuperscript{e} position après \all{Heutzutage} ; \all{wird ... verbracht} = passif présent, participe final ; \all{dem Bildschirm} = datif après \all{vor} (position) ; \all{Wenn ... würden, wären ... würden} = irréel du présent correctement formé avec \all{sich bewegen} (verbe réfléchi : ne pas oublier \all{sich}).
\textbf{Conclusion} : quatre phrases, quatre pièges évités : passif, datif, Konjunktiv II, verbe réfléchi.
\end{exoresolu}

\courssub{Méthode 3 : Traduire un texte de l'allemand vers le français (version)}

\begin{methode}[Réussir la version]
\begin{enumerate}
\item \textbf{Lire le texte entier} deux fois : dégager le sens général avant de traduire le premier mot.
\item \textbf{Découper chaque phrase allemande} : repérer le verbe conjugué, puis chercher en fin de phrase le participe, l'infinitif ou la particule séparable qui le complète.
\item \textbf{Identifier les structures} : passif (\all{wird/werden + Partizip II} = « est/sont + participe »), relatives (traduire la relative après son antécédent), subordonnées (verbe final).
\item \textbf{Traduire en bon français} : un mot allemand n'a pas toujours un équivalent unique ; adapter (\all{die Volksgesundheit} = « la santé publique », \all{sich erholen} = « se reposer, récupérer »).
\item \textbf{Relire le français} : il doit être naturel, correct, sans germanisme (pas de « *il est devenu dit » pour \all{es wurde gesagt}, mais « on a dit / il a été dit »).
\end{enumerate}
\textbf{Réflexe} : pour le passif allemand, le français préfère souvent « on » : \all{In der Schule wird viel über Umweltschutz gesprochen} = « On parle beaucoup de protection de l'environnement à l'école ».
\end{methode}

\begin{exoresolu}[Version guidée : le bien-être au travail]
\textbf{Énoncé (type Bac)} : Traduisez en français :

\all{In vielen deutschen Firmen wird heute mehr auf das Wohlbefinden der Angestellten geachtet. Kurze Pausen, die regelmäßig gemacht werden, verbessern die Konzentration und verringern den Stress. Ein Firmenarzt erklärt: „Wer sich gesund ernährt und genug schläft, ist leistungsfähiger.“ Wenn mehr Betriebe solche Maßnahmen einführen würden, gäbe es weniger kranke Arbeitnehmer.}
\tcblower
\textbf{Raisonnement} :

\emph{Phrase 1} : \all{wird ... geachtet} = passif avec \all{achten auf + Akkusativ} (« faire attention à ») ; \all{das Wohlbefinden} = « le bien-être » ; \all{die Angestellten} = « les employés/salariés ». Le français naturel préfère « on ».

\emph{Phrase 2} : relative \all{die regelmäßig gemacht werden} = passif dans une relative : « qui sont faites régulièrement » ; \all{verringern} = « réduire » ; \all{die Konzentration} et \all{den Stress} : attention, deux cas différents (nominatif / accusatif) mais en français deux COD.

\emph{Phrase 3} : \all{Firmenarzt} = « médecin du travail / médecin d'entreprise » ; \all{Wer ... , ist ...} = « Celui qui... » ; \all{sich gesund ernähren} = « se nourrir sainement / manger sainement » ; \all{leistungsfähiger} = « plus performant, plus efficace » (comparatif).

\emph{Phrase 4} : \all{würden einführen} et \all{gäbe es} = Konjunktiv II $\rightarrow$ conditionnel français ; \all{die Maßnahmen} = « les mesures » ; \all{der Betrieb} = « l'entreprise » (faux ami : pas « l'opération »).

\textbf{Traduction rédigée} :

« Aujourd'hui, dans de nombreuses entreprises allemandes, on accorde plus d'attention au bien-être des salariés. De courtes pauses, qui sont faites régulièrement, améliorent la concentration et réduisent le stress. Un médecin du travail explique : “Celui qui mange sainement et dort suffisamment est plus performant.” Si davantage d'entreprises introduisaient de telles mesures, il y aurait moins de salariés malades. »

\textbf{Justifications} : passif allemand rendu par « on » (phrase 1) et par un passif français (phrase 2) ; Konjunktiv II rendu par le conditionnel (« introduisaient... il y aurait ») ; \all{gäbe es} = « il y aurait » (forme irrégulière de \all{geben}).
\textbf{Conclusion} : la version exige de repérer le passif et l'irréel avant de traduire, puis de formuler un français naturel.
\end{exoresolu}

\courssub{Méthode 4 : Réinvestir la grammaire du chapitre (passif, relatives, irréel)}

\begin{methode}[Transformer et manipuler les structures]
\begin{enumerate}
\item \textbf{Passif} : actif \all{Die Firma verschmutzt den Fluss} $\rightarrow$ passif \all{Der Fluss wird (von der Firma) verschmutzt}. Le COD de l'actif devient sujet ; le verbe se conjugue avec \all{werden + Partizip II}. Au passé : \all{Der Fluss wurde verschmutzt} (Präteritum) ou \all{ist verschmutzt worden} (Perfekt).
\item \textbf{Relatives} : le pronom relatif prend le \cle{genre et le nombre} de son antécédent et le \cle{cas} de sa fonction dans la relative : \all{der Wald, der... / den... / dem...} ; verbe en fin de proposition.
\item \textbf{Irréel (Konjunktiv II)} : pour les souhaits et hypothèses : \all{Wenn ich mehr Zeit hätte, würde ich Sport treiben} ; formes à connaître par cœur : \all{hätte, wäre, gäbe, könnte, sollte, müsste, würde + Infinitiv}.
\item \textbf{Vérifier systématiquement} : place du verbe (2\textsuperscript{e} position en principale, finale en subordonnée), accord du participe au passif (jamais d'accord en allemand !), majuscules.
\end{enumerate}
\end{methode}

\begin{exoresolu}[Transformations grammaticales]
\textbf{Énoncé (type Bac)} :
\begin{enumerate}
\item Mettez au passif : \all{Die Jugendlichen sammeln den Müll am Fluss.}
\item Complétez avec le pronom relatif correct : \all{Der Arzt, \_\_\_\_ wir gestern besucht haben, arbeitet im Krankenhaus von Douala.}
\item Mettez à l'irréel : \all{Wenn wir weniger Plastik kaufen, schützen wir die Umwelt.}
\end{enumerate}
\tcblower
\textbf{Solutions détaillées} :

\emph{1. Passif} : le COD \all{den Müll} devient sujet au nominatif : \all{der Müll} ; \all{sammeln} $\rightarrow$ \all{wird gesammelt} (présent, 3\textsuperscript{e} personne du singulier car \all{der Müll} est singulier). Le sujet actif devient complément d'agent avec \all{von} : \all{Der Müll am Fluss wird von den Jugendlichen gesammelt.} (« Les déchets au bord du fleuve sont ramassés par les jeunes. »)

\emph{2. Relatif} : antécédent \all{der Arzt} (masculin singulier) ; dans la relative, il est COD de \all{besuchen} (\all{wir haben wen besucht?}) $\rightarrow$ accusatif masculin : \all{den}. Réponse : \all{Der Arzt, den wir gestern besucht haben, arbeitet im Krankenhaus von Douala.} (« Le médecin que nous avons consulté hier travaille à l'hôpital de Douala. ») Piège : ne pas écrire \all{der} sous l'influence du français « que ».

\emph{3. Irréel} : présent $\rightarrow$ Konjunktiv II : \all{kaufen} $\rightarrow$ \all{kauften} (ou \all{würden kaufen}) ; \all{schützen} $\rightarrow$ \all{würden schützen}. Réponse : \all{Wenn wir weniger Plastik kauften, würden wir die Umwelt schützen.} (« Si nous achetions moins de plastique, nous protégerions l'environnement. ») Vérification : verbe conjugué en fin de la subordonnée introduite par \all{wenn}, et en 2\textsuperscript{e} position de la principale (\all{würden wir}).
\textbf{Conclusion} : chaque transformation exige deux contrôles : la forme (morphologie) et la place du verbe (syntaxe).
\end{exoresolu}

\courssub{Méthode 5 : Rédiger un tract (Flugblatt) de sensibilisation}

\begin{methode}[Concevoir un tract efficace]
\begin{enumerate}
\item \textbf{Un slogan} court et frappant, souvent à l'impératif ou avec \all{wir} : \all{Rettet unseren Fluss!} / \all{Gesund leben -- besser leben!}
\item \textbf{Des phrases courtes} : un tract se lit en 30 secondes. Éviter les subordonnées longues.
\item \textbf{Des impératifs} pour l'appel à l'action : \all{Trink mehr Wasser!} (tu), \all{Bringen Sie Ihren Müll zur Sammelstelle!} (vous poli), \all{Machen wir mit!} (nous).
\item \textbf{Des listes à puces} avec des infinitifs ou des noms : \all{Weniger Plastik -- mehr Leben!}
\item \textbf{Des informations pratiques} à la fin : date, lieu, contact (\all{Wann? Wo? Wer?}).
\end{enumerate}
\textbf{Impératifs à maîtriser} : \all{Komm! / Kommt! / Kommen Sie!} ; \all{Sei vorsichtig! / Seid vorsichtig! / Seien Sie vorsichtig!} (seul \all{sein} a une forme spéciale \all{Seien Sie}).
\end{methode}

\begin{exoresolu}[Rédaction d'un tract]
\textbf{Énoncé (type Bac)} : Ton club de santé organise une « Journée du sport et de l'alimentation saine » au lycée. Rédige un tract (8--10 lignes) pour inviter les élèves : slogan, trois arguments, programme, appel à participer.
\tcblower
\textbf{Raisonnement} : registre direct et motivant ; impératifs à la 2\textsuperscript{e} personne du pluriel (\all{kommt, macht, bleibt}) car on s'adresse aux élèves ; phrases courtes ; informations pratiques (jour, heure, lieu).

\textbf{Tract rédigé} :

\all{\textbf{BEWEGE DICH -- BLEIB GESUND!}}

\all{Liebe Schülerinnen und Schüler!}

\all{Sport macht Spaß, stärkt den Körper und hilft beim Lernen. Wer sich regelmäßig bewegt, schläft besser und ist weniger gestresst. Und wer Obst statt Chips isst, tut etwas für seine Gesundheit!}

\all{Deshalb laden wir euch ein zu unserem \textbf{Tag des Sports und der gesunden Ernährung}:}

\all{-- 9.00 Uhr: Fußballturnier auf dem Schulhof}

\all{-- 11.00 Uhr: Vortrag über gesunde Ernährung}

\all{-- 13.00 Uhr: Gemeinsames Mittagessen mit Obst und Gemüse}

\all{\textbf{Wann?} Samstag, den 15. März \quad \textbf{Wo?} Im Schulhof eures Gymnasiums}

\all{\emph{Kommt zahlreich -- eure Gesundheit wartet nicht!}}

\textbf{Justifications} : impératifs corrects (\all{bewege dich} réfléchi à la 2\textsuperscript{e} pers. du singulier, \all{kommt} au pluriel) ; \all{Wer sich regelmäßig bewegt, schläft besser} = tournure « celui qui... » avec verbe en 2\textsuperscript{e} position de la principale ; \all{Obst statt Chips} = \all{statt} + génitif/usage courant ; majuscules à tous les noms (\all{Körper, Obst, Mittagessen, Schulhof}).
\textbf{Conclusion} : le tract remplit les quatre critères : slogan, arguments, programme, appel à l'action.
\end{exoresolu}

\courssub{Méthode 6 : Relire et corriger sa production}

\begin{methode}[La relecture en cinq passages]
\begin{enumerate}
\item \textbf{Passage 1 -- Majuscules} : tout nom allemand prend une majuscule (\all{die Gesundheit, das Wasser, der Sport}), même au milieu de la phrase.
\item \textbf{Passage 2 -- Place du verbe} : verbe conjugué en 2\textsuperscript{e} position dans la principale ; en fin dans les subordonnées (\all{weil, dass, wenn, obwohl}) et les relatives ; participe/infinitif/particule en fin de principale.
\item \textbf{Passage 3 -- Cas et articles} : vérifier chaque groupe nominal : sujet (nominatif), COD (accusatif), COI et prépositions à cas fixe (datif : \all{mit, von, zu, bei, nach, aus}).
\item \textbf{Passage 4 -- Orthographe} : Umlaute (\all{ä, ö, ü}), \all{ß} après voyelle longue (\all{groß, Fuß}), mais \all{ss} après voyelle brève (\all{muss, Wasser}).
\item \textbf{Passage 5 -- Sens} : aucun mot français oublié, aucun calque (\all{*Ich bin einverstanden mit} $\rightarrow$ \all{Ich bin mit ... einverstanden}).
\end{enumerate}
\end{methode}

\begin{exoresolu}[Correction d'un texte fautif]
\textbf{Énoncé (type Bac)} : Le texte suivant contient huit erreurs. Retrouvez-les et corrigez-les.

\all{Viele menschen in unserer Stadt trinkt zu wenig Wasser. Die Ärzte sagt, dass man mindestens zwei Liter pro Tag trinken soll. Wenn man genug Wasser trinken würde, man wäre weniger müde. Das Gesundheit ist wichtiger als Geld, weil man kann sie nicht kaufen.}
\tcblower
\textbf{Corrections détaillées} :
\begin{enumerate}
\item \all{menschen} $\rightarrow$ \all{Menschen} : majuscule obligatoire à tout nom.
\item \all{trinkt} $\rightarrow$ \all{trinken} : le sujet \all{viele Menschen} est pluriel.
\item \all{sagt} $\rightarrow$ \all{sagen} : sujet pluriel \all{die Ärzte}.
\item \all{Wenn man genug Wasser trinken würde, man wäre...} $\rightarrow$ \all{Wenn man genug Wasser trinken würde, wäre man weniger müde} : après la subordonnée (qui occupe la 1\textsuperscript{re} position), le verbe de la principale vient immédiatement : inversion sujet-verbe obligatoire.
\item \all{Das Gesundheit} $\rightarrow$ \all{Die Gesundheit} : \all{die Gesundheit} est féminin.
\item \all{weil man kann sie nicht kaufen} $\rightarrow$ \all{weil man sie nicht kaufen kann} : verbe conjugué en fin de subordonnée ; l'infinitif \all{kaufen} précède le modal \all{kann} en position finale.
\end{enumerate}
\textbf{Texte corrigé} :

\all{Viele Menschen in unserer Stadt trinken zu wenig Wasser. Die Ärzte sagen, dass man mindestens zwei Liter pro Tag trinken soll. Wenn man genug Wasser trinken würde, wäre man weniger müde. Die Gesundheit ist wichtiger als Geld, weil man sie nicht kaufen kann.}

\textbf{Conclusion} : la relecture méthodique (majuscules, verbe, cas, orthographe) permet de récupérer des points facilement perdus.
\end{exoresolu}

\begin{attention}[Les 5 erreurs qui coûtent des points au Bac]
\begin{enumerate}
\item \textbf{La majuscule oubliée} : écrire \all{*die gesundheit}, \all{*das wasser}. En allemand, \emph{tout} nom prend la majuscule. C'est la faute la plus sanctionnée en production écrite.
\item \textbf{Le verbe mal placé} : \all{*Weil ich bin krank} au lieu de \all{Weil ich krank bin} ; \all{*Gestern ich habe Sport gemacht} au lieu de \all{Gestern habe ich Sport gemacht}. Principale = verbe en 2\textsuperscript{e} position ; subordonnée = verbe en fin.
\item \textbf{Les faux amis et calques du français} : \all{*aktual} ne veut pas dire « actuel » (\all{aktuell} oui, mais \all{die Aktualität} = l'actualité) ; « passer du temps » = \all{Zeit verbringen} (pas \all{*passieren}) ; « assister à » = \all{teilnehmen an} ; \all{der Betrieb} = l'entreprise, pas l'opération ; \all{become} n'existe pas : « devenir » = \all{werden}.
\item \textbf{Les cas confondus} : \all{*mit den Arzt} au lieu de \all{mit dem Arzt} ; \all{*der Arzt, der wir besucht haben} au lieu de \all{den wir besucht haben}. Toujours identifier la fonction avant de choisir l'article.
\item \textbf{Le passif et le Konjunktiv II mal formés} : \all{*Der Fluss ist verschmutzt} (état) ne remplace pas \all{Der Fluss wird verschmutzt} (action) ; \all{*wenn ich würde haben} est impossible : on dit \all{wenn ich hätte}. Formes à réciter : \all{hätte, wäre, gäbe, könnte, sollte, müsste}.
\end{enumerate}
\textbf{Réflexe final} : avant de rendre la copie, relire une dernière fois uniquement pour les verbes et les majuscules : c'est là que se gagnent les derniers points.
\end{attention}

\newpage

\coursec{Exercices}

\courssub{Je vérifie mes connaissances}

\begin{exercice}[QCM : grammaire et vocabulaire]
Entoure la bonne réponse. (1 point par question)
\begin{enumerate}
\item \all{Der Wald \_\_\_\_\_\_ durch den Sturm zerstört.} \\
a) \all{wird} \quad b) \all{werden} \quad c) \all{ist} \quad d) \all{hat}
\item \all{Das ist der Arzt, \_\_\_\_\_\_ ich vertraue.} \\
a) \all{den} \quad b) \all{dem} \quad c) \all{der} \quad d) \all{die}
\item Le contraire de \all{gesund} est : \\
a) \all{müde} \quad b) \all{krank} \quad c) \all{stark} \quad d) \all{fit}
\item \all{Wenn ich mehr Zeit hätte, \_\_\_\_\_\_ ich mehr Sport.} \\
a) \all{mache} \quad b) \all{machte} \quad c) \all{würde machen} \quad d) \all{habe gemacht}
\item \all{die Mülltrennung} signifie : \\
a) la collecte des ordures \quad b) le tri des déchets \quad c) la pollution \quad d) le recyclage de l'eau
\end{enumerate}
\end{exercice}

\begin{exercice}[Vrai ou faux ? Justifie ta réponse]
Réponds par \all{richtig} ou \all{falsch} et justifie chaque réponse par une phrase complète en allemand.
\begin{enumerate}
\item Au passif Präsens, on conjugue \all{haben} + Partizip II.
\item Dans une relative introduite par une préposition, le pronom relatif se met au cas exigé par la préposition.
\item \all{Wenn ich du wäre, würde ich zum Arzt gehen} est une phrase irréelle au présent.
\item Un article informatif commence toujours par la conclusion de l'auteur.
\item Le nom \all{die Umwelt} est féminin.
\end{enumerate}
\end{exercice}

\begin{exercice}[Vocabulaire : environnement, santé, bien-être]
Complète chaque phrase avec le mot qui convient, choisi dans la liste : \all{der Abfall, die Luftverschmutzung, sich erholen, die Krankenkasse, der Stress, die Bewegung}.
\begin{enumerate}
\item Nach der Krankheit muss er \_\_\_\_\_\_.
\item \_\_\_\_\_\_ ist ein großes Problem in den Großstädten wie Douala.
\item Wer viel Sport treibt, hat genug \_\_\_\_\_\_.
\item In Deutschland zahlt \_\_\_\_\_\_ einen Teil der Arztkosten.
\item Man soll \_\_\_\_\_\_ nicht in den Fluss werfen.
\item Zu viel \_\_\_\_\_\_ macht krank.
\end{enumerate}
\end{exercice}

\begin{exercice}[Conjugaison : passif et Konjunktiv II]
Complète le tableau avec la forme correcte du verbe entre parenthèses.
\begin{center}
\begin{tabular}{|l|l|l|}
\hline
\rowcolor{popblueL} Phrase & Temps / Mode & Ta réponse \\
\hline
\all{Die Straße \_\_\_\_\_\_ (reparieren).} & Passiv Präsens & \\
\hline
\all{Die Bäume \_\_\_\_\_\_ (pflanzen).} & Passiv Perfekt & \\
\hline
\all{Der Patient \_\_\_\_\_\_ (operieren).} & Passiv Präteritum & \\
\hline
\all{Ich \_\_\_\_\_\_ (können) mehr Sport machen.} & Konjunktiv II & \\
\hline
\all{Wir \_\_\_\_\_\_ (sein) froh, wenn ...} & Konjunktiv II & \\
\hline
\end{tabular}
\end{center}
\end{exercice}

\courssub{Je m'entraîne}

\begin{exercice}[Thème : phrases à traduire]
Traduis les phrases suivantes en allemand. Attention au passif, aux relatives et au Konjunktiv II.
\begin{enumerate}
\item Les déchets doivent être triés. (passif + modal)
\item Le médecin, chez qui je vais chaque année, est très gentil. (relative avec préposition)
\item À ta place, je ferais plus de sport. (Konjunktiv II de conseil)
\item Beaucoup d'arbres ont été plantés dans notre quartier à Yaoundé.
\item Si nous protégions l'environnement, nous serions en meilleure santé.
\end{enumerate}
\end{exercice}

\begin{exercice}[Transformation : de l'actif au passif]
Mets les phrases suivantes au passif, au même temps.
\begin{enumerate}
\item \all{Die Schüler reinigen den Schulhof.}
\item \all{Die Firma hat den Fluss verschmutzt.}
\item \all{Man muss die Umwelt schützen.}
\item \all{Der Arzt untersuchte den Patienten.}
\item \all{Die Stadt wird neue Parks bauen.}
\end{enumerate}
\end{exercice}

\begin{exercice}[Texte à trous : un article informatif]
Complète l'article suivant avec les mots de la liste. Attention : deux mots sont en trop.\\
Liste : \all{werden, der, dem, würden, mehr, weniger, Gesundheit, Umwelt, können}

\begin{textebox}[Gesund leben in der Großstadt]
In vielen Städten leiden die Menschen unter Stress und schlechter Luft. Die \_\_\_\_\_\_ der Einwohner ist in Gefahr. Wenn die Autos weniger Benzin verbrauchen \_\_\_\_\_\_, wäre die Luft sauberer. Außerdem sollten \_\_\_\_\_\_ Parks angelegt \_\_\_\_\_\_. In \_\_\_\_\_\_ Park kann man sich gut erholen. Jeder Einwohner kann etwas für die \_\_\_\_\_\_ tun: zu Fuß gehen, den Müll trennen und \_\_\_\_\_\_ fernsehen.
\end{textebox}
\end{exercice}

\begin{exercice}[Appariements : le vocabulaire du bien-être]
Relie chaque mot allemand à sa traduction française.
\begin{center}
\begin{tabular}{|l|l|}
\hline
\rowcolor{popblueL} Deutsch & Français \\
\hline
1. \all{sich entspannen} & a. le mode de vie \\
\hline
2. \all{die Ernährung} & b. se détendre \\
\hline
3. \all{der Lebensstil} & c. l'alimentation \\
\hline
4. \all{die Sucht} & d. équilibré \\
\hline
5. \all{ausgewogen} & e. la dépendance \\
\hline
6. \all{vorbeugen} & f. prévenir \\
\hline
\end{tabular}
\end{center}
\end{exercice}

\courssub{Je me prépare au Bac}

\begin{exercice}[Compréhension écrite et questions de langue]
Lis le texte suivant, puis réponds aux questions.

\begin{textebox}[Jugend und Gesundheit in der Schweiz]
In der Schweiz bewegen sich viele Jugendliche zu wenig. Eine Untersuchung, die in mehreren Kantonen durchgeführt wurde, zeigt: Fast jeder dritte Jugendliche treibt weniger als eine Stunde Sport pro Woche. Die Gründe sind bekannt: Viele junge Leute sitzen stundenlang vor dem Handy oder dem Computer. Dabei wäre es so einfach, etwas zu ändern! Wenn die Schüler öfter zu Fuß zur Schule gingen, würden sie sich besser fühlen. Auch die Ernährung spielt eine große Rolle: Fast Food und süße Getränke werden von vielen Jugendlichen täglich konsumiert. Gesundheitsexperten empfehlen, dass in den Schulen mehr über gesunde Ernährung informiert werden soll. In einigen Schulen wurden schon Obstautomaten aufgestellt, die von den Schülern gut angenommen werden. „Wenn ich weniger Zeit am Handy verbringen würde, hätte ich mehr Energie für Sport“, sagt ein 17-jähriger Schüler aus Bern.
\end{textebox}

\textbf{A. Compréhension (réponds en allemand, par des phrases complètes)}
\begin{enumerate}
\item Was zeigt die Untersuchung in der Schweiz?
\item Nenne zwei Gründe für die Bewegungsarmut der Jugendlichen.
\item Was empfehlen die Gesundheitsexperten?
\end{enumerate}
\textbf{B. Vrai ou faux ? Justifie par une citation du texte.}
\begin{enumerate}
\item Alle Jugendlichen in der Schweiz machen genug Sport.
\item Die Obstautomaten haben keinen Erfolg.
\end{enumerate}
\textbf{C. Questions de langue}
\begin{enumerate}
\item Relève deux phrases au passif et justifie le choix du temps.
\item \all{„eine Untersuchung, die in mehreren Kantonen durchgeführt wurde“} : analyse la relative (pronom, cas, antécédent).
\item Relève une phrase au Konjunktiv II et explique sa valeur (irréel du présent).
\end{enumerate}
\end{exercice}

\begin{exercice}[Thème et version]
\textbf{A. Thème.} Traduis en allemand (environ 5 phrases) :

\begin{textebox}[Texte à traduire]
Dans ma ville, l'environnement n'est pas assez protégé. Les déchets sont jetés partout dans les rues. Si les habitants triaient leurs ordures, la ville serait plus propre. Le maire, à qui j'ai écrit une lettre, a promis des poubelles nouvelles. À ta place, je participerais aussi à la journée de nettoyage de samedi.
\end{textebox}

\textbf{B. Version.} Traduis en français le texte suivant (environ 10 lignes) :

\begin{textebox}[Mülltrennung in Deutschland]
In Deutschland wird der Müll seit vielen Jahren getrennt. In fast jedem Haushalt stehen mehrere Tonnen: eine für Papier, eine für Plastik und eine für den Restmüll. Glas wird in großen Containern gesammelt, die in jeder Straße zu finden sind. Wer den Müll nicht richtig trennt, muss manchmal eine Strafe zahlen. Viele Deutsche sind stolz auf dieses System, denn viele Abfälle können recycelt werden. Experten sagen: Wenn in allen Ländern der Müll so getrennt würde, gäbe es weniger Umweltprobleme. Auch in Kamerun wird in einigen Städten schon über Mülltrennung diskutiert.
\end{textebox}

\textbf{C. Justification grammaticale.} Après ta traduction, explique en français :
\begin{enumerate}
\item le temps du passif dans \all{wird der Müll getrennt} et dans \all{Glas wird gesammelt} ;
\item la forme et la valeur de \all{würde} et \all{gäbe es} dans la phrase des experts.
\end{enumerate}
\end{exercice}

\begin{exercice}[Production écrite type Bac et reprise d'erreurs]
\textbf{A. Production écrite.} Rédige en allemand un article informatif de \textbf{120 à 150 mots} intitulé \all{„Gesund leben in meiner Stadt“}. Respecte le plan imposé :
\begin{enumerate}
\item \textbf{Einleitung} : présente le problème (santé et environnement dans ta ville) ;
\item \textbf{Hauptteil} : décris deux problèmes concrets et propose deux solutions (utilise au moins une phrase au passif et une phrase au Konjunktiv II) ;
\item \textbf{Schluss} : donne ton avis personnel et un conseil aux jeunes.
\end{enumerate}
Barème indicatif : respect du plan et du nombre de mots (4 pts), correction grammaticale : passif, relatives, Konjunktiv II, ordre des mots (8 pts), richesse du vocabulaire (4 pts), orthographe et majuscules (4 pts).

\textbf{B. Reprise d'erreurs.} Le tract suivant contient \textbf{10 fautes} (majuscules, ordre des mots, cas, passif, conjugaison). Trouve-les et corrige-les.

\begin{textebox}[Tract à corriger]
Liebe bürger von Yaoundé!\\
Unsere Stadt wird jeden Tag mehr verschmutzt. Der müll wird in die Flüsse geworfen, und viele Menschen werden krank. Wir müssen etwas tun! Am Samstag organisiert unser Verein eine Aktion: die Straßen werden von uns gereinigt werden. Wenn du kommst, wir würden sehr froh sein. Bring deine Freunde mit! Für mehr Informationen ruft der Präsident an. Zusammen können wir unsere umwelt schützen und gesünder leben!
\end{textebox}
\end{exercice}

\newpage

\coursec{Corrigés des exercices}

\coursec{Corrigés détaillés des exercices — Leçon AL15}

\courssub{Exercice 1 : QCM — Grammaire et vocabulaire}
\begin{corrige}[Exercice 1]
\begin{enumerate}
\item \textbf{a) \all{wird}}. \all{Der Wald wird durch den Sturm zerstört.} (Passif procès \emph{Vorgangspassiv} au présent : \all{werden} + Partizip II). \all{ist} marquerait un passif d'état (\emph{Zustandspassiv}), moins naturel avec un agent \all{durch den Sturm} qui insiste sur l'action.
\item \textbf{b) \all{dem}}. \all{Das ist der Arzt, dem ich vertraue.} Le verbe \all{vertrauen} (faire confiance) régit le datif (\all{vertrauen + Dativ}), donc le pronom relatif est \all{dem} (masculin singulier datif).
\item \textbf{c) \all{krank}}. \all{gesund} (en bonne santé) ↔ \all{krank} (malade). \all{müde} = fatigué, \all{stark} = fort, \all{fit} = en forme.
\item \textbf{d) \all{würde machen}}. Phrase irréelle au présent (\emph{Konjunktiv II}) : \all{Wenn ich mehr Zeit hätte, würde ich mehr Sport machen.} Pour les verbes faibles comme \all{machen}, on utilise la forme \all{würde + infinitif} (et non \all{machte}, qui est l'indicatif prétérit).
\item \textbf{e) \all{die Mülltrennung}}. \all{die Mülltrennung} (féminin) = séparation/tri des ordures. \all{die Müllabfuhr} = collecte, \all{die Umweltverschmutzung} = pollution, \all{die Wasseraufbereitung} = traitement de l'eau.
\end{enumerate}
\end{corrige}

\courssub{Exercice 2 : Vrai ou faux ? Justifications}
\begin{corrige}[Exercice 2]
\begin{enumerate}
\item \all{falsch}. \all{Im Passiv Präsens konjugiert man „werden“ + Partizip II.} (Ex. : \all{Die Straße wird repariert.}) \all{haben} s'utilise pour le parfait \emph{actif} ; le passif parfait se forme avec \all{sein} + Partizip II + \all{worden} (ex. : \all{ist repariert worden}).
\item \all{richtig}. \all{In einem Relativsatz mit Präposition steht das Relativpronomen im Kasus, den die Präposition fordert.} (Ex. : \all{der Arzt, dem ich vertraue} — \all{vertrauen} + Dativ).
\item \all{richtig}. \all{„Wenn ich du wäre, würde ich zum Arzt gehen“} est une phrase conditionnelle irréelle au présent (\emph{Konjunktiv II} : \all{wäre} / \all{würde gehen}), exprimant un conseil ou une hypothèse non réalisée.
\item \all{falsch}. \all{Ein informativer Artikel beginnt mit einer Einleitung (Thema, These), nicht mit dem Fazit.} La structure standard est : \all{Einleitung – Hauptteil – Schluss}.
\item \all{richtig}. \all{die Umwelt} est un nom féminin (\all{die}), pluriel \all{die Umwelten} (rare), souvent utilisé au singulier.
\end{enumerate}
\end{corrige}

\courssub{Exercice 3 : Vocabulaire — Environnement, santé, bien-être}
\begin{corrige}[Exercice 3]
\begin{enumerate}
\item Nach der Krankheit muss er \textbf{\all{sich erholen}}. (se remettre / récupérer)
\item \textbf{\all{Die Luftverschmutzung}} ist ein großes Problem in den Großstädten wie Douala. (la pollution de l'air)
\item Wer viel Sport treibt, hat genug \textbf{\all{Bewegung}}. (activité physique / mouvement)
\item In Deutschland zahlt \textbf{\all{die Krankenkasse}} einen Teil der Arztkosten. (la caisse d'assurance maladie)
\item Man soll \textbf{\all{den Abfall}} nicht in den Fluss werfen. (les déchets / ordures)
\item Zu viel \textbf{\all{Stress}} macht krank. (le stress)
\end{enumerate}
\end{corrige}

\courssub{Exercice 4 : Conjugaison — Passif et Konjunktiv II}
\begin{corrige}[Exercice 4]
\begin{center}
\begin{tabular}{|l|l|l|}
\hline
\rowcolor{popblueL} Phrase & Temps / Mode & Réponse \\
\hline
\all{Die Straße \textbf{wird repariert}.} & Passiv Präsens & \all{werden} + Partizip II \\
\hline
\all{Die Bäume \textbf{sind gepflanzt worden}.} & Passiv Perfekt & \all{sein} (présent) + Partizip II + \all{worden} \\
\hline
\all{Der Patient \textbf{wurde operiert}.} & Passiv Präteritum & \all{werden} (prétérit) + Partizip II \\
\hline
\all{Ich \textbf{könnte} mehr Sport machen.} & Konjunktiv II & \all{können} → \all{könnte} (forme simple) ou \all{würde machen} \\
\hline
\all{Wir \textbf{wären} froh, wenn ...} & Konjunktiv II & \all{sein} → \all{wären} \\
\hline
\end{tabular}
\end{center}
\emph{Note :} Au passif parfait, l'auxiliaire est \all{sein} (et non \all{haben}) + Partizip II + \all{worden} (sans \emph{ge-}).
\end{corrige}

\courssub{Exercice 5 : Thème — Phrases à traduire}
\begin{corrige}[Exercice 5]
\begin{enumerate}
\item \all{Der Müll muss getrennt werden.} (Passif + modal : \all{müssen} + \all{getrennt werden})
\item \all{Der Arzt, \textbf{zu dem} ich jedes Jahr gehe, ist sehr nett.} (Relatif + préposition : \all{zu} + Dativ → \all{zu dem} ; ordre des mots : verbe à la fin du relatif). \emph{Variante :} \all{Der Arzt, bei dem ich mich jedes Jahr behandeln lasse, ...}
\item \all{An deiner Stelle würde ich mehr Sport machen.} (Konjunktiv II de conseil : \all{an deiner Stelle} + \all{würde + infinitif}). Variante : \all{Wenn ich du wäre, würde ich mehr Sport machen.}
\item \all{In unserem Viertel in Yaoundé sind viele Bäume gepflanzt worden.} (Passif Perfekt : \all{sind ... gepflanzt worden}). Variante Präteritum : \all{wurden viele Bäume gepflanzt.}
\item \all{Wenn wir die Umwelt schützen würden, wären wir gesünder.} (Irréel présent : \all{würden schützen} / \all{wären} ; \all{schützen} verbe faible → \all{würde}-Forme).
\end{enumerate}
\end{corrige}

\courssub{Exercice 6 : Transformation — De l'actif au passif}
\begin{corrige}[Exercice 6]
\begin{enumerate}
\item \all{Der Schulhof wird (von den Schülern) gereinigt.} (Présent)
\item \all{Der Fluss ist (von der Firma) verschmutzt worden.} (Perfekt : \all{hat verschmutzt} → \all{ist ... worden})
\item \all{Die Umwelt muss geschützt werden.} (Modal + passif infinitif : \all{Man} disparaît, \all{muss} reste)
\item \all{Der Patient wurde (von dem Arzt) untersucht.} (Prétérit)
\item \all{Neue Parks werden (von der Stadt) gebaut werden.} (Futur I passif : \all{werden} + Partizip II + \all{werden})
\end{enumerate}
\end{corrige}

\courssub{Exercice 7 : Texte à trous — Article informatif}
\begin{textebox}[Gesund leben in der Großstadt — Texte complété]
In vielen Städten leiden die Menschen unter Stress und schlechter Luft. Die \textbf{\all{Gesundheit}} der Einwohner ist in Gefahr. Wenn die Autos weniger Benzin verbrauchen \textbf{\all{würden}}, wäre die Luft sauberer. Außerdem sollten \textbf{\all{mehr}} Parks angelegt \textbf{\all{werden}}. In \textbf{\all{dem}} Park kann man sich gut erholen. Jeder Einwohner kann etwas für die \textbf{\all{Umwelt}} tun: zu Fuß gehen, den Müll trennen und \textbf{\all{weniger}} fernsehen.
\end{textebox}

\begin{corrige}[Exercice 7]
\textbf{Mots proposés (dont 2 en trop) :} \all{der, können, Gesundheit, würden, mehr, werden, dem, Umwelt, weniger}. \\
\textbf{Mots en trop :} \all{der}, \all{können}. \\
\textbf{Justifications :}
\begin{itemize}
\item \all{Gesundheit} : sujet, féminin, nominatif singulier (\all{die Gesundheit}).
\item \all{würden} : Konjunktiv II dans la proposition \emph{si} (\all{wenn}-Satz), verbe faible \all{verbrauchen} → \all{würden verbrauchen}.
\item \all{mehr} : adjectif comparatif (attribut / déterminant) pour \all{Parks}.
\item \all{werden} : passif modal après \all{sollten} (\all{sollten ... angelegt werden}).
\item \all{dem} : Datif singulier masculin après préposition \all{In} + article défini (\all{in dem} = \all{im}).
\item \all{Umwelt} : accusatif féminin singulier après préposition \all{für} (\all{für die Umwelt}).
\item \all{weniger} : adverbe comparatif modifiant \all{fernsehen}.
\end{itemize}
\end{corrige}

\courssub{Exercice 8 : Appariements — Vocabulaire du bien-être}
\begin{corrige}[Exercice 8]
\begin{center}
\begin{tabular}{|l|l|}
\hline
\rowcolor{popblueL} Deutsch & Français \\
\hline
1. \all{sich entspannen} & \textbf{b. se détendre} \\
\hline
2. \all{die Ernährung} & \textbf{c. l'alimentation} \\
\hline
3. \all{der Lebensstil} & \textbf{a. le mode de vie} \\
\hline
4. \all{die Sucht} & \textbf{e. la dépendance} \\
\hline
5. \all{ausgewogen} & \textbf{d. équilibré} \\
\hline
6. \all{vorbeugen} & \textbf{f. prévenir} \\
\hline
\end{tabular}
\end{center}
\end{corrige}

\courssub{Exercice 9 : Compréhension écrite et questions de langue}
\begin{corrige}[Exercice 9]
\textbf{A. Compréhension}
\begin{enumerate}
\item \all{Die Untersuchung zeigt, dass fast jeder dritte Jugendliche weniger als eine Stunde Sport pro Woche treibt.}
\item \all{Erstens sitzen viele junge Leute stundenlang vor dem Handy oder Computer. Zweitens konsumieren sie täglich Fast Food und süße Getränke.}
\item \all{Die Gesundheitsexperten empfehlen, dass in den Schulen mehr über gesunde Ernährung informiert werden soll.}
\end{enumerate}

\textbf{B. Vrai ou faux ?}
\begin{enumerate}
\item \all{falsch}. Beleg: \all{„Fast jeder dritte Jugendliche treibt weniger als eine Stunde Sport pro Woche.“}
\item \all{falsch}. Beleg: \all{„die von den Schülern gut angenommen werden.“}
\end{enumerate}

\textbf{C. Questions de langue}
\begin{enumerate}
\item \textbf{Passifs relevés dans le texte support :}
\begin{itemize}
\item \all{„... eine Untersuchung, die in mehreren Kantonen \textbf{durchgeführt wurde}“} → Passif Prétérit (\emph{Vorgangspassiv}) : action passée accomplie (l'enquête a eu lieu).
\item \all{„Fast Food und süße Getränke \textbf{werden ... konsumiert}“} → Passif Présent : fait général, habitude actuelle.
\item \all{„\textbf{soll ... informiert werden}“} → Passif modal (infinitif) : recommandation, obligation morale.
\item \all{„Obstautomaten \textbf{wurden ... aufgestellt}“} → Passif Prétérit : action passée ponctuelle.
\item \all{„die von den Schülern gut \textbf{angenommen werden}“} → Passif Présent : état actuel, succès continu.
\end{itemize}
\item \textbf{Analyse de la relative :} \all{„eine Untersuchung, die in mehreren Kantonen durchgeführt wurde“}
\begin{itemize}
\item Antécédent : \all{eine Untersuchung} (féminin, singulier).
\item Pronom relatif : \all{die} (Nominatif féminin singulier, fonction = sujet de la relative).
\item Cas : Nominatif (le pronom est le sujet de \all{wurde durchgeführt}).
\item Verbe : \all{wurde durchgeführt} (Passif Prétérit), placé en fin de proposition.
\end{itemize}
\item \textbf{Konjunktiv II relevé :} \all{„Wenn die Schüler öfter zu Fuß zur Schule \textbf{gingen}, \textbf{würden} sie sich besser \textbf{fühlen}.“} (ou phrase de l'élève : \all{„Wenn ich weniger Zeit am Handy \textbf{verbringen würde}, \textbf{hätte} ich mehr Energie...“}).
\textbf{Valeur :} Irréel du présent (\emph{Irrealis der Gegenwart}). La condition n'est pas réalisée dans le présent (les élèves ne vont pas à pied, l'élève passe trop de temps sur son téléphone), on exprime une hypothèse contraire à la réalité.
\end{enumerate}
\end{corrige}

\courssub{Exercice 10 : Thème et Version}
\begin{corrige}[Exercice 10]
\textbf{A. Thème (Français → Allemand) — Proposition}
\all{In meiner Stadt wird die Umwelt nicht genug geschützt. Der Müll wird überall auf den Straßen geworfen. Wenn die Einwohner ihren Müll trennten, wäre die Stadt sauberer. Der Bürgermeister, dem ich einen Brief geschrieben habe, hat neue Mülltonnen versprochen. An deiner Stelle würde ich auch an der Aufräumaktion am Samstag teilnehmen.}

\textbf{B. Version (Allemand → Français) — Texte source et traduction}
\begin{textebox}[Texte source pour la version]
In Deutschland wird der Müll seit vielen Jahren getrennt. In fast jedem Haushalt gibt es mehrere Mülltonnen: eine für Papier, eine für Plastik und eine für Restmüll. Glas wird in großen Containern gesammelt, die in jeder Straße stehen. Wer den Müll nicht richtig trennt, muss manchmal ein Bußgeld zahlen. Viele Deutsche sind stolz auf dieses System, denn viel Müll kann recycelt werden. Experten sagen: „Wenn in allen Ländern der Müll so getrennt würde, gäbe es weniger Umweltprobleme.“ Auch in Kamerun diskutiert man in manchen Städten schon über die Mülltrennung.
\end{textebox}

\textbf{Traduction proposée :}
En Allemagne, les déchets sont triés depuis de nombreuses années. Dans presque chaque foyer, il y a plusieurs poubelles : une pour le papier, une pour le plastique et une pour les déchets résiduels. Le verre est collecté dans de grands conteneurs que l'on trouve dans chaque rue. Celui qui ne trie pas correctement les déchets doit parfois payer une amende. Beaucoup d'Allemands sont fiers de ce système, car beaucoup de déchets peuvent être recyclés. Les experts disent : si dans tous les pays les déchets étaient triés de cette façon, il y aurait moins de problèmes environnementaux. Aussi au Cameroun, dans certaines villes, on discute déjà du tri des déchets.

\textbf{C. Justification grammaticale (extrait du texte source)}
\begin{enumerate}
\item \all{wird der Müll getrennt} / \all{Glas wird gesammelt} / \all{kann ... recycelt werden} : \textbf{Passif Présent (\emph{Vorgangspassiv})}. Forme : \all{werden} (présent) + Partizip II. Emploi : description de processus généraux, habituels, faits permanents (système de tri en Allemagne). Le sujet grammatical (\all{der Müll}, \all{Glas}, \all{viel Müll}) subit l'action.
\item \all{würde} / \all{gäbe es} : \textbf{Konjunktiv II (Irréel du présent)}.
\begin{itemize}
\item \all{würde} : forme du KII de l'auxiliaire \all{werden} (servant à former le KII des verbes faibles : \all{würde trennen}). Ici dans l'apodose (conséquence) de la phrase conditionnelle des experts.
\item \all{gäbe es} : forme du KII de \all{es gibt} (il y a). \all{gäbe} = KII de \all{geben} (3e pers. sg.) + \all{es} (sujet expletif). La phrase \all{„Wenn in allen Ländern der Müll so getrennt würde, gäbe es weniger Umweltprobleme“} exprime une hypothèse irréelle au présent (le tri n'est pas fait partout) et sa conséquence hypothétique.
\end{itemize}
\end{enumerate}
\end{corrige}

\courssub{Exercice 11 : Production écrite type Bac et reprise d'erreurs}
\begin{textebox}[Gesund leben in meiner Stadt — Proposition de rédaction modèle (env. 135 mots)]
\textbf{Einleitung}
In Yaoundé, der Hauptstadt Kameruns, sind Gesundheit und Umwelt eng miteinander verbunden. Leider ist die Situation besorgniserregend: Die Luft ist durch Abgase und Müllverbrennung stark verschmutzt, und viele Viertel haben keine grünen Erholungsräume.

\textbf{Hauptteil}
Ein erstes Problem ist die wilde Müllentsorgung. \textbf{Der Müll wird oft einfach in die Flüsse oder an den Straßenrand geworfen} (Passiv), was Krankheiten wie Cholera begünstigt. \textbf{Wenn mehr öffentliche Mülltonnen aufgestellt würden und die Mülltrennung eingeführt würde} (Konjunktiv II), \textbf{wäre} die Stadt sauberer. Ein zweites Problem ist der Bewegungsmangel. Viele Jugendliche sitzen den ganzen Tag vor Bildschirmen. \textbf{Würde die Stadt sichere Radwege und Sportplätze bauen} (Konjunktiv II), \textbf{könnten} sich die jungen Leute besser erholen.

\textbf{Schluss}
Meiner Meinung nach muss die Stadtverwaltung endlich handeln, aber jeder Bürger trägt auch Verantwortung. Mein Rat an die Jugend: Bewegt euch mehr, esst ausgewogen und setzt euch für eine saubere Umwelt ein – eure Gesundheit hängt davon ab!
\end{textebox}

\begin{corrige}[Exercice 11 — Production écrite]
\textbf{Barème indicatif (sur 20) :}
\begin{itemize}
\item Respect du plan (3 parties), nombre de mots (120–150), cohérence : \textbf{4 pts}
\item Correction grammaticale (passif, relatif, Konjunktiv II, ordre des mots V2/fin de phrase, cas) : \textbf{8 pts}
\item Richesse lexicale (vocabulaire précis : \all{Luftverschmutzung, Mülltrennung, Bewegungsmangel, Erholungsräume, Cholera, Radwege}...) : \textbf{4 pts}
\item Orthographe (majuscules noms, Umlaute, ß, ponctuation) : \textbf{4 pts}
\end{itemize}

\textbf{Points de vigilance (erreurs fréquentes à éviter) :}
\begin{itemize}
\item Oublier la majuscule aux noms (\all{umwelt} → \all{Umwelt}).
\item Ordre des mots : verbe en 2e position phrase principale, à la fin subordonnée (\all{wenn ... würde}).
\item Passif : confusion \all{werden/sein/worden}, omission de \all{worden} au Perfekt.
\item Relatifs : mauvais cas après préposition (\all{mit dem} pas \all{mit den}).
\item Konjunktiv II : utiliser \all{würde + infinitif} pour verbes faibles, ne pas écrire \all{*machte} pour \all{machen}.
\item Faux-amis : \all{aktuell} = actuel (pas \emph{actuellement}), \all{konsequent} = cohérent/conséquent.
\end{itemize}
\end{corrige}

\courssub{Exercice 11 (suite) : Reprise d'erreurs — Tract corrigé}
\begin{textebox}[Tract corrigé — Version finale]
\textbf{Liebe Bürger} von Yaoundé!\\
Unsere Stadt wird jeden Tag \textbf{immer mehr} verschmutzt. Der \textbf{Müll} wird in die Flüsse geworfen, und viele Menschen werden krank. Wir müssen etwas tun! Am Samstag organisiert unser Verein eine Aktion: die Straßen werden \textbf{gereinigt}. Wenn du kommst, \textbf{würden wir uns sehr freuen}. Bring deine Freunde mit! Für mehr Informationen \textbf{ruft den Präsidenten an}. Zusammen können wir unsere \textbf{Umwelt} schützen und gesünder leben!
\end{textebox}

\begin{corrige}[Exercice 11 — Reprise d'erreurs (10 corrections)]
\textbf{Liste des 10 corrections attendues :}
\begin{enumerate}
\item \all{bürger} → \textbf{\all{Bürger}} (majuscule obligatoire pour les noms).
\item \all{mehr} → \textbf{\all{immer mehr}} (style : comparatif progressif) \emph{ou} correction grammaticale stricte : \all{müll} → \textbf{\all{Müll}} (majuscule, voir correction 3).
\item \all{müll} → \textbf{\all{Müll}} (majuscule).
\item \all{werden ... gereinigt werden} → \textbf{\all{werden gereinigt}} (Passif Présent suffisant pour futur proche programmé ; \all{werden ... werden} = Futur I passif, lourd et incorrect ici).
\item \all{wir würden} → \textbf{\all{würden wir}} (ordre V2 en proposition principale après subordonnée \all{wenn}).
\item \all{sehr froh sein} → \textbf{\all{uns sehr freuen}} (idiomatique ; \all{froh sein} exige la préposition \all{über}).
\item \all{ruft der Präsident an} → \textbf{\all{ruft den Präsidenten an}} (Accusatif \all{den Präsidenten} objet direct de \all{anrufen} ; particule \all{an} en fin de phrase).
\item \all{umwelt} → \textbf{\all{Umwelt}} (majuscule nom).
\item \all{Liebe Bürger} → \textbf{\all{Liebe Bürger!}} (ponctuation : point d'exclamation pour l'apostrophe).
\item \all{am samstag} → \textbf{\all{Am Samstag}} (majuscule en début de phrase / nom propre jour).
\end{enumerate}
\emph{Note :} Toute correction grammaticalement justifiée (cas, temps, ordre, majuscule, particule) est acceptée dans la limite de 10 points.
\end{corrige}

\newpage

\coursec{Fiche bilan}

\begin{popfigure}
\begin{center}
\begin{tikzpicture}[mindmap, grow cyclic, every node/.style={concept, font=\small},
root concept/.append style={concept color=popblue, font=\small\bfseries, minimum size=2.6cm},
level 1/.append style={level distance=3.4cm, sibling angle=60, minimum size=2.2cm, font=\footnotesize},
level 2/.append style={level distance=1.9cm, sibling angle=45, minimum size=1.6cm, font=\scriptsize}]
\node[root concept]{Umwelt, Gesundheit, Wohlbefinden}
child[concept color=popgreen]{ node {Lexique}
  child { node {die Umwelt} }
  child { node {die Gesundheit} }
}
child[concept color=poporange]{ node {Passiv}
  child { node {werden + Partizip II} }
  child { node {von / durch} }
}
child[concept color=poppurple]{ node {Relativsätze}
  child { node {der, die, das} }
  child { node {Kasus nach Funktion} }
}
child[concept color=poppink]{ node {Irrealis}
  child { node {Konjunktiv II} }
  child { node {wenn ... würde} }
}
child[concept color=popgold]{ node {Artikel / Flugblatt}
  child { node {Titel, Vorspann, Text} }
  child { node {Slogan, Imperativ} }
}
child[concept color=popblue!60]{ node {Übersetzung}
  child { node {Thème : fr $\to$ all} }
  child { node {Version : all $\to$ fr} }
};
\end{tikzpicture}
\end{center}
\legende{Carte mentale de la leçon : environnement, santé et bien-être}
\end{popfigure}

\begin{aretenir}[L'essentiel de la leçon]
\textbf{1. Lexique clé à connaître par cœur}

\begin{itemize}
\item \all{die Umwelt} (l'environnement, sans pluriel), \all{der Umweltschutz} (la protection de l'environnement), \all{die Verschmutzung, -en} (la pollution), \all{der Müll} (les déchets, sans pluriel), \all{recyceln} (recycler), \all{Energie sparen} (économiser de l'énergie), \all{das Klima} (le climat), \all{die Erderwärmung} (le réchauffement climatique), \all{der Abfall, die Abfälle} (le déchet), \all{die Mülltrennung} (le tri des déchets).
\item \all{die Gesundheit} (la santé, sans pluriel), \all{gesund / krank} (en bonne santé / malade), \all{der Arzt, die Ärzte} (le médecin), \all{das Krankenhaus, die Krankenhäuser} (l'hôpital), \all{die Krankheit, -en} (la maladie), \all{sich erholen} (se rétablir, se reposer), \all{vorbeugen} (prévenir, verbe à particule séparable : \all{Er beugt vor.}), \all{die Impfung, -en} (la vaccination).
\item \all{das Wohlbefinden} (le bien-être, sans pluriel), \all{der Stress} (le stress), \all{sich entspannen} (se détendre), \all{Sport treiben} (faire du sport, sans article), \all{die Ernährung} (l'alimentation, sans pluriel), \all{ausgewogen} (équilibré), \all{genug schlafen} (dormir assez), \all{regelmäßig} (régulièrement).
\end{itemize}

\textbf{2. Grammaire à connaître par cœur}

\begin{center}
\begin{tabularx}{\linewidth}{|l|X|X|}
\hline
\rowcolor{popblueL} Point & Règle & Exemple \\
\hline
Passif (Präsens) & \all{werden} conjugué + Partizip II en fin de phrase ; agent introduit par \all{von} + datif & \all{Der Wald wird geschützt.} « La forêt est protégée. » \\
\hline
Passif au passé & Präteritum : \all{wurde(n)} + Partizip II ; Perfekt : \all{ist / sind} + Partizip II + \all{worden} & \all{Die Fabrik wurde geschlossen.} « L'usine a été fermée. » ; \all{Die Fabrik ist geschlossen worden.} « L'usine a été fermée. » \\
\hline
Proposition relative & pronom relatif = genre et nombre de l'antécédent, cas selon sa fonction dans la relative ; verbe conjugué en fin & \all{Der Arzt, der hier arbeitet, kommt aus Douala.} « Le médecin qui travaille ici vient de Douala. » \\
\hline
Irréel (Konjunktiv II) & \all{hätte}, \all{wäre}, \all{würde} + infinitif en fin ; dans la subordonnée avec \all{wenn}, le verbe conjugué va en fin & \all{Wenn ich Zeit hätte, würde ich Sport treiben.} « Si j'avais le temps, je ferais du sport. » \\
\hline
\end{tabularx}
\end{center}

\textbf{3. Méthodes express}

\begin{itemize}
\item \textbf{Article informatif} : titre accrocheur $\to$ chapô (qui ? quoi ? où ? quand ?) $\to$ corps en paragraphes (faits, chiffres, citations) $\to$ conclusion ouverte. Style neutre, passif fréquent.
\item \textbf{Tract} : slogan court + impératifs (\all{Rette die Umwelt!} « Sauve l'environnement ! »), phrases brèves, pronoms \all{wir / ihr}, appel à l'action final.
\item \textbf{Thème (fr $\to$ all)} : repérer le temps, choisir actif ou passif, placer le verbe conjugué en 2\up{e} position, vérifier genre et cas des noms.
\item \textbf{Version (all $\to$ fr)} : traduire d'abord le sens global, repérer les verbes à particule séparable et les subordonnées, reformuler en français naturel.
\item \textbf{Relecture} : majuscule à tous les noms, place du verbe, accord des articles, Umlaute et \all{ß}.
\end{itemize}
\end{aretenir}

\begin{aretenir}[Checklist avant le Bac]
\begin{itemize}
\item[$\square$] Je connais le lexique de l'environnement, de la santé et du bien-être avec les articles et les pluriels.
\item[$\square$] Je sais former le passif au présent, au Präteritum et au Perfekt (\all{wird / wurde / ist worden} + Partizip II).
\item[$\square$] Je sais introduire l'agent avec \all{von} + datif et l'instrument ou le moyen avec \all{durch} + accusatif.
\item[$\square$] Je choisis correctement le pronom relatif (genre, nombre, cas) et je place le verbe conjugué en fin de relative.
\item[$\square$] Je maîtrise le Konjunktiv II pour exprimer l'irréel : \all{hätte, wäre, würde} + infinitif.
\item[$\square$] Je respecte la place du verbe conjugué en 2\up{e} position et l'infinitif ou le participe en fin de phrase.
\item[$\square$] Je mets une majuscule à tous les noms allemands et j'écris correctement les Umlaute et le \all{ß}.
\item[$\square$] Je connais la structure d'un article informatif (titre, chapô, corps, conclusion) et d'un tract (slogan, impératifs, appel à l'action).
\item[$\square$] Je sais repérer et traduire les verbes à particule séparable (\all{vorbeugen, stattfinden} : \all{Die Aktion findet statt.} « L'action a lieu. »).
\item[$\square$] Je me méfie des faux amis : \all{die Krankheit} = la maladie (et non « la crainte ») ; \all{bekommen} = recevoir (et non « devenir » = \all{werden}) ; \all{aktiv} = actif, mais « actuel » se dit \all{aktuell}.
\item[$\square$] Je relis ma production : orthographe, cas, ordre des mots, ponctuation.
\item[$\square$] Je gère mon temps : lecture du sujet, plan rapide, rédaction, relecture.
\end{itemize}
\end{aretenir}

\findecours

\end{document}
